\chapter{文本/图像到3D生成}

\paragraph*{本章阅读锚点}
3DGS 进入生成任务后，核心矛盾从“如何重建已知相机的场景”变成“如何从文本、单图、视频或噪声中产生一个三维一致、可渲染、可编辑的高斯资产”。早期锚点是 DreamGaussian 和 GaussianDreamer：用 SDS 或 3D 扩散初始化，把 2D 扩散的语义先验蒸馏到高斯里。后续方法逐渐减少逐对象优化，转向直接生成高斯：把高斯属性编码成图像、UV、函数、token、体素 latent、splatter image 或 diffusion denoiser 的输出。阅读本章时重点看每个方法改变了哪一环：是初始化几何、定义高斯潜空间、把扩散模型从 2D 改到 3D、用视频扩散保证多视角一致，还是把静态生成扩展为 4D 动画。

\section{3D 生成模型与 GAN}

这一节的主线是从“用扩散优化高斯”走向“让生成模型直接输出高斯”。方法之间最大的差异在表示：DiffSplat 把高斯属性看成可被图像扩散模型处理的 latent；DiffusionGS 把高斯预测烘进 denoiser；DiffGS 把离散高斯改写成连续函数；UVGS、Can3Tok、ARGS 和 SplatFlow 则分别用 UV 图、场景 token、层次树或多视图 flow 来管理无结构高斯。



% METHOD: 2406.02968
\paragraph{GSGAN（层次多尺度高斯 GAN）（2024.06）}
GSGAN~\cite{2406_02968} 把 3D Gaussians 用作 3D GAN 表示。朴素高斯 GAN 训练不稳定，因为没有 densification 和尺度管理；GSGAN 因此采用粗到细层次结构，精细高斯的位置受粗级高斯约束，尺度随层级递减。生成器从随机 latent 生成多尺度高斯，再由 2D discriminator 对渲染图做对抗训练。相对扩散方法，GSGAN 展示了另一条路线：用层次正则让高斯可端到端对抗生成。

具体实现中，生成器不直接吐出一堆彼此独立的高斯，而是先生成粗层级锚点，再在每个粗高斯的局部坐标系中生成细层级子高斯。子高斯的位置由父高斯的位置、旋转和尺度变换得到，因此不会在训练早期到处漂移；子高斯尺度还被约束为随层级变细单调减小，相当于把 3DGS 中依赖启发式 densification 的“粗到细增长”写进生成器结构。渲染时仍使用高斯 rasterization，把给定相机下的图像送给 2D 判别器；判别器学习真实图像分布和相机一致性，生成器则学习产生能跨视角渲染一致的高斯层次。这个设计的关键不在于换了一个 GAN 损失，而在于让 GAN 输出的高斯天生带有可训练的尺度组织。
\textbf{实现细节：} GSGAN（2024.06） 直接在对抗训练框架中生成3D高斯，无需NeRF中间表示。但方法发现，朴素生成器（直接合成一组无约束的高斯）在对抗训练中面临两个问题：训练早期所有高斯从渲染图像中消失（训练不稳定），以及收敛后出现长条形高斯导致的视觉伪影。根本原因在于3DGS的非可微场景依赖特性（如SfM初始化和启发式致密化）无法直接迁移到GAN框架中。为解决这些问题，GSGAN引入分层多尺度高斯表示。层级从粗到细编号为 $l \in \{0, \ldots, L-1\}$，不同层级的精细高斯由其对应的粗层级高斯参数化——精细高斯的位置被约束在粗高斯的局部坐标系内，且尺度随层级加深单调递减： 其中 $\boldsymbol{\mu}^{(l)}$、$s^{(l)}$ 和 $R^{(l)}, S^{(l)}$ 分别为第 $l$ 层高斯的位置、尺度和旋转/缩放矩阵，$\Delta\boldsymbol{\mu}^{(l)}$ 是由 tanh 激活约束的局部位置偏移。尺度差使用 softplus 取负值确保单调递减。直觉上，粗层级建模场景的大尺度结构，精细层级逐步补充细节。层级间通过注意力层（底层用自注意力，高层用KNN（K-Nearest Neighbors，K近邻）局部注意力，$k=16$）建模高斯交互，KNN正则化损失鼓励邻居高斯具有一致的属性，进一步稳定训练。
\textbf{方法信息：} 代码：\url{https://github.com/hse1032/gsgan}。



% METHOD: 2406.06367
\paragraph{MVGamba（状态空间多视图高斯重建器）（2024.06）}
MVGamba~\cite{2406_06367} 把 3D 内容生成统一为 state space sequence modeling。它沿用“先生成多视图，再重建高斯”的路线：MVDream 或 ImageDream 负责从文本/单图产生多视图图像，后端再把这些图像转成 3DGS 或 2DGS。问题在于，多视图图像 token 和高斯 token 数量都很长，普通 Transformer 跨视图建模代价高，容易把计算花在二次注意力上。

MVGamba 的改变是用 SSM 层替换重建器里的全局注意力。实现时，多视图图像特征先按视角和空间位置展开为序列，SSM 以线性复杂度沿序列传播上下文，让一个视角的可见边界、遮挡关系和纹理信息能影响其他视角。随后轻量 Gaussian decoder 预测每个高斯的位置、尺度、旋转、不透明度和颜色；使用 2DGS/surfel 分支时，高斯盘面还能更直接贴合表面。训练时，预测出的高斯会被渲染回输入视角和目标视角，通过图像重建损失、深度或几何一致性约束来更新解码器，而不是直接监督某个固定顺序的高斯列表。相对 DSplats 依赖扩散模型逐步去噪高斯，MVGamba 更像一个快速前馈重建器：多视图扩散负责想象，SSM 负责把跨视图证据高效汇合成显式高斯。
\textbf{实现细节：} MVGamba（2024.06） 是一种基于状态空间模型（State Space Model, SSM）的轻量级高斯重建模型。已有的基于 Transformer 的高斯重建模型往往面临多视角不一致和纹理模糊的问题——MVGamba 将此归因于 Transformer 架构在追求强大表达力时牺牲了多视角信息传播的效率。MVGamba 的核心架构采用了类 RNN 的状态空间模型替代 Transformer，其高斯重建器在生成高斯序列的同时，通过因果上下文传播多视角信息： \textit{白话解释}：公式 (相应公式)描述了 SSM 的核心递推关系。隐状态$\mathbf{h}_t$在时间步$t$由前一步隐状态$\mathbf{h}_{t-1}$经矩阵$\bar{\mathbf{A}}$变换后加上当前输入$\mathbf{x}_t$经$\bar{\mathbf{B}}$变换的结果。输出$\mathbf{y}_t$（即高斯属性）由隐状态经$\bar{\mathbf{C}}$投影得到。$\bar{\mathbf{A}}$的因果性保证了多视角信息的有序传播，实现跨视图自精化（cross-view self-refinement）。MVGamba 统一了单图、稀疏图和文本驱动的3D生成任务。
\textbf{方法信息：} 代码：\url{https://github.com/skyworkai/mvgamba}。
\textbf{设计动机与边界：} 相对原始 3DGS 以独立高斯和图像重建误差为主的设定，MVGamba 把问题落在以下可实现的改动上：MVGamba采用两阶段管线：(1) 利用现成多视图扩散模型（MVDream、ImageDream）从文本或单图生成多视图图像。(2) SSM-based多视图高斯重建器将多视图图像转换为3D高斯。核心架构创新：将多视图图像tokens展开为序列，通过SSM层以线性复杂度传播因果上下文，实现跨视图信息自精炼。相比Transformer的O(n\^{}2)复杂度，SSM的O(n)复杂度允许处理更长的高斯序列。包含轻量高斯解码器预测每个高斯的属性。使用2DGS作为基础表示以获得更准确的几何。训练目标包括重建损失和多视图一致性损失。其设计目的不是单独增加一个模块，而是用上述表示、监督或渲染约束消除原流程在目标任务中的具体失配。提出MVGamba，基于RNN风格状态空间模型(SSM)的多视图高斯重建器。解决现有方法中多视图信息传播与计算成本之间的妥协问题。SSM以线性复杂度传播因果上下文实现跨视图自精炼，同时生成长序列高斯进行细节建模。模型大小仅SOTA的0.1倍，在所有3D生成场景中超越SOTA。适用时还应保留论文给出的边界：依赖预训练多视图扩散模型的质量；SSM的序列建模可能限制全局信息交互。
% METHOD: 2410.19657
\paragraph{DiffGS（函数化高斯扩散）（2024.10）}
DiffGS~\cite{2410_19657} 把离散、无序的 3DGS 改写成三个连续函数：高斯概率函数预测某个位置是否为中心，颜色函数预测颜色，变换函数预测旋转和缩放。高斯 VAE 将这些函数编码成 latent，再在函数 latent 上训练扩散模型。生成后，八叉树引导采样从高概率区域提取任意数量高斯。相对 UVGS 和 DiffSplat 的“把高斯属性排列成可扩散数据”，DiffGS 的锚点是函数空间：先生成连续场，再采样离散高斯。
\textbf{实现细节：} DiffGS（2024.10） 的核心洞察是绕过离散障碍：不直接对离散高斯点集做扩散，而是将3DGS分解为三个连续函数，在函数的连续表示上训练潜在扩散模型。这一思路类似于用 SDF（符号距离函数）表示离散网格。连续函数不依赖点的数量和排列顺序，天然兼容扩散模型的生成范式。具体而言，方法定义三个解耦的高斯泼溅函数：  \par  概率函数 $\mathrm{GauPF}(\mathbf{x})$：建模空间中任意位置 $\mathbf{x}$ 存在高斯基元的概率密度。直觉上，离所有高斯中心越远的点，被高斯占据的概率越低。真值由查询点到最近高斯中心的最小距离定义，经截断函数过滤极端值后，再用连续映射将距离转换为 $[0,1]$ 区间的概率值 \par  颜色函数 $\mathrm{GauCF}(\mathbf{x})$：给定已知的高斯基元位置 $\mathbf{x}$，预测该位置的颜色属性 $c$ \par  变换函数 $\mathrm{GauTF}(\mathbf{x})$：给定已知的高斯基元位置 $\mathbf{x}$，预测旋转 $r$、缩放 $s$ 和不透明度 $o$ 等几何变换属性  三函数的解耦设计反映了不同的建模逻辑：GauPF 从3D空间中探索几何结构（"这个位置是否存在高斯基元？"），而 GauCF 和 GauTF 在已知几何的基础上预测外观和变换属性（"这个位置的高斯长什么样？"）。这种分离使得几何生成和属性预测可以独立优化。高斯VAE与潜在扩散。三个连续函数虽然化解了离散性障碍，但它们仍然是高维的神经场表示，无法直接作为扩散模型的输入。DiffGS 训练一个高斯变分自编码器（VAE）来完成压缩。GS编码器将整个3DGS形状编码为一个潜向量 $z$，三平面解码器将 $z$ 解码为特征三平面 $t \in \mathbb{R}^{H \times W \times C \times 3}$（三个正交特征平面 $t_{XY}$、$t_{XZ}$、$t_{YZ}$）。对于3D空间中任意查询点 $\mathbf{x}$，将其投影到三个平面并三线性插值得到特征 $f$，再分别送入三个神经预测器 $\phi_{\mathrm{pf}}$、$\phi_{\mathrm{cf}}$、$\phi_{\mathrm{tf}}$ 预测概率、颜色和变换。VAE 的训练目标包含属性预测的 $L_1$ 损失和潜空间的 KL 散度正则化。
\textbf{方法信息：} 引用很多，Semantic Scholar 约 58 次。
\textbf{设计动机与边界：} 相对原始 3DGS 以独立高斯和图像重建误差为主的设定，DiffGS 把问题落在以下可实现的改动上：方法核心是将离散3DGS转化为连续函数表示：1)高斯概率函数(GauPF)：预测3D位置是高斯中心的概率，编码几何信息。2)高斯颜色函数(GauCF)：给定3D位置预测高斯颜色属性。3)高斯变换函数(GauTF)：给定3D位置预测高斯的旋转和缩放。训练高斯VAE将函数编码为压缩潜向量，在潜空间训练潜扩散模型(LDM)生成新的3DGS。生成后通过八叉树引导采样算法：先在粗八叉树层级采样高概率位置作为候选，再用优化细化位置，可提取任意数量的3D高斯。支持无条件生成和多种条件生成（文本/图像/部分3DGS/点云）。其设计目的不是单独增加一个模块，而是用上述表示、监督或渲染约束消除原流程在目标任务中的具体失配。提出DiffGS，首个基于函数解耦的3DGS扩散生成模型。将离散无结构的3DGS表示为三个连续函数（概率/颜色/变换），在函数空间训练潜扩散模型，通过八叉树引导采样从生成的函数中提取任意数量的3D高斯。支持无条件生成、文本/图像/部分3DGS条件生成和点云到高斯转换等多种任务。适用时还应保留论文给出的边界：生成质量受限于LDM训练数据规模；大规模场景生成困难；生成的高斯可能存在局部几何伪影；推理速度受采样步骤影响。
% METHOD: 2411.14384
\paragraph{DiffusionGS（把高斯烘进扩散去噪器）（2024.11）}
DiffusionGS~\cite{2411_14384} 解决的是另一类断裂：许多单图到 3D 方法先用 2D 多视图扩散生成若干视角，再把这些图像交给 3D 重建网络。这个流程中，扩散模型本身没有 3D 表示，视角稍微变化就可能生成不一致的纹理、结构或遮挡。DiffusionGS 的做法是把 3DGS 烘进扩散去噪器，让 denoiser 在每个去噪时间步直接预测一组 pixel-aligned Gaussians，而不是先预测干净图像再事后重建。

训练时，模型输入一个干净条件视图、若干带噪目标视图、对应相机条件和扩散时间步。Transformer 主干输出多张高斯属性图，每个像素对应一个候选高斯，因此总数固定为 $(N+1)HW$。每个高斯的中心不直接自由回归，而是沿该像素射线用 near/far 深度插值得到；尺度、旋转、颜色和不透明度由网络输出。由于没有逐高斯真值，监督不在属性空间做，而是把预测高斯从训练视角和新视角光栅化，再用多视图渲染图像损失训练 denoiser。这样，扩散过程中的中间变量本身就是可渲染的 3DGS。

为了同时覆盖物体和场景，DiffusionGS 还设计了 scene-object mixed training。物体数据来自 Objaverse、MVImgNet，场景数据来自 RealEstate10K、DL3DV；直接混合会因为尺度、相机轨迹和深度范围差异过大而不收敛，所以方法控制视图采样、相机条件、成像深度和高斯点云范围，并提出 RPPC（Reference-Point Plucker Coordinates）作为相机条件，让网络更明确地感知射线相对世界原点的深度关系。推理时，给定单张输入图像，模型沿目标相机轨迹生成高斯表示，论文报告约 6 秒完成物体生成或场景重建。

相对 DiffSplat，DiffusionGS 不需要先把高斯属性压进 VAE latent 再扩散，而是直接把“去噪一步要输出什么”改成高斯点云；相对 LGM/GRM 等重建式方法，它把多视图一致性约束提前到扩散模型内部。这个区别使它更适合作为“3DGS-native denoiser”的锚点。
\textbf{锚点细化：} DiffusionGS：针对``先生成多视角图像再重建 3D''两阶段范式视角不一致时易崩溃的缺陷，用\textbf{单阶段 3D 扩散在每个去噪时间步直接输出像素对齐的 3D 高斯点云}替代中间 2D 图像，核心思想是\textbf{高斯中心位置由像素射线参数化保证空间定位，扩散过程天然强制视角一致性}；配合场景-物体混合训练策略（Objaverse+MVImgNet 物体先验、RealEstate10K+DL3DV 场景先验）和相对位姿条件化（RPPC）统一参数化方式，6 秒推理、PSNR 提升超 2 dB。
\textbf{实现细节：} DiffusionGS（2024.11） 指出先前方法的两阶段范式（先生成多视角图像，再重建3D）存在两个根本缺陷：多视角扩散模型缺乏3D模型，视角不一致时容易崩溃；且提示视角方向受限。方法提出单阶段3D扩散模型，在每个去噪时间步直接输出像素对齐的3D高斯点云 $G_\theta = \{g_t^{(k)}(\mu_t^{(k)}, \Sigma_t^{(k)}, o_t^{(k)}, c_t^{(k)})\}$，而非中间的2D图像，天然保证了视角一致性。高斯中心位置通过 $\mu_t^{(k)} = o^{(k)} + u_t^{(k)} d^{(k)}$ 参数化，其中距离 $u_t^{(k)} = w_t^{(k)} u_{near} + (1-w_t^{(k)}) u_{far}$ 由网络预测权重控制。为同时学习物体和场景级别的3D先验，方法开发场景-物体混合训练策略，并设计参考点Pl\"ucker坐标（RPPC）编码每个射线上距世界坐标系原点最近的点，帮助模型感知深度和3D几何。
\textbf{设计动机与边界：} 相对原始 3DGS 以独立高斯和图像重建误差为主的设定，Baking Gaussian Splatting into Diffusion Denoiser for Fast and Scalable Single-stage Image 把问题落在以下可实现的改动上：DiffusionGS是一个单阶段3D扩散模型。核心设计包括：(1) 3D扩散去噪器：不同于传统扩散模型预测噪声或2D多视图扩散模型，DiffusionGS在每个去噪时间步直接预测像素对齐的3D Gaussian基元，通过高并行光栅化和可扩展的成像范围，在~6秒内完成推理；(2) 像素对齐Gaussian：每个Gaussian的位置由像素射线的原点和方向参数化，数量固定为(N+1)HW，确保3D一致性；(3) 场景-物体混合训练策略：解决直接混合场景和物体数据导致的不收敛问题，通过控制选定的视图分布、相机条件、Gaussian点云和成像深度来适应两类数据，学习通用的几何和纹理先验；(4) 多视图生成：不同于仅从给定视图预测3D表示的方法，DiffusionGS沿相机轨迹生成其他视图，更好地感知几何来重建场景，无需深度估计器。其设计目的不是单独增加一个模块，而是用上述表示、监督或渲染约束消除原流程在目标任务中的具体失配。提出首个将3DGS嵌入扩散去噪器的单阶段3D扩散模型，通过像素对齐Gaussian实现3D一致性，场景-物体混合训练实现泛化，速度提升5x以上适用时还应保留论文给出的边界：论文未明确列出局限，但像素对齐Gaussian的数量与图像分辨率相关，可能限制高分辨率输出。
\textbf{方法信息：} Adobe 参与了这篇论文；代码：\url{https://github.com/caiyuanhao1998/open-diffusiongs}；GitHub 关注度高，约 853 stars。该仓库说明其开源实现是重新实现和重新训练版本，与 Adobe 原始内部版本存在差异。



% METHOD: 2411.16443
\paragraph{SplatFlow（多视图 Rectified Flow 到高斯）（2024.11）}
SplatFlow~\cite{2411_16443} 统一文本驱动的 3DGS 生成和编辑。DiffusionGS 让扩散 denoiser 直接输出高斯属性，但主要围绕单图输入和目标视角重建；SplatFlow 的改变是把图像、深度、相机位姿和高斯解码放进同一个多视图 rectified flow 流程中。模型在 latent space 里生成一组多视图表示，再由 Gaussian Splatting Decoder 转成 pixel-aligned 3DGS。

实现上，第一部分是多视图 Rectified Flow 模型。它借用 Stable Diffusion 3 的编码器并冻结共享，在潜空间同时预测多视图图像 latent、深度和相机条件，使不同视角不是各自独立生成，而是在同一个流模型里对齐。第二部分是 GSDecoder：它基于 SD3 decoder 改造，加入跨视图注意力，让一个视角的纹理和几何线索能影响其他视角；可学习深度特征和 ray embedding 则把每个像素的射线位置送进解码器，弥补冻结编码器损失的细节。推理时，文本或图像条件先得到多视图 latent，再一次性解码出像素对齐高斯；做编辑时，方法利用 inversion 把已有 3DGS 投回生成空间，再通过 inpainting 修改局部区域。相对 DiffusionGS，SplatFlow 更强调“生成、相机、深度和编辑”在同一多视图流框架中闭合。
\textbf{设计动机与边界：} 相对原始 3DGS 以独立高斯和图像重建误差为主的设定，SplatFlow 把问题落在以下可实现的改动上：SplatFlow包含两个核心组件：(1) 多视图Rectified Flow（RF）模型：基于Latent Diffusion Model架构，冻结并共享Stable Diffusion 3的编码器，在潜在空间联合生成多视图相机位姿、深度图和图像，解决真实场景中多样化场景尺度和复杂相机轨迹的挑战。(2) Gaussian Splatting Decoder（GSDecoder）：基于SD3解码器架构改进，增加跨视图注意力机制，将多视图潜在表示解码为像素对齐的3DGS。为补偿冻结编码器损失的细节信息，引入可学习的深度特征和射线嵌入作为额外输入。损失函数结合MSE、LPIPS和vision-aided GAN loss。推理时利用无需训练的inversion（基于SD3的inversion）和inpainting（Repaint技术）实现3DGS编辑。其设计目的不是单独增加一个模块，而是用上述表示、监督或渲染约束消除原流程在目标任务中的具体失配。提出SplatFlow统一框架，首次实现文本驱动的直接3DGS生成与编辑。多视图RF模型同时生成图像、深度和位姿，GSDecoder高效转换为3DGS。利用无需训练的inversion和inpainting技术支持3DGS编辑，在统一框架内完成物体编辑、新视角合成、相机位姿估计等多种3D任务。适用时还应保留论文给出的边界：生成质量受限于2D扩散模型的能力；长序列多视图一致性仍有挑战；GSDecoder依赖冻结编码器可能损失精细空间细节
\textbf{方法信息：} 代码：\url{https://github.com/gohyojun15/splatflow}；GitHub 关注度较高，约 142 stars。



% METHOD: 2412.09648
\paragraph{DSplats（去噪过程中的高斯重建器）（2024.12）}
DSplats~\cite{2412_09648} 把 LGM 或 Splatter 类高斯重建器与 latent diffusion denoiser 结合。噪声图像先映射到 latent space，3D-aware denoising network 在相机条件下进行多视图去噪，并从去噪 latent 中预测 3D Gaussian primitives。训练同时包含扩散去噪损失和渲染重建损失。相对普通 U-Net 多视图扩散，显式高斯表示给 denoiser 强三维归纳偏置；相对 DiffusionGS，它更像在现有多视图重建器骨干中替换为 diffusion prior。
\textbf{实现细节：} DSplats（2024.12） 重新审视了高斯重建模型的骨干架构：传统方法使用像素级 U-Net，与预训练的 Latent Diffusion Model（如 Stable Diffusion）不兼容。DSplats 将骨干替换为预训练的 LDM，并引入 VAE 编码器/解码器在像素空间与隐空间之间映射。然而，标准 LDM 仅针对图像生成优化，因此 DSplats 对其进行多项适应性改造：添加噪声编码器处理多视角输入、在去噪网络中注入3D感知的相机位姿条件，以及通过可微高斯光栅化将隐空间的中间特征解码为显式3D高斯。整个框架以单次前向推理的方式联合完成多视角去噪和3D重建。
\textbf{方法信息：} Apple 参与了这篇论文。
\textbf{设计动机与边界：} 相对原始 3DGS 以独立高斯和图像重建误差为主的设定，DSplats 把问题落在以下可实现的改动上：DSplats将高斯重建器（如LGM、Splatter）与传统U-Net骨干替换为Latent Diffusion Model（如Stable Diffusion）。整体框架：(1) 噪声编码器：将输入图像映射到潜在空间；(2) 3D感知去噪网络：在潜在空间执行多视图扩散去噪，同时结合相机位姿条件注入3D感知；(3) 高斯Splatting：从去噪后的潜在表示预测3D高斯基元集合。训练时联合优化扩散损失（去噪潜在）和渲染损失（视图重建）。使用VAE在图像和潜在空间之间映射。对标准潜在扩散模型进行了多项修改以适应重建任务。去噪网络中的显式3D表示提供了强归纳偏置，确保多视图几何一致性。其设计目的不是单独增加一个模块，而是用上述表示、监督或渲染约束消除原流程在目标任务中的具体失配。提出DSplats，首个在多视图扩散去噪过程中直接使用高斯Splat重建器的方法。将预训练Latent Diffusion Model整合到重建器骨干中预测3D高斯，显式3D表示为去噪网络提供强归纳偏置，确保几何一致的新视角生成。在GSO数据集上PSNR 20.38、SSIM 0.842、LPIPS 0.109。适用时还应保留论文给出的边界：计算资源需求较高（扩散模型+3DGS）；生成质量受训练数据分布限制；对复杂场景的生成能力有限。
% METHOD: 2412.16253
\paragraph{ExCellGen（单样本交互式 3D 场景生成）（2024.12）}
ExCellGen~\cite{2412_16253} 处理的是“只有一个真实场景样本时如何生成同风格新布局”。LucidDreamer 和 BloomScene 主要依赖扩散模型从文本或图像外推场景，ExCellGen 则从已有视频或无人机素材先重建一个 3DGS 场景，再把这个场景当作可复用的素材库。第一步是语义增强重建：标准 3DGS 表示外观和几何，DINO 特征附加到高斯上，用于交互式区域选择和语义引导；这些高维特征再压到低维，避免生成器直接处理过重的语义向量。

第二步把不规则高斯转成稀疏体素抽象。每个体素存储外观统计和语义特征，轻量生成元胞自动机把当前体素及其邻域特征作为输入，输出下一轮体素状态和语义属性，因此它学习的是“局部块如何生长和组合”，不是直接回归每个高斯的坐标。生成结果还不是最终高斯，而是一张“应该放什么类型块”的稀疏地图；块重映射模块再从原始 3DGS 中寻找相似局部高斯块，把真实场景的纹理、植被、建筑边缘等外观统计拷贝回来。最后渲染时仍使用原始高斯块的细节，只改变它们在新布局中的排列和组合。这个设计的增量在于，它不用训练一个大规模场景扩散模型，而是把单个场景拆成可语义检索、可局部重排的高斯块，适合交互式编辑和少样本扩展。
\textbf{实现细节：} ExCellGen（2024.12） 从单个真实世界样本（手持视频或无人机拍摄）出发生成3D场景，解决"从单例学习生成"的独特问题。与需要大规模3D数据集训练的生成方法不同，ExCellGen将3D场景生成建模为基于样例的合成（exemplar-based synthesis）任务——灵感源自纹理合成的经典范式——仅需单个真实场景作为训练数据。方法采用两阶段管线：第一阶段将不规则高斯抽象为稀疏特征化体素网格，用轻量级生成模型在亚秒级合成完整体素；第二阶段通过基于块的重新映射从原始3D高斯中恢复高质量外观。整个管线实现了从单个真实样本到可泛化场景生成的飞跃。语义增强的3D高斯重建。输入阶段首先用3DGS重建输入场景。具体实现采用DINO特征：从多视角图像中提取768维2D特征图，通过PCA降维至8维，再通过将可微渲染方程中的颜色 $c^{(i)}$ 替换为特征 $f^{(i)}$ 将其提升至3D空间，经优化后每个高斯获得一致的语义特征。这些语义特征同时服务于两个目的：一是为体素化和生成提供语义引导，二是通过k-means聚类实现交互式场景区域选择——用户可在GUI中调整聚类数量并直接操纵3D高斯来选取感兴趣的区域，解决了真实场景中背景杂乱、物体边界模糊的分割难题。体素化与特征编码。为将不规则的3D高斯集合转化为便于生成的规整表示，方法将其抽象为稀疏特征化体素网格 $\mathcal{V} = \{v_i\}_{i=1}^N$，其中 $N$ 为占据体素数，每个体素由位置索引 $p_i \in \mathbb{R}^{}3$ 和特征向量 $z_i$ 组成。体素化过程的关键参数如下：仅在不透明度超过阈值 $\alpha=0.1$ 的位置实例化体素；精细分辨率 $r_f = 64^{}3$ 用于生成表达力丰富的场景细节，粗糙分辨率 $r_c = 16^{}3$ 通过下采样获得，作为条件信号。为维持体素间的多样性，每个体素选取其所包含的不透明度最高的高斯的特征（而非简单平均），提取其球谐系数（外观，4维）和DINO语义特征（4维），各自独立PCA降维后重归一化，拼接为8维联合特征向量。这一稀疏体素表示在保留场景外观和语义信息的同时，将高斯的数量级压缩至数千个占据体素，大幅降低了生成任务的计算负担。
\textbf{方法信息：} 代码：\url{https://github.com/clementjambon/excellgen}。
\textbf{设计动机与边界：} 相对原始 3DGS 以独立高斯和图像重建误差为主的设定，ExCellGen 把问题落在以下可实现的改动上：ExCellGen的两阶段流程：(1)语义增强3DGS重建：用3DGS重建输入场景，附加DINO语义特征实现区域选择和生成引导，通过PCA降维到8维；(2)稀疏体素抽象：将不规则高斯抽象为稀疏体素网格，每个体素编码外观和语义特征向量；(3)生成元胞自动机（GCA）：轻量级生成模型，以粗体素输入为条件，每个场景训练不到10分钟，生成新体素布局仅需0.5-2秒；(4)块重映射：将GCA生成的体素特征通过块匹配从原始3DGS场景中找到最相似的高斯块进行合成，恢复输入场景的外观统计。支持交互式GUI中的区域选择和生成控制。其设计目的不是单独增加一个模块，而是用上述表示、监督或渲染约束消除原流程在目标任务中的具体失配。提出ExCellGen框架，从单个手持视频或无人机拍摄素材，在10分钟内训练完成，0.5-2秒生成3D场景。利用3DGS高质量重建，GCA生成稀疏体素特征，块重映射从样本高斯中恢复外观统计，实现完全可控的交互式3D生成。适用时还应保留论文给出的边界：生成质量受限于样本场景的复杂度；GCA为每个场景单独训练；块重映射可能产生不自然的拼接边界。
% METHOD: 2501.05427
\paragraph{Zero-1-to-G（多属性图像分解的直接 3D 生成）（2025.01）}
Zero-1-to-G~\cite{2501_05427} 把 splatter image 的高维属性拆成多个三通道属性图像，例如颜色、尺度、旋转、不透明度和位置，使它们更接近预训练 2D 扩散模型的输入分布。模型在 Stable Diffusion 中加入跨视角注意力交换不同视角信息，并加入跨属性注意力保持属性间相关。相对 DiffSplat 的统一 latent，Zero-1-to-G 更强调属性分解：让每类高斯属性以 2D 扩散更擅长的形式生成。
实现上，它先用 LGM 一类前馈网络产生较平滑的 Splatter Image，而不是逐对象优化出带大量高频伪影的属性图；然后把约 14 通道的高斯属性拆成 5 组更接近 RGB 图像分布的属性图，并微调 VAE 解码器以减少高斯属性在编码/解码中的损失。扩散主干仍继承 Stable Diffusion 的能力，但在注意力层里同时做两种交互：跨视角注意力让不同视角共享几何信息，跨属性注意力让位置、尺度、旋转、不透明度和颜色不要各自漂移。这样，模型生成的不是普通多视图图像，而是一组可以直接组合成 3DGS 的属性图。
\textbf{方法信息：} Apple 参与了这篇论文。
\textbf{设计动机与边界：} 相对原始 3DGS 以独立高斯和图像重建误差为主的设定，Zero-1-to-G 把问题落在以下可实现的改动上：Zero-1-to-G的核心洞察是将14通道的Splatter Image分解为5个3通道属性图像（RGB颜色、尺度、旋转、不透明度、位置），每个属性图像在预训练2D扩散模型的分布范围内。训练流程：（1）使用LGM神经网络（而非逐对象优化拟合）生成平滑的Splatter Image，避免高频伪影导致VAE编解码质量下降。（2）微调预训练VAE解码器增强解码质量。（3）在Stable Diffusion框架中引入跨视角注意力层实现不同视角间信息交换，跨属性注意力机制维持属性间相关性。支持在隐空间中进行扩散训练，高效利用预训练2D扩散先验。其设计目的不是单独增加一个模块，而是用上述表示、监督或渲染约束消除原流程在目标任务中的具体失配。提出将高斯溅射分解为多视角属性图像的创新方法，将3D生成问题重新定义为2D扩散框架。引入跨视角和跨属性注意力层捕捉复杂相关性并强制3D一致性。首个有效利用预训练2D扩散先验的直接图像到3D生成模型。适用时还应保留论文给出的边界：依赖LGM生成训练数据质量；分解过程可能丢失高斯间细微关联；对复杂拓扑结构的生成能力有限。
% METHOD: 2501.06714
\paragraph{F3D-Gaus（循环聚合的前馈 3D 感知生成）（2025.01）}
F3D-Gaus~\cite{2501_06714} 的目标不是为每个输入单独做 SDS 优化，而是从 ImageNet 等单目图像集合中学一个前馈 3D-aware 生成器。它继承像素对齐高斯的思路：U-Net 读取 RGB-D 输入，在每个像素附近预测一组 Gaussian 属性，包括 3D 坐标、不透明度、颜色、尺度和旋转；其中 3D 坐标由输入深度加预测偏移得到，因此高斯一开始就贴着输入视图的几何展开。

困难在于，单目数据没有真实新视角监督，直接训练会让模型只会复现输入视角。F3D-Gaus 的改变是加入 cycle-aggregative self-supervision：模型先从一个视角预测高斯并渲染到另一个视角，再把这个渲染结果作为新的输入继续预测，要求循环回来的视图仍能对齐原图。多个循环视角产生的像素对齐高斯会被聚合起来，使单视角看不到的区域能由其他推断视角补充。相对 Zero-1-to-G 的扩散直接生成，F3D-Gaus 更像从单目图像集合中自监督训练一个可一次前向输出高斯的重建器。
\textbf{实现细节：} F3D-Gaus（2025.01） 面向一个更困难的问题：仅从单目数据集（如ImageNet，含420K训练图像）学习3D感知生成，无需多视角或动态视频数据。多视角数据集获取成本高、需要专用设备，而单目数据集可从网络大量获取，具有更好的可扩展性。方法提出基于像素对齐高斯（pixel-aligned Gaussian Splatting）的前馈管线：U-Net backbone 以 RGB-D 输入，预测像素对齐的高斯属性图（位置偏移、不透明度、颜色、缩放、旋转），深度图由现成的单目深度估计模型提供。核心挑战在于单目数据集仅提供正面监督，侧面视角的渲染质量无法直接验证。为此，方法引入循环聚合（cycle-aggregative）自监督策略：先从正面输入预测 $GS_0$，渲染侧面视角得到 $(I_1, D_1)$，再用该侧面输入预测 $GS_1$，将两个GS在高斯级别做互补聚合——在重叠区域取属性平均，在非重叠区域保留各自的高斯基元。重建损失 $\mathcal{L}_{recon} = \|I_0 - I_0^{GT}\|_1 + \|D_0 - D_0^{GT}\|_1$ 保证正面视角质量，聚合损失确保跨视角一致性。此外，方法利用视频扩散模型先验（CogVideoX）进行几何感知精炼，在宽视角场景中提升3D纹理质量。在ImageNet全量评估上，F3D-Gaus在FID和非平坦度分数（NFS）上均优于G3DR基线。
\textbf{方法信息：} 代码：\url{https://github.com/w-ted/f3d-gaus}。



% METHOD: 2501.07574
\paragraph{uCO3D（大规模物体视频与 3DGS 数据资产）（2025.01）}
UnCommon Objects in 3D~\cite{2501_07574} 是数据集工作。它提供大规模高分辨率物体视频、相机位姿、深度、稀疏点云、标题和 3DGS 重建，覆盖大量物体类别和完整 360 度视角。其技术意义在于，生成式 3DGS 不只需要模型，也需要高质量多视图和高斯资产作为训练分布。相对 DiffSplat 的方法本身，uCO3D 提供可训练 LRM、MVPlanes 或高斯生成模型的数据基础。

它的增量不在某个网络模块，而在数据形态。每个物体序列不仅给 RGB 视频，还提供相机位姿、深度、稀疏点云、人工核验的文本标题以及 3DGS 重建结果；完整环绕覆盖减少了只看正面或少数视角带来的偏差。对 3DGS 生成模型来说，这意味着训练样本可以同时服务于多视图监督、前馈重建、跨模态检索和高斯资产学习。论文还把该数据用于 LRM、MVPlanes 等模型训练，说明数据多样性和标注质量会直接影响后续高斯生成器的泛化能力。
\textbf{实现细节：} UnCommon Objects in 3D（2025.01） 专注于罕见物体类别的 3D 生成，通过分布外泛化扩展了重建模型的适用范围。
\textbf{实现补充：} 为避免只把它读成名称或数据点，还应把实现接口明确拆开：输入是什么，输出是高斯参数、图像、轨迹、标注还是评测指标，哪些参数、网络模块、训练损失或渲染接口真正接到原始 3DGS / Scaffold-GS / 4D-GS 管线上，以及剪枝、分裂、采样、解码、优化或评测协议在哪一步发生。若该条目是survey、benchmark 或 dataset，这里的“实现”对应论文集合、任务划分、相机/轨迹/标注组织、评价指标和基线配置，而不是虚构新的 rasterizer；若是算法条目，则对应高斯位置、尺度、旋转、不透明度、特征或 MLP 解码如何进入训练和渲染流程。
\textbf{设计动机与边界：} 相对原始 3DGS 以独立高斯和图像重建误差为主的设定，UnCommon Objects in 3D 把问题落在以下可实现的改动上：uCO3D数据集的特点：(1)规模和多样性：超过1000个物体类别，是MVImgNet和CO3Dv2的数倍，覆盖日常物品、交通工具、家具等。(2)高质量标注：每个物体配有3D相机位姿、深度图、稀疏点云，以及人工验证的标题描述和3DGS重建。(3)严格质量控制：对采集视频和3D标注进行多重质量检查，确保数据可靠性。(4)完整360度覆盖：每个物体确保全方位视角覆盖，避免视角偏差。数据集在训练LRM、MVPlanes等3D大模型时，相比MVImgNet和CO3Dv2均取得更优结果，验证了数据质量和多样性的重要性。其设计目的不是单独增加一个模块，而是用上述表示、监督或渲染约束消除原流程在目标任务中的具体失配。构建了uCO3D数据集，这是迄今最大规模的公开高分辨率物体视频集合，覆盖1000+物体类别，提供完整360度覆盖和严格质量控制。每个物体配备3DGS重建、标注、标题，训练3D大模型在uCO3D上取得优于MVImgNet和CO3Dv2的效果。适用时还应保留论文给出的边界：尽管规模最大，但某些细粒度类别样本仍然有限；3D标注依赖SfM/COLMAP，可能存在累积误差。
\textbf{方法信息：} Meta AI 参与了这篇论文；官方代码和数据工具：\url{https://github.com/facebookresearch/uco3d}；GitHub 关注度很高，约 1.3k stars。
% METHOD: 2501.16764
\paragraph{DiffSplat（图像扩散模型原生生成高斯）（2025.01）}
DiffSplat~\cite{2501_16764} 是从“逐样本优化”转向“模型一次性生成高斯”的重要锚点。DreamGaussian 和 GaussianDreamer 仍然要对每个提示词单独优化一套 3DGS；LGM、GRM 等前馈方法虽然快，但通常依赖上游多视图图像生成，若这些多视图不一致，后面的重建也会继承裂缝和漂移。DiffSplat 的改变是把高斯属性整理成图像扩散模型能够处理的结构化网格，让大规模文本到图像扩散模型的先验直接服务于 3DGS 生成。

它先构建训练用的“splat grid”。给定若干带位姿的多视图 RGB 图像，轻量重建网络从每个像素回归一个高斯基元的属性，包括颜色、深度或位置、尺度、旋转和不透明度；位置由像素、深度和相机参数反投影得到。为了让这个回归更稳定，训练时还把坐标图和法线图作为辅助几何输入，但这些辅助图只用于训练重建器，不是最终生成阶段必须给用户准备的输入。回归出的多视图 splat grid 仍然可以被 3DGS 光栅化器从任意视角渲染，因此它不是普通 2D 特征图，而是一组带相机几何含义的高斯属性图。

第二步是把 splat grid 放进扩散潜空间。DiffSplat 不把 12 维左右的高斯属性粗暴拆成几张 RGB 图，而是微调图像扩散模型的 VAE，使它能编码和解码高斯属性网格；初始化时复用预训练 VAE 的卷积权重，再扩展通道以适配高斯属性。生成模型随后在 splat latent 上训练，损失由两部分组成：标准扩散损失让 latent 分布学会“像图像扩散一样生成”，任意视角的渲染损失让解码后的高斯在 3D 空间保持一致。这样，模型推理时可以从文本或图像条件采样 splat latent，解码为高斯网格，再直接渲染为 3DGS 资产。

这条路线的价值在于兼容性。因为底座仍是 SD、SDXL、PixArt 等图像扩散结构，ControlNet、IP-Adapter、文本条件和图像条件等 2D 生成技术都可以迁移到 3D 高斯生成中。相对 DiffusionGS 把高斯输出写进 denoiser 的每个去噪步，DiffSplat 更强调“高斯属性图像化 + VAE 潜空间 + 渲染一致性损失”的组合；它适合作为 2D 扩散先验迁移到高斯 latent 的锚点。
\textbf{锚点细化：} DiffSplat：的策略是直接利用大规模文本到图像扩散模型的丰富2D先验，其核心发现在于：图像扩散模型已隐式学习了3D几何估计能力（如深度、法线预测），因此可以将3D高斯属性视为"特殊风格的图像"来生成。关键设计是将3D高斯组织为多视角泼溅网格（multi-view splat grids），每个像素对齐一个高斯基元的全部属性，网格结构与2D图像兼容。为适配预训练图像扩散模型的潜在空间，方法引入"泼溅潜变量"（splat latents）：将高斯属性的多个通道通过VAE压缩到与图像潜变量相同维度的空间，仅需复制预训练VAE卷积权重的前4个通道即可完成初始化。训练流程分两步：先训练一个仅42M参数的轻量级Transformer重建模型（不到LGM/GRM的10\%），在0.1秒内从多视角图像回归高斯网格，用于在G-Objaverse的265K物体上高效构建3D数据集；然后在泼溅潜变量上微调各种预训练扩散模型（SD1.5、SDXL、PixArt-$\Sigma$、SD3），同时使用扩散损失和3D渲染损失确保跨视角一致性。整个训练仅需8 A100 GPU约3天。由于与图像扩散模型兼容，DiffSplat能够无缝迁移ControlNet、IP-Adapter等2D技术到3D生成。
\textbf{实现细节：} DiffSplat（2025.01） 的策略是直接利用大规模文本到图像扩散模型的丰富2D先验，其核心发现在于：图像扩散模型已隐式学习了3D几何估计能力（如深度、法线预测），因此可以将3D高斯属性视为"特殊风格的图像"来生成。关键设计是将3D高斯组织为多视角泼溅网格（multi-view splat grids），每个像素对齐一个高斯基元的全部属性，网格结构与2D图像兼容。为适配预训练图像扩散模型的潜在空间，方法引入"泼溅潜变量"（splat latents）：将高斯属性的多个通道通过VAE压缩到与图像潜变量相同维度的空间，仅需复制预训练VAE卷积权重的前4个通道即可完成初始化。整个训练仅需8 A100 GPU约3天。由于与图像扩散模型兼容，DiffSplat能够无缝迁移ControlNet、IP-Adapter等2D技术到3D生成。
\textbf{设计动机与边界：} 相对原始 3DGS 以独立高斯和图像重建误差为主的设定，DiffSplat 把问题落在以下可实现的改动上：DiffSplat的三大核心组件：(1) 数据集构建 - 结构化溅射重建：受通用3DGS重建技术启发，用一组结构化多视图高斯溅射网格表示3D物体。给定V张位姿图像，小网络F\_regress在0.1秒内回归逐像素溅射属性（颜色、位置、尺度、旋转、不透明度）。引入坐标图和法线图作为辅助几何引导。(2) 溅射潜在空间：将3D高斯溅射属性视为特殊风格的图像，可直接用预训练图像扩散模型的VAE编码为潜在表示。(3) 生成模型：在潜在空间中训练扩散模型，结合标准扩散损失和3D渲染损失（从任意视角渲染并计算图像损失）确保3D一致性。兼容各种图像扩散模型架构（SD/SDXL/PixArt等）。其设计目的不是单独增加一个模块，而是用上述表示、监督或渲染约束消除原流程在目标任务中的具体失配。提出DiffSplat框架，通过驯服大规模文本到图像扩散模型原生生成3D高斯溅射。提出轻量级重建模型即时生成多视图高斯溅射网格用于数据集构建。结合扩散损失和3D渲染损失确保3D一致性。无缝适配图像扩散模型的丰富技术到3D领域。适用时还应保留论文给出的边界：生成质量受限于2D扩散模型的3D几何理解能力。结构化溅射重建可能损失细节。训练需要GPU资源。
\textbf{方法信息：} ByteDance 参与了这篇论文；代码：\url{https://github.com/chenguolin/diffsplat}；GitHub 关注度高，约 530 stars；引用很多，Semantic Scholar 约 53 次。



% METHOD: 2502.01846
\paragraph{UVGS（把无结构高斯变成 UV 图）（2025.02）}
UVGS~\cite{2502_01846} 解决 3DGS 无结构导致难以直接使用 2D 生成模型的问题。DiffSplat 把高斯属性做成多视图 splat grid，再训练 VAE 和扩散模型；UVGS 的改变更像把对象高斯重新排版：先把 3DGS 对象放进包围球，将每个高斯中心从笛卡尔坐标转成球坐标，再用方位角和极角落到 2D UV 网格。每个 UV 像素不只是颜色，而是存储位置、旋转、缩放、不透明度和颜色等高斯属性，因此这张图仍然能被解码回可渲染的 3DGS。

为了让这种高维属性图能被图像模型处理，UVGS 再用多分支网络分别处理几何、外观和不透明度等异构字段，并压缩到低维共享特征空间。这样一来，标准 VAE、latent diffusion 和 2D 图像修补模型可以直接在 UV 表示上工作：生成阶段采样或编辑的是 UV 图，渲染阶段再把 UV 像素还原为高斯集合。相对 DiffSplat 的 splat latent，UVGS 的关键增量是显式建立“高斯集合到 2D 参数域”的可逆组织方式，适合对象级生成、条件生成和局部修补。
\textbf{实现细节：} LucidDreamer（2023.11） 与GaussianDreamer同为SDS优化方法，但核心差异在于：GaussianDreamer聚焦物体级生成并依赖3D扩散模型初始化，LucidDreamer则面向场景级生成且完全不需要3D训练数据。动机。场景级生成比物体级更具挑战性——场景范围大、内容多样，无法用简单几何初始化；LucidDreamer 通过"边走边看"的渐进扩展策略，利用扩散模型的 inpainting 能力逐步填充未见区域。方法提出一种域无关的管线，交替执行两个步骤。（1）Dreaming：以已有3D点云为几何引导，通过扩散模型（Stable Diffusion）的inpainting从新的相机视角生成多视角图像，再用单目深度估计（如ZoeDepth）将像素反投影到3D空间，形成新的点云。直觉上，这像是"边走边看"，从一个初始视角出发，逐步向未观测方向扩展场景。（2）Alignment：将新生成的3D点云与已有场景进行对齐融合，通过相机轨迹规划确保扩展方向的连贯性。重复Dreaming-Alignment循环直到场景覆盖足够大的范围后，将完整3D点云作为3DGS优化的初始化，通过3DGS的$\alpha$-混合渲染来填充点云中因深度估计误差产生的空隙。LucidDreamer生成的3DGS场景细节丰富且无领域约束。
\textbf{方法信息：} Meta 参与了这篇论文。
\textbf{设计动机与边界：} 相对原始 3DGS 以独立高斯和图像重建误差为主的设定，UVGS 把问题落在以下可实现的改动上：UVGS（UV Gaussian Splatting）通过球面映射将3DGS的无结构点集转换为结构化的2D表示。具体流程：（1）将3DGS对象内切到以其中心为球心的包围球中；（2）将每个高斯的位置从笛卡尔坐标转换为球坐标，通过方位角和极角映射到2D UV图的网格位置；（3）每个UV像素存储对应高斯的所有属性（位置、旋转、缩放、不透明度、颜色），形成14通道的图像式UV图。为降低维度，设计了多分支网络将异构特征压缩到低维（如3通道）共享特征空间。压缩后的UVGS可直接作为RGB图像处理。关键发现是标准VAE（配合潜在扩散模型训练）无需额外训练即可直接泛化到这种新表示。通过增加UV分辨率可容纳更多高斯，具有良好的可扩展性。该表示立即解锁了3DGS的无条件生成、条件生成和修补等应用。其设计目的不是单独增加一个模块，而是用上述表示、监督或渲染约束消除原流程在目标任务中的具体失配。提出UVGS表示，通过球面映射将无结构的3DGS转换为结构化的2D UV图，使现有的2D扩散模型可以直接用于3DGS生成，无需额外训练即可实现无条件/条件生成和修补。适用时还应保留论文给出的边界：球面映射对非封闭或非凸对象可能产生变形；UV分辨率与高斯数量的平衡需要手动调整；生成质量受限于2D扩散模型的能力。
% METHOD: 2502.16475
\paragraph{Dragen3D（种子点控制的 Anchor-GS 生成）（2025.02）}
Dragen3D~\cite{2502_16475} 不把重点放在复用 2D 扩散先验，而是把“可控几何骨架”放进 3DGS 生成。DiffSplat 和 Zero-1-to-G 主要关心如何把高斯属性转成扩散模型能生成的表示；Dragen3D 先训练 Anchor-GS VAE，将 3D 几何和外观压成一组 anchor latents，并保证这些 latent 可以解码回 3DGS。随后，单图分支只需要生成稀疏 seed points 作为粗几何骨架，再通过 Seed-Anchor 映射把 seed points 送入 anchor latent 空间。

生成阶段使用 Rectified Flow 学习从 seed-point latent 到 anchor latent 的映射。因为 seed points 是稀疏、低维、可解释的几何控制点，它们比完整高斯集合更容易保持整体形状；用户拖拽 seed points 时，变化通过 anchor latents 传播到最终高斯，从而实现几何编辑。相对纯自动生成方法，Dragen3D 的增量是把 3DGS 的生成和交互式形变放在同一套 anchor 表示中。
\textbf{实现细节：} 可编辑和可控制路线。Dragen3D（2025.02） 引入锚点-高斯变分自编码器（Anchor-GS VAE）和种子点驱动策略，使用户能直接在3D空间编辑生成结果的几何形状。具体流程为：首先训练 Anchor-GS VAE 将复杂3D信息编码为紧凑隐空间表示；然后从输入图像生成一组稀疏3D种子点作为粗几何骨架；接着通过整流流模型（Rectified Flow）将种子点信息映射到 Anchor-GS VAE 的隐空间，解码为3DGS。用户通过拖拽种子点即可改变最终生成的3D形状——种子点的变形会更新对应的隐编码，从而产生几何一致的变形输出。这种"稀疏控制点驱动"的交互方式极大地降低了3D内容编辑的门槛。3DisGS（2025.04） 首次实现了无监督可解释的3DGS重建。该方法提出双分支重建架构：几何分支使用风格化折叠模块（style-based folding module）从2D网格生成3D点云，外观分支基于三平面（triplane）表示生成纹理特征，两者结合实现粗粒度的几何-外观解耦。在细粒度层面，3DisGS 通过解耦表示学习（DRL）的编码器-适配器模块，自动发现3D场景中潜在语义因素（如"家具""装饰""人物"等）的独立表示。损失函数中引入互信息最小化，确保各语义因素之间相互独立。
\textbf{方法信息：} Tencent 参与了这篇论文。
\textbf{设计动机与边界：} 相对原始 3DGS 以独立高斯和图像重建误差为主的设定，Dragen3D 把问题落在以下可实现的改动上：Dragen3D包含三个核心组件：(1)Anchor-GS VAE：将3D几何和外观编码为紧凑的anchor latents，可解码回3DGS。VAE的训练使得在隐空间进行高效生成成为可能。(2)种子点驱动策略：从单张图像生成稀疏种子点作为粗几何表示，由于种子点稀疏分布容易学习，确保了几何一致性。种子点通过Seed-Anchor映射模块映射到VAE的隐空间。(3)种子点驱动变形：用户可直接拖拽种子点来变形最终3DGS的几何，变化通过anchor latents传播。使用Rectified Flow Model建立种子点隐空间到anchor隐空间的映射。整个过程不依赖2D扩散先验，实现多视图几何一致的3D生成。其设计目的不是单独增加一个模块，而是用上述表示、监督或渲染约束消除原流程在目标任务中的具体失配。首次实现不依赖2D扩散先验的几何可控3DGS生成和编辑，通过稀疏种子点实现多视图几何一致性和用户拖拽控制。适用时还应保留论文给出的边界：单图输入对遮挡区域的推理仍有限；种子点数量和分布影响生成质量；复杂拓扑结构的处理需改进。
% METHOD: 2503.13272
\paragraph{Generative Gaussian Splatting（视频扩散特征场生成）（2025.03）}
Generative Gaussian Splatting~\cite{2503_13272} 将 3DGS 表示直接集成到 latent video diffusion。它发现，在 feature space 而非 latent space 学习 3D 表示，更能兼容视频扩散的 v-prediction 训练和三维一致性。方法从 3DGS 参数化的特征场渲染特征图，再映射到 latent space 生成多视图；渲染的特征图也可解码成 3D 场景。相对 Lyra 的 RGB 分支自蒸馏，GGS 更强调在视频扩散内部插入显式 3D feature field，避免生成多视图之间几何漂移。
\textbf{方法信息：} Meta 参与了这篇论文。
\textbf{设计动机与边界：} 相对原始 3DGS 以独立高斯和图像重建误差为主的设定，Generative Gaussian Splatting 把问题落在以下可实现的改动上：GGS方法的核心设计：(1) 位姿条件图像到视频架构：基于预训练U-Net视频扩散模型，条件化于Plucker嵌入的相机位姿和多个参考图像。仅对目标视图应用前向扩散，参考视图以零噪声级连接。(2) 3D约束集成：在特征空间而非潜在空间学习3D表示。3DGS参数化的特征场渲染特征图，通过解码器映射到潜在空间。这解决了两个挑战——潜在空间本身不是3D一致的，且v-预测框架需要建模噪声而非去噪输入。(3) 双路径生成：渲染的特征图既可解码为多视图图像，也可直接上采样为3D辐射场。额外训练一个解码器直接从生成的特征图预测解码后的3D场景。生成后可用标准3D重建算法进一步细化。其设计目的不是单独增加一个模块，而是用上述表示、监督或渲染约束消除原流程在目标任务中的具体失配。提出GGS方法，首次将3DGS表示直接集成到预训练潜在视频扩散模型中。核心发现：在特征空间而非潜在空间学习3D表示可以同时兼容v-预测训练和有意义的3D表示生成。GGS从特征场渲染特征图再解码到潜在空间，显著改善生成多视图图像的3D一致性。相比无3D表示的类似模型，FID在RealEstate10K和ScanNet+上改善约20%。适用时还应保留论文给出的边界：训练需要大量GPU资源；生成质量仍受限于视频扩散模型的能力；对复杂几何（如细小结构）的重建可能不够精确。
% METHOD: 2508.01464
\paragraph{Can3Tok（场景级高斯 token 化）（2025.08）}
Can3Tok~\cite{2508_01464} 面向 scene-level 3DGS 潜空间。它用 PerceiverIO 式交叉注意力，把大量高斯编码为固定数量 canonical tokens；编码时对高斯中心做 Fourier 位置编码，并加入结构化体素 anchors 减轻学习负担。VAE 的 KL 正则让 token 成为可生成 latent，解码器再由 latent 交叉注意力恢复高斯参数。相对 UVGS 的 2D 图重排，Can3Tok 面向复杂场景和数量很大的高斯集合，目标是给文本到场景或图像到场景生成一个可扩展 token 空间。
\textbf{实现细节：} Can3Tok（2025.08） 是首个面向场景级3DGS的VAE。场景级生成面临独特挑战：不同场景尺度不一致（COLMAP的SfM初始化使得全局场景尺度和每个高斯基元的缩放值均不具有度量意义），高斯数量差异巨大（单个场景可达数十万基元），且3DGS数据本质上是高度非结构化的——异构特征（几何、外观、光照混杂）与点云式的不规则性并存。Can3Tok从架构设计和数据处理两个层面系统性地解决了这些难题。架构设计。Can3Tok采用基于Transformer的VAE编码器-解码器结构。编码器的输入是经过傅里叶位置编码（Fourier positional encoding，最大频带 $L_B=51$）的高斯基元参数 $G \in \mathbb{R}^{N \times (2L_B+C)}$，同时为每个高斯附加以其最近体素中心坐标的傅里叶编码作为结构化锚点（体素分辨率 $V=40$），引入规整的几何先验以降低编码器学习负担。由于单个场景的高斯数量 $N$ 超过万级，直接自注意力计算代价过高，编码器采用PerceiverIO式的交叉注意力：输入经线性层映射为$key$和$value$，与一个可学习规范查询（learnable canonical query）进行交叉注意力运算。该规范查询初始化为规则体素网格坐标及其描述子 $\mathrm{query} \in \mathbb{R}^{M \times (P+Q)}$，其中 $M$ 为规范体素数量，$P$ 为位置维度，$Q$ 为描述子维度（$Q=256$，隐层维度768），并在训练过程中持续优化。这种设计不仅大幅降低计算开销（将 $O(N^{}2)$ 自注意力降为 $O(NM)$ 交叉注意力），还在非结构化3DGS数据上引入了结构化的几何先验。交叉注意力输出后，经8个自注意力块进一步挖掘潜在标记之间的亲和关系。
\textbf{方法信息：} Adobe Research 参与了这篇论文；仓库：\url{https://github.com/zerg-overmind/can3tok}，当前 GitHub 约 42 stars，但核验时仓库根目录主要是 README、LICENSE 和图片目录，完整可运行代码暂未确认。
\textbf{设计动机与边界：} 相对原始 3DGS 以独立高斯和图像重建误差为主的设定，Can3Tok 把问题落在以下可实现的改动上：Can3Tok架构包含编码器和解码器：(1)编码器：对3DGS中心应用傅里叶位置编码，附加结构化体素锚点减少编码器学习负担。通过线性层映射输入为key/value，使用PerceiverIO风格的交叉注意力将10K+高斯基元标记化为固定数量的规范token（规范查询由规则体素网格初始化）。潜在空间通过VAE的KL散度正则化；(2)解码器：从潜在表示通过交叉注意力恢复高斯基元参数；(3)数据处理管线：解决场景级3DGS的尺度不一致问题，包括归一化和统一坐标系处理。整个框架训练在DL3DV-10K数据集上，学习通用的场景级潜在空间，可对接扩散模型实现条件生成。其设计目的不是单独增加一个模块，而是用上述表示、监督或渲染约束消除原流程在目标任务中的具体失配。提出Can3Tok，首个3D场景级VAE，能将大量高斯基元编码为低维潜在嵌入，有效捕获语义和空间信息。设计通用3D场景数据处理管线解决尺度不一致问题。在DL3DV-10K上验证，对比方法无法收敛而Can3Tok成功泛化到新场景。展示了图像到3DGS和文本到3DGS生成应用。适用时还应保留论文给出的边界：重建精度受VAE压缩瓶颈限制；场景级生成的细节质量仍有提升空间；潜在空间离散化可能损失高频细节。
% METHOD: 2509.19296
\paragraph{Lyra（视频扩散自蒸馏到 3DGS）（2025.09）}
Lyra~\cite{2509_19296} 不依赖真实多视角训练数据，而是把视频扩散模型中的隐含 3D 知识蒸馏到 3DGS 解码器。它在标准 RGB decoder 旁增加 3DGS decoder 分支，同一个视频 latent 同时解码成 RGB 视频和高斯场景；3DGS 分支由 RGB 分支生成的视频监督训练。这样 3D 解码器可以完全使用视频扩散模型产生的合成数据训练。相对 DiffSplat 需要高斯数据集，Lyra 的增量是从视频模型自蒸馏三维表示。
具体流程建立在 camera-controlled video diffusion model 上。训练时冻结预训练视频扩散模型和 RGB decoder，只训练新增的 3DGS decoder：给定文本、单图或单目视频条件，扩散模型先产生带相机轨迹的视频 latent，RGB 分支把它解码成多帧图像；3DGS 分支从同一 latent 直接预测高斯的位置、尺度、旋转、不透明度和外观特征。监督不是高斯参数真值，而是把这些高斯沿同一相机轨迹渲染回图像，再与 RGB 分支输出的视频帧对齐，相当于用视频模型自己的多视角结果提供重建损失和跨视角约束。由于监督来自视频模型自身，Lyra 不需要真实 3DGS 数据集，也不需要真实多视角采集。推理时可以跳过 RGB 分支，直接由 3DGS decoder 输出静态 3D 场景；输入是视频时，分支进一步输出动态高斯表示，把时变几何和外观放进 4DGS。相对 GGS 在视频扩散内部插入特征场，Lyra 更像把视频模型的 2D/多视角想象力蒸馏成一个前馈高斯解码器。
\textbf{锚点细化：} Lyra：提出从视频扩散模型中蒸馏隐式3D知识到显式3DGS表示，消除了对多视角训练数据的需求。方法基于GEN3C视频扩散模型构建师生框架：教师模型（视频扩散模型）生成相机控制的RGB视频，学生模型（3DGS解码器）在同一潜空间中操作，由教师的RGB输出进行监督。3DGS解码器将视频VAE的潜变量解码为3D高斯参数，整个训练仅需视频扩散模型生成的合成数据，无需任何真实3D或多视角标注。推理时，模型可从文本提示或单张图像合成3D场景并实时渲染。框架还扩展到从单目输入视频生成动态3D场景，展示时间维度的生成能力。
\textbf{实现细节：} Lyra（2025.09） 提出从视频扩散模型中蒸馏隐式3D知识到显式3DGS表示，消除了对多视角训练数据的需求。方法基于GEN3C视频扩散模型构建师生框架：教师模型（视频扩散模型）生成相机控制的RGB视频，学生模型（3DGS解码器）在同一潜空间中操作，由教师的RGB输出进行监督。3DGS解码器将视频VAE的潜变量解码为3D高斯参数，整个训练仅需视频扩散模型生成的合成数据，无需任何真实3D或多视角标注。推理时，模型可从文本提示或单张图像合成3D场景并实时渲染。框架还扩展到从单目输入视频生成动态3D场景，展示时间维度的生成能力。自蒸馏架构。Lyra的核心思路是利用相机控制视频扩散模型 $V$ 的3D一致性来监督3DGS解码器 $D_s$，形成自蒸馏循环。教师模型 $V$ 基于GEN3C构建：它首先从输入图像构建时空3D缓存（将RGB图像通过深度估计反投影为彩色点云），然后将缓存从任意相机姿态渲染为RGB图像和遮挡掩码，作为结构化视觉引导。视频VAE将RGB视频 $x \in \mathbb{R}^{3 \times H \times W}$ 编码到潜空间 $z = \mathcal{E}(x) \in \mathbb{R}^{\tau \times C \times h \times w}$，其中潜通道维度 $C=16$，时间压缩因子 $\kappa_t=8$，空间压缩因子 $\kappa_s=8$。给定输入图像 $I$ 和采样的相机轨迹 $\{Camera^{v}\}_{v=1}^{V}$，教师模型生成去噪后的视频潜变量 $z^v = V(I, \{Camera^{v}\}_{v=1}^{V})$。潜变量沿两条分支解码：预训练RGB解码器 $\mathcal{D}_r$ 生成RGB帧 $x^v_r = \mathcal{D}_r(z^v)$，3DGS解码器 $D_s$ 输出显式3D高斯 $\mathcal{G}$。学生模型的渲染视角 $\hat{x}^v_r = \mathrm{Render}(\mathcal{G}, \{Camera^{v}\}_{v=1}^{V})$ 被监督以匹配教师的RGB帧 $x^v_r$，形成自蒸馏闭环。为扩大视角覆盖范围，每个输入图像采样 $V=6$ 条相机轨迹，每条轨迹包含 $L=121$ 个相机姿态，总计726个输入视角。基于潜空间的3DGS解码器。此前的前馈重建方法（如GS-LRM、AnySplat）在像素空间操作，受限于GPU内存仅能处理少量低分辨率视角。
\textbf{设计动机与边界：} 相对原始 3DGS 以独立高斯和图像重建误差为主的设定，Lyra 把问题落在以下可实现的改动上：Lyra框架的核心是在视频扩散模型的潜在空间中直接操作3DGS解码器。具体设计：（1）架构：在标准RGB解码器旁并行增加3DGS解码器分支。给定单张图像或视频输入，采样相机轨迹作为条件，去噪视频潜在表示，然后通过两个并行分支解码。RGB分支合成视频序列，3DGS分支直接从潜在表示生成3DGS场景表示。（2）自蒸馏训练：3DGS解码器由RGB解码器的输出监督，无需真实3D数据。这意味着3DGS解码器可以完全用视频扩散模型生成的合成数据训练。（3）扩展到动态场景：框架进一步支持从单目输入视频生成动态3D场景，通过4D高斯表示建模时变几何和外观。推理时支持从文本提示或单图生成3D场景并进行实时渲染。其设计目的不是单独增加一个模块，而是用上述表示、监督或渲染约束消除原流程在目标任务中的具体失配。提出自蒸馏框架，将视频扩散模型中隐含的3D知识蒸馏到显式3DGS表示中，消除对多视角训练数据的需求。在视频扩散模型中增加3DGS解码器分支，通过RGB解码器的输出监督训练，支持从文本或单图生成静态和动态3D场景。适用时还应保留论文给出的边界：生成质量受限于视频扩散模型的能力；对复杂几何结构的恢复可能不够精确；大规模场景生成能力未充分验证。
\textbf{方法信息：} NVIDIA 参与了这篇论文；代码：\url{https://github.com/nv-tlabs/lyra}；GitHub 关注度很高，约 2.1k stars。


% METHOD: 2601.19247
\paragraph{TIGaussian（三模态高斯对齐）（2026.01）}
TIGaussian~\cite{2601_19247} 不只是生成模型，也是文本-图像-3D 对齐框架。它用多分支 3DGS tokenizer 解耦几何属性和外观属性，先采样和分组高斯局部 patch，再编码融合为结构化 3D embedding。图像侧通过扩散增强多视图融合补偿单视图 3D 感知不足；文本侧通过 3D-to-text projection 对齐文本嵌入，并用 InfoNCE 训练。相对 UniGS 的预训练框架，TIGaussian 更强调 3DGS tokenizer 和扩散增强图像特征。
\textbf{实现细节：} TIGaussian（2026.01） 面向视觉-语言-3D跨模态对齐，提出多分支3DGS标记器将高斯的内在属性解耦为结构化潜在嵌入 $F_G^I$：分别处理几何属性（位置、缩放、旋转）和外观属性（不透明度、颜色），而非将所有属性混合编码。配合双向跨模态对齐策略：(1) 图像-3D对齐利用扩散先验的多视角特征融合解决视角模糊问题，生成3D感知视觉特征；(2) 文本-3D对齐通过自适应投影模块将空间特征 $F_G^I$ 映射到文本嵌入空间 $F_G^T$，使用InfoNCE对比损失联合优化。在跨模态检索、零样本分类和场景识别等多个任务上达到最优性能。问题动机与现有方法局限。
\textbf{方法信息：} Alibaba 参与了这篇论文；仓库：\url{https://github.com/ruin-jiarun/tigaussian}，当前约 5 stars，核验时主要是 README 和 LICENSE，完整代码暂未确认。
\textbf{设计动机与边界：} 相对原始 3DGS 以独立高斯和图像重建误差为主的设定，TIGaussian 把问题落在以下可实现的改动上：框架包含三个协调组件：(1)多分支3DGS分词器：将高斯的几何属性（位置、尺度、旋转）和外观属性（颜色、SH系数、不透明度）分离开来，使用FPS+KNN采样至1024个高斯后分组为局部patch，分别编码再融合，减少属性间耦合，构建结构化隐嵌入F\_G\^{}I。(2)扩散增强多视图融合：利用预训练扩散模型的隐式3D先验，补偿单视图图像的3D感知不足，生成3D感知视觉特征F\_I\^{}mv。(3)3D-文本投影：通过投影模块将空间特征F\_G\^{}I映射到F\_G\^{}T，确保与文本嵌入的模态不变一致性。使用InfoNCE对比损失分别训练3D-文本和3D-图像对齐。其设计目的不是单独增加一个模块，而是用上述表示、监督或渲染约束消除原流程在目标任务中的具体失配。提出TIGaussian三模态对齐框架，利用3DGS特性增强跨模态对齐。核心包含：(1)多分支3DGS分词器将高斯属性解耦为紧凑隐表示；(2)扩散增强多视图融合机制解决图像-3D对齐中的视角模糊；(3)3D-文本投影模块自适应映射3D特征到文本嵌入空间。在零样本分类、跨模态检索、少样本探测和场景识别上达到SOTA。适用时还应保留论文给出的边界：FPS+KNN下采样可能丢失细粒度3D信息；对大规模场景的处理效率待优化；多分支分词器的分支设计依赖手工选择。
% METHOD: 2601.04090
\paragraph{Gen3R（重建 latent 与视频扩散联合生成）（2026.01）}
Gen3R~\cite{2601_04090} 是 2026 年扩展召回中补入的相邻生成路线。它不直接以 3DGS 为输出，而是把前馈重建模型的几何先验和视频扩散模型的外观先验放到同一个生成过程里：VGGT 这类重建模型提供几何 token，adapter 把这些 token 变成 geometric latents；预训练视频扩散模型提供 appearance latents；训练时用正则让两类 latent 对齐，使模型在生成 RGB 视频时也能同步输出相机位姿、深度图和全局点云。
实现上，Gen3R 的输入可以是单图或多图条件，输出不是单纯视频，而是一组彼此配准的 RGB 帧、camera poses、depth maps 和 global point cloud。它的模块关系更接近“生成模型 + 前馈重建器”的耦合：视频扩散负责补全外观和未观测区域，VGGT token 经过编码和 adapter 后约束几何结构，联合采样阶段同时生成外观 latent 与几何 latent。放在本节的原因是，它和 Lyra、CamPilot、GGS 处在同一条从视频扩散走向显式 3D 输出的路径上；差别是 Gen3R 当前证据指向点云和深度等几何输出，而不是直接输出可渲染的高斯参数。
\textbf{方法信息：} 作为 2026 年世界生成扩展召回条目纳入；相对原始 3DGS，它缺少可直接光栅化的高斯参数输出和 densification/渲染优化闭环，因此本文只把它作为相邻生成/重建路线保守处理；未发现本地可靠代码或外部高关注度证据。
\textbf{设计动机与边界：} 相对原始 3DGS 以独立高斯和图像重建误差为主的设定，Gen3R 把问题落在以下可实现的改动上：方法从 VGGT token 中训练 adapter 得到几何 latent，并将其正则化到与预训练视频扩散模型的外观 latent 对齐；生成阶段同时采样外观与几何 latent，输出 RGB 视频以及相机位姿、深度图和全局点云。它不是 3DGS 方法，但与 3DGS 生成章节中的 Lyra、GGS 等前馈生成路线相邻。其设计目的不是单独增加一个模块，而是用上述表示、监督或渲染约束消除原流程在目标任务中的具体失配。把基础重建模型和视频扩散模型耦合起来，使生成先验和几何重建互相约束。适用时还应保留论文给出的边界：未显式使用 3DGS 或 Gaussian Splatting，作为相邻生成/重建路线保留。
% METHOD: 2601.16214
\paragraph{CamPilot（相机奖励反馈中的 3D Gaussian 解码）（2026.01）}
CamPilot~\cite{2601_16214} 解决 camera-controlled video diffusion 的相机控制不稳问题。普通奖励反馈如果每次都把 latent 解码成 RGB 视频再评估相机一致性，计算开销很高，而且奖励通常只看 2D 帧，缺少几何判断。CamPilot 的改变是在视频 latent 后接 camera-aware 3D decoder：输入是视频 latent 和相机 pose，输出是一组 3D Gaussians，再把这些高斯从目标相机渲染成 novel views，用渲染结果来评估相机控制是否正确。
具体训练流程把相机 pose 同时作为条件和投影参数。如果视频 latent 与相机轨迹不一致，3D Gaussian decoder 生成的几何在新视角渲染时会出现模糊、错位或结构变形；这些渲染误差就可以转化为 reward。奖励由 novel-view 像素一致性和可见性项组成，可见性项只监督由几何 warping 确定的区域，避免遮挡和未观测部分把反馈带偏。这样，优化目标不再只是“视频看起来像”，而是让 latent、相机和可渲染几何在同一坐标关系下闭合。相对 Lyra 把视频扩散蒸馏成 3DGS 输出，CamPilot 更像把 3DGS decoder 用作相机控制奖励模型。
\textbf{方法信息：} 作为 2026 年世界生成扩展召回条目纳入；当前本地证据未确认完整公开代码。
\textbf{设计动机与边界：} 相对原始 3DGS 以独立高斯和图像重建误差为主的设定，CamPilot 把问题落在以下可实现的改动上：方法在视频扩散模型的 latent 空间后接 camera-aware 3D decoder：输入视频 latent 与相机 pose，输出 3D Gaussian 表示；相机 pose 同时作为条件和投影参数，若 latent 与相机不一致，渲染的新视角会出现模糊或几何失真。训练时用渲染 novel views 与真值视图之间的像素一致性作为 reward，并加入可见性项只监督几何 warping 得到的确定区域。其设计目的不是单独增加一个模块，而是用上述表示、监督或渲染约束消除原流程在目标任务中的具体失配。把相机对齐误差转化为 3D Gaussian 渲染质量和像素一致性奖励，避免每次奖励评估都解码完整 RGB 视频。适用时还应保留论文给出的边界：主要目标是视频扩散相机控制，不是独立 3DGS 场景重建系统；代码状态未核验。
% METHOD: 2602.21929
\paragraph{Geometry-as-context（显式几何作为视频生成上下文）（2026.02）}
Geometry-as-context~\cite{2602_21929} 不是直接 3DGS 方法，而是面向 scene-consistent video generation 的相邻路线。它关注的问题和 Lyra 2.0 接近：相机沿轨迹探索场景时，视频生成模型容易忘记历史区域或在自回归过程中累积错误。它的做法是在每一步生成中先估计当前视图重建所需的显式几何，再用这些几何作为上下文去模拟和恢复新视角图像，让视频模型不只依赖短窗口内的 RGB 历史。
实现流程是自回归的：给定相机轨迹和已有帧，模型先输出当前视图的几何信息，再根据该几何渲染或模拟 novel-view 线索，随后恢复目标 RGB 帧。训练阶段用文本上下文区分“当前要输出几何还是 RGB”，并随机丢弃 geometry context，使模型在推理时即使上下文不完整也能继续生成。它还引入 camera-gated attention，把相机位姿更明确地注入网络。本文把它作为补充条目，是因为它体现了 2026 年视频世界生成中“显式 3D 上下文”的趋势；但由于摘要和本地证据没有显示其使用 3DGS，高斯参数、渲染器和优化损失不在这里强行补写。
\textbf{方法信息：} 作为 2026 年世界生成扩展召回条目纳入；相对原始 3DGS，它使用显式几何上下文而不是高斯基元、协方差和 splatting 渲染器，本文按相邻视频生成方法保守处理。
\textbf{设计动机与边界：} 相对原始 3DGS 以独立高斯和图像重建误差为主的设定，Geometry-as-context 把问题落在以下可实现的改动上：方法在每一步生成中先估计用于 3D 重建的几何，再模拟并恢复从 3D 场景渲染出的 novel view 图像；训练阶段通过文本上下文区分几何或 RGB 输出，并随机丢弃 geometry context 以保证推理时可以只输出 RGB。其设计目的不是单独增加一个模块，而是用上述表示、监督或渲染约束消除原流程在目标任务中的具体失配。把当前视图重建所需几何作为上下文喂给视频生成模型，并用 camera gated attention 强化相机位姿利用。适用时还应保留论文给出的边界：未显式使用 3DGS；作为视频生成到三维重建的相邻路线而非核心 3DGS 方法。
% METHOD: 2604.00494
\paragraph{ARGS（层次树自回归高斯生成）（2026.04）}
ARGS~\cite{2604_00494} 把 next-scale autoregressive prediction 引入 3DGS 物体生成。DiffSplat 和 Can3Tok 把高斯变成扩散 latent 或 token 序列，但它们仍要处理无结构集合的生成顺序；ARGS 的改变是先把一个训练好的高斯模型压成从粗到细的层次树，再学习“如何把粗节点展开成细节点”。这样生成过程不再是一点一点预测全部高斯，而是沿树的尺度层级并行细化。

具体流程先做 Gaussian simplification：按照协方差体积、透明度和渲染贡献找出影响较小的高斯，把它与邻近高斯用矩匹配方式合并，尽量保持合并前后的均值、协方差和外观统计一致。反复合并会自然形成二叉树，根节点是最粗场景，叶节点接近原始高斯。训练时把这个逆过程作为监督，tree-based Transformer 在每个尺度预测下一层节点的位置、不透明度、SH 特征和协方差；叶节点可以注意到自己的祖先节点，因此局部细节不会脱离粗形状。推理时如果只需要粗预览，可以提前停止在中间尺度；如果需要高保真结果，再继续展开到叶层。相对普通自回归点预测，ARGS 的核心收益是把生成步数从线性点数压到层次数，并且把细节预算变成可控参数。
\textbf{实现细节：} 自回归模型在2D图像上已取得突破性进展（如下一尺度预测范式），将其扩展到3D高斯生成面临两个核心难题：(1) 高斯是无序的3D基元，不存在自然的标记序列；(2) 不同物体的高斯数量差异悬殊，难以统一建模。ARGS（Auto-Regressive Gaussian Splatting）（2026.04） 的解决思路是：先将高斯集简化为一棵层次化二叉树，再反转该过程作为生成路径，从而将无序高斯转化为有序的粗到细生成序列。基于矩守恒的简化原则。方法的起点是高斯的矩性质——零阶矩 $m_0$ 对应总质量（不透明度与协方差行列式的函数），一阶矩 $m_1$ 指向高斯中心。ARGS的简化原则是：在每一步中，选取协方差矩阵行列式最小的高斯（即空间最紧凑、对渲染贡献最小的那个），将其与加权距离最近的邻居高斯融合。之所以选择行列式最小的而非最远的，是因为行列式直接度量高斯的体积大小，体积越小则投影面积越小、视觉影响越低，从而保证简化对渲染质量的损失最小。融合操作本身并不简单叠加属性——为正确处理两个高斯的空间重叠，ARGS引入了一组辅助"交叉高斯"：其不透明度 $\alpha_c = \alpha_1 \alpha_2$，协方差 $\Sigma_c = (\Sigma_1^{-1} + \Sigma_2^{-1})^{-1}$，均值 $\mathbf{u}_c = \Sigma_c(\Sigma_1^{-1}\mathbf{u}_1 + \Sigma_2^{-1}\mathbf{u}_2)$。融合后的零阶和一阶矩分别从两个源高斯减去交叉高斯的贡献（$m_{0,3} = m_{0,1} + m_{0,2} - m_{0,c}$，$m_{1,3} = m_{1,1} + m_{1,2} - m_{1,c}$），新均值由矩比值恢复 $\mathbf{u}_3 = m_{1,3}/m_{0,3}$，协方差和球谐特征则以零阶矩为权重做加权平均。整个简化过程用堆优先队列驱动，复杂度接近线性。从线性序列到层次化树结构。直接拟合简化序列意味着自回归需要 $n-1$ 个串行步骤（每步生成一个高斯），对于动辄数万高斯的物体完全不可行。ARGS的关键改进是将线性简化序列重组为层次化二叉树：以最终融合后的单一高斯为根节点，每次融合操作的逆过程即为一个父节点分裂为两个子节点，递归展开直到所有叶节点均为原始高斯。由此，$n$ 个高斯被组织为 $\mathcal{O}(\log n)$ 层的树，每层包含该尺度下的全部高斯节点。
\textbf{方法信息：} Lightspeed Studios at Tencent 参与了这篇论文；完整公开代码暂未确认。
\textbf{设计动机与边界：} 相对原始 3DGS 以独立高斯和图像重建误差为主的设定，ARGS 把问题落在以下可实现的改动上：核心流程：(1)高斯简化：基于协方差矩阵行列式排序，使用堆优先队列迭代合并最小贡献高斯与最近邻，基于矩加权公式保持渲染保真度；(2)层次树构建：合并过程自然形成二叉树结构，根为最粗糙表示，叶为原始高斯；(3)自回归生成：逆简化过程，从根节点（粗糙）逐步生成到叶节点（精细），只需O(log n)步；(4)树基Transformer：自回归预测树结构，叶节点可关注内部祖先节点增强结构一致性。每步预测新的高斯属性（位置、不透明度、SH特征、协方差）。其设计目的不是单独增加一个模块，而是用上述表示、监督或渲染约束消除原流程在目标任务中的具体失配。提出ARGS框架，首次将自回归下一尺度预测范式成功扩展到3DGS物体生成。层次树结构实现O(log n)步高效生成，树基Transformer保证多尺度结构一致性，可控细节层次和视觉保真度。适用时还应保留论文给出的边界：生成质量依赖简化策略；目前聚焦单个物体；长序列生成可能累积误差
% METHOD: 2604.13036
\paragraph{Lyra 2.0（长轨迹可探索 3D 世界生成）（2026.04）}
Lyra 2.0~\cite{2604_13036} 是 Lyra 系列从“单次前馈生成 3DGS/4DGS 场景”走向“长轨迹可探索 3D 世界”的后续版本。Lyra 1.0 的锚点是把视频扩散模型的 latent 蒸馏到 3DGS decoder；Lyra 2.0 的问题变成：当相机沿很长路径反复进入、离开并重访同一空间时，视频模型会忘记离开时间窗口的区域，也会在自回归生成中逐步累积外观和几何漂移。因此它不只问“如何从一个 latent 前馈出高斯场景”，而是把 camera-controlled video generation、历史几何检索和前馈重建放进同一条世界生成管线。
具体流程先用相机控制视频模型生成场景漫游视频，再用前馈重建模型把视频提升为可实时渲染和仿真的 3D 输出。为了缓解 spatial forgetting，系统为已经生成的帧维护 per-frame 3D geometry，但这些几何不直接替代生成模型的外观合成，而是作为信息路由层使用：给定当前目标视角，先从历史中检索空间上相关的帧，再建立与目标视角的 dense correspondences，把过去看见过的结构重新送回模型上下文。为了缓解 temporal drifting，训练时构造 self-augmented histories，让模型看到自己在长轨迹生成中产生的退化历史，并学习修正这些退化输出，而不是把小误差一路传下去。这样得到的长视频轨迹再用于 fine-tune feed-forward reconstruction models，使最终 3D 场景在大视角变化、位置重访和长距离探索下更稳定。与 GGS 在视频扩散内部插入 3D feature field 相比，Lyra 2.0 更强调视频世界生成的记忆机制；与 Lyra 1.0 相比，它的核心增量是把局部 latent-to-3D 解码扩展为跨时间、跨位置的持久世界生成。
从实现角度看，它的输入不是一组固定多视图图像，而是带相机轨迹、历史帧和目标视角的长序列条件；中间模块先编码历史几何特征，再把检索到的帧投影或对应到当前视角；解码端输出当前时间窗口的视频帧和可交给前馈重建器的多视角证据。训练目标围绕两个约束展开：一类损失维持目标视角和历史重访区域的外观一致性，另一类约束让自增强历史中的退化结果在后续优化中被修正。最终 3D 表示的参数、尺度和可渲染结构由前馈重建分支生成，而不是在每个新场景上长时间逐样本优化。
\textbf{方法信息：} NVIDIA 参与了这篇论文；代码：\url{https://github.com/nv-tlabs/lyra}，官方 Lyra 仓库包含 Lyra 1.0 和 Lyra 2.0；GitHub 关注度很高，约 2.1k stars。arXiv 摘要没有显式写出 3DGS 字样，本文按 Lyra 系列后续版本和官方仓库证据纳入。



\section{文本到 3D 生成}

文本到 3D 的早期锚点是 DreamGaussian 与 GaussianDreamer：用 2D 扩散 SDS 优化高斯，但用 3D 扩散、点云或 mesh 先验帮助几何不塌。后续方法扩展到大场景、PBR 材质、组合对象、植物、城市和全景室内。


% METHOD: 2310.08529
\paragraph{GaussianDreamer（3D 扩散初始化加 SDS）（2023.10）}
GaussianDreamer~\cite{2310_08529} 是 DreamGaussian 之后的初始化锚点。DreamGaussian 证明了 3DGS 适合承载 SDS，但它从球体内随机高斯开始，几何要完全由 2D 扩散梯度拉出来，容易遇到 Janus、多物体结构混乱和局部最优。GaussianDreamer 的核心改变是：SDS 仍然保留，但初始 3DGS 不再随机生成，而是先由 3D 扩散模型给出粗几何。

具体流程分两段。第一段用 Shap-E 这类文本到 3D 扩散模型，或 MDM 这类文本到动作/人体姿态模型，根据提示词生成粗 3D 资产，再从 mesh 或人体模型中采样点云。点云被转成 3D Gaussians 的初始中心；为了避免初始点过稀，方法在已有点附近做 noisy point growing，随机补充新高斯；为了避免所有点颜色过于单调，还对初始颜色做扰动。这样得到的不是最终结果，而是一套已经有大体形状和空间分布的高斯起点。

第二段回到 2D 扩散监督。训练时随机采样相机视角，渲染当前 3DGS，再用 Stable Diffusion 的 SDS 梯度更新高斯的位置、尺度、旋转、不透明度和颜色。由于粗几何已经来自 3D 扩散，SDS 主要负责补外观细节和局部几何，而不是从零组织整体结构。论文还把同一思路用于文本到头像：先用动作或人体相关 3D 先验产生初始姿态，再用 2D 扩散细化外观。相对 DreamGaussian，GaussianDreamer 的增量非常明确：它把“随机初始化 + SDS”改成“3D 扩散点云初始化 + SDS”，用 3D 一致性换取更稳定的优化起点。
\textbf{锚点细化：} GaussianDreamer：修改了DreamGaussian/SDS的初始化环节。它用3D扩散模型的输出替代随机初始化作为3DGS的起点，而SDS优化框架本身保持不变。
\textbf{实现细节：} GaussianDreamer 提出了一种将3D扩散模型的3D一致性与2D扩散模型的丰富细节相结合的策略。方法分两个阶段：首先利用3D扩散模型（如Shap-E）的输出作为3DGS的初始化先验，提供粗糙但全局一致的几何；然后通过SDS优化，用2D扩散模型（Stable Diffusion）逐步丰富几何和外观细节。为增强初始化的多样性，方法引入噪声点生长（noisy point growing）和颜色扰动（color perturbation）操作： 其中$\mathbf{x}_i$是3D扩散模型输出的初始点位置，$\sigma_n$控制扰动幅度。这种两阶段设计在15分钟内即可生成高质量的3D实例或3D虚拟形象，实现了显著的速度提升。
\textbf{设计动机与边界：} 相对原始 3DGS 以独立高斯和图像重建误差为主的设定，GaussianDreamer 把问题落在以下可实现的改动上：GaussianDreamer的框架分为两阶段。第一，3D扩散模型初始化：使用Shap-E（文本到3D）或MDM（文本到动作）等3D扩散模型生成粗略3D实例，从中提取点云作为3D高斯的初始化。引入噪声点增长操作在初始化点周围随机添加新高斯补充细节，颜色扰动操作随机扰动初始颜色增加多样性。第二，2D扩散模型优化：使用SDS（Score Distillation Sampling）损失，通过预训练的2D扩散模型（如Stable Diffusion）优化3D高斯的属性。每个训练步骤随机采样相机姿态渲染图像，计算SDS梯度更新高斯参数。由于3D扩散模型提供了良好的几何先验和3D高斯本身具有的类点云结构先验，优化过程非常高效，15分钟即可完成。其设计目的不是单独增加一个模块，而是用上述表示、监督或渲染约束消除原流程在目标任务中的具体失配。提出GaussianDreamer框架，通过3D Gaussian Splatting桥接2D和3D扩散模型。3D扩散模型（Shap-E/MDM）提供几何一致的粗初始化，2D扩散模型通过SDS损失优化几何和外观细节。引入噪声点增长和颜色扰动操作增强初始化高斯。15分钟内从文本生成高质量3D实例，支持实时渲染。适用时还应保留论文给出的边界：生成质量受3D和2D扩散模型能力限制；SDS损失仍存在多面生成等问题；对复杂文本提示的生成效果有限。
\textbf{方法信息：} Huawei 参与了这篇论文；代码：\url{https://github.com/wuguanjun0627/gaussiandreamer}；引用很多，Semantic Scholar 约 309 次。



% METHOD: 2311.13384
\paragraph{LucidDreamer（Dreaming 与 Alignment 交替扩展场景）（2023.11）}
LucidDreamer~\cite{2311_13384} 把 3DGS 生成从单物体推进到场景级。GaussianDreamer 的对象通常是一个中心物体或头像，优化目标是让一个封闭资产符合文本；LucidDreamer 面对的是室内、室外、动漫、乐高等开放场景，不能依赖固定领域的 3D 扫描训练集。它的路线不是对整场景做一次 SDS，而是像沿相机路径行走一样，不断生成新视角、提升到 3D、再接回已有点云。

输入可以是文本、RGB 或 RGBD。若只有文本，先用 Stable Diffusion 生成第一张图，再用 ZoeDepth 等单目深度模型估计深度；若已有 RGB 或 RGBD，则直接从输入图像和深度开始。第一张 RGBD 被反投影成初始点云。随后沿预设相机轨迹移动相机，把已有点云投影到新视角，已知区域由点云投影给出，未知区域形成 mask。Dreaming 步骤用 Stable Diffusion inpainting 补全这些未知区域，并估计补全图像的深度。

难点在于每次深度估计只有相对尺度，直接反投影会让新点云漂在旧场景前后。LucidDreamer 的 Alignment 步骤根据新旧投影重叠区域估计深度尺度，并在 3D 空间微调新点位置，使新点云和已有点云接起来。这个“投影、修复、估深、对齐、合并”的循环重复多次后，形成大点云；最后再用这份大点云作为 3DGS 初始化，通过高斯优化填补点云空洞并提升连续渲染质量。相对 GaussianDreamer，它不依赖 3D 扩散初始化单个物体，而是把 2D 扩散 inpainting 当作沿路径扩展场景的局部生成器。
\textbf{锚点细化：} LucidDreamer~\cite{2311_13384}针对先前方法受限于特定领域的问题，提出了一种完全无需3D扫描数据训练的域无关场景生成管线。方法交替执行两个步骤：（1）Dreaming——以已有3D点云为几何引导，通过扩散模型的inpainting从新的相机视角生成一致的多视角图像，再通过深度估计将图像提升为3D点云；（2）Alignment——将新生成的3D点云与已有场景进行对齐融合。最终获得的完整3D点云作为3DGS优化的初始化。LucidDreamer生成的3DGS场景细节丰富且无领域约束。
\textbf{方法信息：} 代码：\url{https://github.com/luciddreamer-cvlab/luciddreamer}；GitHub 关注度高，约 1523 stars；引用很多，Semantic Scholar 约 257 次。



% METHOD: 2405.05768
\paragraph{FastScene（全景高斯室内场景生成）（2024.05）}
FastScene~\cite{2405_05768} 从文本快速生成室内 3D 场景。它先用 Diffusion360 生成全景图，并估计全景深度；然后对新相机姿态做粗视角合成，将全景投到球面再重投影到新视角。为了避免直接修复整张全景带来失真，PNVI 逐步修复 cubemap 的六个透视面，再回投到全景。最后通过多视角投影把全景转成透视图，用标准 3DGS 重建。相对 LucidDreamer 的普通透视轨迹，FastScene 利用全景保证全局一致。

具体管线有三个关键转换。第一，文本提示先生成一张 $360^\circ$ 全景图，EGformer 给出对应深度；这一步让房间布局、墙面连续性和大物体位置先在一个全局坐标中确定下来。第二，CVS 根据等距投影把全景像素提升到球面，再从新的相机位置重投影，得到带空洞的新全景视角；PNVI 不直接修补整张全景，而是把全景拆成 cubemap 的六个透视面，逐面 inpainting 后再投回 equirectangular 全景，减少全景边缘模糊和畸变。第三，MVP 把若干干净全景转成 COLMAP/3DGS 能处理的透视图，随后训练标准 3DGS。这个增量的实际意义是把“大场景生成的全局一致性”交给全景，把“3D 可浏览资产”交给 3DGS，而不是沿窄视场轨迹一步步外推。
\textbf{实现细节：} FastScene（2024.05） 的核心贡献在于将全景修复问题从全景图空间转移到立方体贴图空间，从而规避等距柱状投影的严重畸变。其流程分为三步：首先用 Diffusion360 生成全景图并通过 EGformer 估计深度图；然后提出粗视角合成（Coarse View Synthesis, CVS）——根据球面全景的等距投影公式将像素投影到 3D 球面坐标，再根据新相机位姿重新投影获得带孔洞的退化视角；最后提出渐进式新视角修复（Progressive Novel View Inpainting, PNVI），将全景分解为 6 个立方体贴图面分别用重训练的 AOT-GAN 修复，再重投影回全景空间。多次迭代 PNVI 获得多个干净新视角后，通过多视角投影生成透视图像并用 3DGS 重建。整个流程在单张 RTX A6000 GPU 上仅需约 15 分钟（全景生成 10 秒 + PNVI 约 2 分钟 + 训练数据准备约 3 分钟 + 3DGS 优化约 10 分钟），是当时最快的全景到 3DGS 场景生成方法。
\textbf{方法信息：} 代码：\url{https://github.com/mr-ma-yikun/fastscene}。
\textbf{设计动机与边界：} 相对原始 3DGS 以独立高斯和图像重建误差为主的设定，FastScene 把问题落在以下可实现的改动上：框架包含三个阶段。阶段一：给定文本提示，使用Diffusion360生成全景图（利用全景的边界连续性和全局信息），然后使用EGformer估计深度图。阶段二：新视角生成。对给定新相机姿态执行粗视角合成（CVS），根据等距投影将全景投影到球面上并重投影到新视角，产生带空洞的损坏全景。然后提出渐进新视角修复（PNVI）：不直接修复全景而是逐步修复6个透视立方体面视图（cubemap），再重投影回全景。多次迭代PNVI获得多个干净的全景新视角。阶段三：由于COLMAP不支持全景输入，通过多视角投影（MVP）将全景转为透视视图，再用标准3DGS进行场景重建。全景的使用确保了全局一致性，渐进修复保证了视角质量。其设计目的不是单独增加一个模块，而是用上述表示、监督或渲染约束消除原流程在目标任务中的具体失配。提出基于全景高斯溅射的快速3D场景生成框架，仅需15分钟生成3D场景（比SOTA快至少1小时），通过全景的几何约束和渐进修复策略确保全局一致性，在速度和质量上均优于现有方法。适用时还应保留论文给出的边界：依赖全景生成和深度估计的质量。COLMAP不支持全景需转换。对复杂几何结构（如多层空间）的处理有限。
% METHOD: 2406.18462
\paragraph{GaussianDreamerPro（mesh 绑定高斯的三阶段文本生成）（2024.06）}
GaussianDreamerPro~\cite{2406_18462} 继承 GaussianDreamer，但把高斯无控制生长的问题改成 mesh 约束。第一阶段仍用 3D 扩散获得粗资产；第二阶段转成 2DGS，并用 2D 扩散优化，使表面几何更稳定；随后从 2DGS 导出带顶点颜色的 mesh；第三阶段把 3D Gaussians 绑定到 mesh 上继续优化。相对 GaussianDreamer 的自由高斯，Pro 版的增量是“先表面化再绑定”，让文本生成资产更可操控、可动画和可仿真。
三阶段的分工很明确。粗阶段继承 GaussianDreamer 的思路，用 3D 扩散先验给出初始几何和外观，避免完全从随机高斯开始。表面化阶段把资产转为 2DGS，让高斯盘贴近表面并继续接受 2D 扩散监督；这一阶段的目标不是最终渲染，而是导出带顶点颜色的 mesh，让几何有可操作的拓扑载体。绑定阶段再把 3D Gaussians 挂到 mesh 上优化，几何约束限制高斯乱长，SDS/ISM 等扩散损失负责补外观。这样输出既保留 3DGS 的快速渲染，又能进入动画、合成或仿真管线。
\textbf{实现细节：} GaussianDreamerPro（2024.06） 继承GaussianDreamer的3D扩散模型初始化思路，但核心改动是将自由增长的3D高斯绑定到逐步演化的几何结构上。动机。在生成任务中，3D 高斯倾向于无控制地向各个方向增长（因为 SDS 的随机监督使梯度方向不一致），导致生成质量远低于重建质量；将高斯绑定到几何结构可约束其运动范围。设计动机源于一个关键观察：在生成任务中，3D高斯倾向于无控制地向各个方向增长（因为SDS的随机监督使梯度方向不一致），导致生成质量远低于重建质量。在重建中确定性梯度指向唯一最优，而在生成中随机梯度使高斯四处扩散。为解决此问题，方法提出几何引导的三阶段框架。第一阶段（粗几何生成）：继承GaussianDreamer，用3D扩散模型（如Shap-E）生成初始网格，但采用2D高斯（2DGS）而非3D高斯作为表示。2D高斯能更好地贴合表面，提供更准确的几何。第二阶段（基础纹理生成）：在2D高斯上用2D扩散模型通过SDS优化外观和几何，得到基本合格的3D资产。第三阶段（高质量纹理增强）：将2D高斯导出为带颜色顶点的网格，然后在网格表面初始化3D高斯并绑定到网格上。这些"几何约束"高斯只能在网格表面附近移动，从而在SDS优化中保持结构一致性。直觉上，前两阶段像是"搭骨架"，第三阶段是在骨架上"贴皮肤"并精修细节。最终输出绑定到网格的3D高斯资产，还可无缝集成到动画、组合和物理模拟等下游管线中。
\textbf{锚点细化：} GaussianDreamerPro~\cite{2406_18462}观察到SDS优化中高斯基元倾向于无控制地增长，导致生成质量远低于重建质量。核心洞察是：将高斯基元绑定到合理的几何结构上，并在生成全过程中让几何结构持续演化。方法按阶段逐步丰富几何和外观，最终输出绑定到网格的3D高斯资产，在细节和质量上显著超越先前方法。生成的资产还可以无缝集成到动画、组合和物理模拟等下游管线中。
\textbf{方法信息：} Huawei 参与了这篇论文；作者关联代码：\url{https://github.com/hustvl/GaussianDreamerPro}；GitHub 关注度高，约 235 stars。



% METHOD: 2410.20723
\paragraph{CompGS（组合文本到 3D）（2024.10）}
CompGS~\cite{2410_20723} 解决多实体文本到 3D。它先利用 2D 组合生成得到构图合理图像，提取每个实体区域，用深度估计提升为各实体的粗 3DGS 初始化。优化时自动把高斯按空间分解为不同实体，在实体级单独用 SDS 保持各自特征，在组合级对整体场景用 SDS 保持交互关系，并动态调整实体位置和尺度。相对单物体 DreamGaussian，CompGS 的增量是把 2D 组合性迁移到 3DGS 的逐实体初始化和优化。
具体流程先让 LLM 从复杂 prompt 中拆出实体级文本，例如“猫”“椅子”“桌面关系”；文本到图像模型生成一张构图合理的组合图，再由分割模型得到每个实体 mask。每个 mask 对应的图像区域送入 image-to-3D 模型，得到粗 mesh 或粗几何，再初始化为该实体的一组高斯。优化阶段交替做两件事：实体级 SDS 只渲染和更新某个实体的高斯，使颜色、形状和小物体细节不被全局 prompt 淹没；组合级 SDS 渲染整个场景，约束多个实体之间的空间关系和语义交互。为了处理大小差异，CompGS 还把每个实体临时缩放到标准体积内优化，再映射回原场景坐标。这个设计的关键是利用 3DGS 可分组、可局部渲染的显式结构，把 2D 的组合性拆成实体初始化和分层优化。
\textbf{实现细节：} 组合式场景生成需要同时处理多个物体的空间关系和物理约束。CompGS（2024.10） 的核心洞察是：3DGS 的显式表示天然支持局部化渲染，可以单独渲染物体 A 和物体 B 所在的子空间，从而在"分而治之"的意义上实现全局组合生成。方法首先从复杂文本提示生成一张"良好组合"的 2D 图像作为参考，将 2D 组合性先验迁移到 3DGS 初始化中。然后设计动态 SDS 优化策略，自动分解高斯基元为不同实体部分，分别进行实体级优化（关注单个物体的几何细节）和组合级优化（关注物体间的交互与遮挡关系）。
\textbf{方法信息：} Stability AI 参与了这篇论文；项目页：\url{https://chongjiange.github.io/compgs.html}；完整公开代码暂未确认。
\textbf{设计动机与边界：} 相对原始 3DGS 以独立高斯和图像重建误差为主的设定，CompGS 把问题落在以下可实现的改动上：方法利用3DGS的显式性质实现组合生成：1)2D组合性初始化：先从复杂文本提示生成构图良好的2D图像，提取每个实体的2D区域，利用深度估计将2D区域提升为粗略3D形状作为各实体的3DGS初始参数。2)动态SDS优化：训练过程中自动将3DGS按空间位置分解为不同实体的子集。在实体级别对每个实体单独应用SDS优化保持各实体特征，在组合级别对整体场景应用SDS保持实体间合理交互关系。动态调整各实体的空间参数（位置/尺度），尤其增强较小实体的细节生成。3)可扩展到可控3D编辑和场景生成。其设计目的不是单独增加一个模块，而是用上述表示、监督或渲染约束消除原流程在目标任务中的具体失配。提出CompGS，首个基于3DGS的组合文本到3D生成框架。两个核心设计：1)2D组合性初始化：将成熟的2D组合性迁移到3DGS参数的逐实体初始化中，确保每个实体的3D先验一致且实体间交互合理；2)动态优化：自动将3DGS分解为不同实体部分，在实体和组合级别分别优化，动态调整空间参数增强小实体的细粒度细节。在T3Bench上生成质量优于现有方法。适用时还应保留论文给出的边界：生成质量受限于2D扩散模型的能力；超过3-4个实体的复杂场景仍有挑战；物体间物理交互建模有限；生成时间较长。
% METHOD: 2411.19654
\paragraph{TexGaussian（octree 高斯 PBR 材质生成）（2024.11）}
TexGaussian~\cite{2411_19654} 不生成新几何，而是给已有 mesh 生成 PBR 材质。GaussianDreamerPro 这类方法把高斯作为文本到 3D 资产的几何和外观表示；TexGaussian 的局部锚点是材质生成：mesh 几何已经给定，高斯只作为可微的三维材质中间层，最后还要烘焙回 UV 贴图。

流程从 mesh 表面采样开始。方法在表面采样大量点并构建 octree，最细叶节点自然贴近物体边界；每个叶节点中心放一个高斯，位置固定在 mesh 表面附近。与普通 3DGS 只预测颜色和透明度不同，TexGaussian 还为每个高斯预测 roughness 和 metallic，用于 PBR 渲染。octree-based 3D U-Net 在这个稀疏层次结构上一次前向回归高斯参数，训练时同时约束 2D 渲染结果和 3D 高斯参数。推理阶段先从多个视角渲染这些 PBR 高斯，再通过可微 mesh 渲染把结果烘焙到输入 mesh 的 UV 空间。它的价值在于把“贴图生成”改写成“表面对齐高斯参数回归”，比逐视角扩散更容易保持材质在 3D 表面上的一致性。
\textbf{实现细节：} TexGaussian（2024.11） 专注于为3D网格快速生成PBR材质。动机。MGM 生成完整 3DGS 资产（含 PBR），但实际工业管线中已有大量 3D 网格只需材质贴图；TexGaussian 假设输入网格已知，仅生成其表面的 PBR 材质，且通过八叉树对齐确保高斯自然贴合网格表面。与MGM的差异：MGM从文本生成完整的3DGS资产（含PBR），而TexGaussian假设已有输入网格，仅生成其表面材质贴图。八叉树对齐的高斯泼溅：在输入网格表面采样$N=100{,}000$个3D点构建八叉树，将3D高斯放置在每个最细叶节点中心（不额外调整位置），确保高斯自然贴合网格表面。每个高斯包含13个参数：1维不透明度、3维缩放、4维旋转、3维反照率、1维粗糙度、1维金属度（不含球谐系数，因PBR图像无光照变化）。训练时采用softplus激活并将缩放乘以0.01，确保初始高斯贴合物体轮廓而非向外扩张。基于八叉树的3D U-Net：使用八叉树卷积神经网络预测高斯参数，采用回归目标而非扩散去噪——输入为八叉树叶节点的几何特征（法线、局部位移）和CLIP文本特征，输出直接为13通道高斯参数。推理时将渲染的多视角图像烘焙到UV空间生成最终贴图。
\textbf{设计动机与边界：} 相对原始 3DGS 以独立高斯和图像重建误差为主的设定，TexGaussian 把问题落在以下可实现的改动上：TexGaussian方法流程：(1)在输入mesh表面采样100,000个3D点，构建八叉树直到最大深度，所有最细叶节点自然位于3D物体边界。(2)在每个最细叶节点中心放置一个3D高斯，位置固定在mesh表面上。每个高斯除颜色和透明度外，增加roughness和metallic两个额外通道用于PBR材质渲染。(3)使用octree-based 3D U-Net预测所有高斯的参数，训练时最小化2D渲染图和3D高斯参数的差异。(4)推理时渲染多视角图像，通过可微分mesh渲染烘焙到输入mesh的UV空间。整个流程是回归训练而非扩散去噪，支持单次前向生成。其设计目的不是单独增加一个模块，而是用上述表示、监督或渲染约束消除原流程在目标任务中的具体失配。首次将octant-aligned 3DGS用于PBR材质生成，通过八叉树结构确保高斯自然位于mesh表面，以回归训练而非扩散去噪实现单次前向生成，在视觉质量和速度上均优于先前方法。适用时还应保留论文给出的边界：生成质量受限于3D U-Net的容量；对极其复杂的几何结构处理有挑战；PBR材质的物理准确性仍有提升空间。
\textbf{方法信息：} Baidu 参与了这篇论文；代码：\url{https://github.com/ymxbj/texgaussian}。



% METHOD: 2412.04470
\paragraph{Turbo3D（少步多视图扩散与潜空间 GS-LRM）（2024.12）}
Turbo3D~\cite{2412_04470} 追求超快速文本到 3D。它先用多视图扩散教师和单视图美学教师做双教师蒸馏，避免少步扩散退化成过简风格；生成器只需少数步产生多视图 latent。随后 latent-space GS-LRM 直接从多视图 latent 重建 3DGS，跳过 VAE 解码并减少 Transformer 序列长度。相对 GaussianDreamer 需要每个对象优化，Turbo3D 的改变是把优化式文本到 3D 压成少步生成和一次重建。
第一阶段基于 DiT 训练文本到多视图扩散教师，然后用分布匹配蒸馏把生成步数压到很少。直接蒸馏容易让结果退化成简单 Objaverse 风格，Turbo3D 因此引入双教师：多视图教师负责跨视角一致，单视图美学教师负责每个视角的照片真实感，两个 KL 目标共同约束少步生成器。第二阶段的 GS-LRM 不再等 VAE 把多视图 latent 解码成像素图，而是直接读 latent 重建 3D Gaussians；输入序列更短，且少了一次像素空间转换。这个增量把“多视图生成”和“高斯重建”压到同一条低延迟链路。
\textbf{实现细节：}
\textbf{方法信息：} Adobe Research 参与了这篇论文；仓库：\url{https://github.com/hzhupku/Turbo3D}，当前约 76 stars，但核验时仓库根目录主要是 README 和 LICENSE，完整可运行代码暂未确认。
\textbf{设计动机与边界：} 相对原始 3DGS 以独立高斯和图像重建误差为主的设定，Turbo3D 把问题落在以下可实现的改动上：两阶段管线：(1) 快速4步4视图扩散生成器：基于DiT架构训练多步文本到多视图扩散模型作为教师。发现直接用分布匹配蒸馏(DMD)到少步会导致复合模式坍塌(生成过于简化的Objaverse风格)。解决方案：双教师蒸馏——多视图教师确保跨视图一致性，单视图教师(在高质量美学图像上训练)确保每视图的照片真实感。损失函数结合MV和SV教师的KL散度。(2) 潜空间GS-LRM：跳过VAE解码步骤，直接将多视图潜表示输入GS-LRM重建3D高斯。输入从像素空间变为潜空间后，transformer序列长度减半，极大提升效率。其设计目的不是单独增加一个模块，而是用上述表示、监督或渲染约束消除原流程在目标任务中的具体失配。提出亚秒级文本到3D生成系统。核心创新：(1) 双教师蒸馏策略：多视图教师教视图一致性，单视图教师教照片真实感，解决复合模式坍塌问题。(2) 潜空间GS-LRM：将高斯重建器输入从像素空间移到潜空间，消除VAE解码时间并将transformer序列长度减半。4步生成4视图潜表示后即时重建3DGS。适用时还应保留论文给出的边界：生成质量受限于4视图信息的完备性；对复杂场景文本描述的理解可能不够精确；潜空间表征的细节保真度可能低于像素空间。
% METHOD: 2501.10462
\paragraph{BloomScene（渐进场景生成与压缩）（2025.01）}
BloomScene~\cite{2501_10462} 可以看作 LucidDreamer 式渐进场景生成的工程化延伸。LucidDreamer 沿相机轨迹反复 inpainting、估深和合并点云，但单向外扩容易累积深度误差，生成出的 3DGS 也会很大。BloomScene 保留“先生成初始图像，再估计深度，再逐步扩展点云并优化为 3DGS”的主线，同时把场景扩展、深度约束和压缩放进同一个框架。

具体来说，PSG 模块先用文本到图像模型或输入图像得到起始视图，单目深度估计器给出初始几何，再采用双向旋转相机策略从起点向两侧扩展。这样未知区域不是只沿一个方向持续外推，而是在相反方向上互相校正，减少修复误差越走越大的问题。DPR 模块在不同尺度约束深度：全局层面对齐场景尺度，局部层面保持平滑，高频层面保留边界和细节。最后，SCC 模块把生成后的无序 anchor 属性放进结构化哈希网格中建模上下文，利用空间相关性压缩冗余属性。相对 LucidDreamer，BloomScene 的关键增量不是换掉 3DGS，而是把渐进生成的深度稳定性和高斯资产大小一起纳入设计。
\textbf{实现细节：} 除了上述方法，还有一些工作聚焦于特定环节的改进。基于词汇丰富度的GS初始化框架 指出简单球体初始化难以处理复杂文本描述的形状，提出将3D高斯聚合到空间均匀体素中，设计全局信息感知（GIP）块和高斯-文本融合（GTF）块，使每个高斯能吸收来自其他区域的空间信息和来自文本的语义信息。Hash3D 提出无需训练的通用3D生成加速方法——观察到相邻时间步和相机角度渲染的特征图存在冗余，通过自适应网格哈希复用特征图，将文本到3D优化加速1.3--4倍。PanoDreamer 利用大语言模型和warp-refine管线生成360度全景，然后提升为3D点云，再用多种方法生成与初始点云一致的额外图像来扩展和精炼，最终通过3DGS创建完整3D场景。BloomScene 提出轻量级结构化3DGS框架用于跨模态场景生成，包含渐进式点云重建和分层深度先验正则化机制，以及结构化哈希网格压缩冗余。Dream3DVG 提出双分支优化框架——辅助3DGS分支桥接文本提示与矢量图形的域差异，通过调度分类器自由引导实现从粗到细的渐进细节控制。GaussVideoDreamer 通过渐进式视频修复策略和不一致感知高斯泼溅（Inconsistency-Aware GS）实现单图3D场景生成。EarthCrafter 将3DGS生成扩展到地理尺度——在50000个600m$\times$600m航拍场景上训练双稀疏潜在扩散模型，支持从语义引导的城市布局生成到无条件地形合成的多种应用。
\textbf{设计动机与边界：} 相对原始 3DGS 以独立高斯和图像重建误差为主的设定，BloomScene 把问题落在以下可实现的改动上：BloomScene包含三个核心模块：(1)跨模态渐进场景生成（PSG）：使用预训练文本到图像模型生成初始图像，单目深度估计器获取深度图。采用双向旋转相机策略（从初始位置向两侧旋转）最小化图像修复累积误差。渐进式进行图像修复和点云更新，逐步构建完整场景。(2)层次化深度先验正则化（DPR）：利用多层级深度精度和平滑约束，包括全局深度对齐、局部深度平滑和高频细节保留三个层次，增强生成场景的真实感和连续性。(3)结构化上下文引导压缩（SCC）：使用结构化哈希网格建模无序锚点属性的上下文，消除结构冗余。通过学习属性间的空间相关性，显著减少高斯模型的存储开销。其设计目的不是单独增加一个模块，而是用上述表示、监督或渲染约束消除原流程在目标任务中的具体失配。提出BloomScene框架，包含跨模态渐进场景生成、层次化深度先验正则化和结构化上下文引导压缩机制三项创新。通过渐进式点云重建和3DGS实现高质量场景生成，利用哈希网格压缩消除结构冗余，显著减少存储开销。适用时还应保留论文给出的边界：依赖预训练文本到图像模型和深度估计器的质量；对非常复杂的户外场景可能产生不连贯区域；渐进生成的速度仍需提升。
\textbf{方法信息：} ByteDance 参与了这篇论文；代码：\url{https://github.com/sparklingh/bloomscene}。



% METHOD: 2505.15528
\paragraph{PlantDreamer（植物几何与纹理的扩散高斯）（2025.05）}
PlantDreamer~\cite{2505_15528} 是领域化文本到 3D。植物的难点不是只有外观复杂，而是枝条、叶片和花序会形成大量细长结构，通用 text-to-3D 的 SDS 优化容易把叶片粘成块，或在不同视角产生不一致纹理。PlantDreamer 仍以 3DGS 作为可优化表示，但把植物领域先验接到初始化、扩散引导和高斯清理三个位置。

具体流程可以从两类输入启动：一类是 L-System 程序生成的植物网格，另一类是真实植物点云。它们先提供大致枝干和叶片空间分布，随后用 depth ControlNet 约束扩散监督，使新视角纹理生成时不偏离已有几何深度。对于特定物种的叶片纹理和颜色分布，PlantDreamer 用 LoRA 微调扩散模型，使 SDS 梯度更符合植物类别而不是泛化物体风格。优化过程中，adaptive Gaussian pruning 会移除异常大的高斯或不稳定高斯，避免少数高斯把细枝叶结构糊成大片。相对 GaussianDreamerPro 的通用资产生成，PlantDreamer 的增量在于把植物几何先验和物种纹理先验显式插入 3DGS 优化链路。
\textbf{实现细节：} 除了通用场景生成，3DGS+文本还渗透到特定领域。SIC3D 引入风格图像条件控制，通过变分风格化分数蒸馏（VSSD）损失从参考图像迁移风格到3DGS。3DBonsai 设计了可训练的3D空间殖民算法生成盆景结构，作为3DGS的几何先验。PlantDreamer 利用深度ControlNet和可调高斯剔除算法生成高保真合成植物模型。SplatFont3D 专注于3D艺术字体生成，通过Glyph2Cloud模块和动态组件分配策略实现部件级风格控制。ClotheDreamer 和GarmentDreamer 分别面向可穿戴3D服装生成，前者提出解耦服装高斯（DCGS）分别优化服装和身体，后者利用3DGS引导和神经纹理场生成多样化的服装几何与纹理。ART3D 将扩散模型与3DGS结合用于艺术场景生成，通过图像语义迁移算法和深度一致性模块桥接艺术与真实图像的域差异。
\textbf{方法信息：} 代码：\url{https://github.com/lewis-stuart-11/plantdreamer}。
\textbf{设计动机与边界：} 相对原始 3DGS 以独立高斯和图像重建误差为主的设定，PlantDreamer 把问题落在以下可实现的改动上：PlantDreamer框架包含以下核心模块：（1）3DGS生成基线：基于SDS损失（Score Distillation Sampling）使用Stable Diffusion 2.1引导3DGS场景优化，通过迭代渲染-去噪-更新过程逐步提升纹理保真度和几何一致性。（2）深度ControlNet：引入深度条件控制确保训练过程中的几何一致性，使生成的3D植物模型结构更准确。（3）LoRA微调：在每种植物30张图像上训练LoRA模块，提供针对特定物种的风格迁移，改善纹理真实感。（4）自适应高斯剔除算法：检测并移除训练过程中出现的异常大高斯，防止平滑或扭曲的表面，改善收敛性。输入支持两种方式：L-System程序化生成的网格或真实世界植物点云。其设计目的不是单独增加一个模块，而是用上述表示、监督或渲染约束消除原流程在目标任务中的具体失配。提出专门针对复杂植物几何和纹理的3DGS生成框架PlantDreamer。集成深度ControlNet确保几何一致性，LoRA微调提升植物纹理真实感，自适应高斯剔除算法改善收敛性。支持纯合成植物生成（通过L-System网格）和真实植物点云增强两种模式，在植物高保真度生成上超越通用text-to-3D模型。适用时还应保留论文给出的边界：依赖输入点云质量；每种植物需要30张图像训练LoRA；程序化L-System生成的网格可能不够真实。
% METHOD: 2508.15169
\paragraph{MeSS（城市 mesh 引导户外场景生成）（2025.08）}
MeSS~\cite{2508_15169} 用城市 mesh 作为几何先验生成户外场景。Stage I 级联外扩 ControlNet，用 depth、semantic 和 normal 控制稀疏关键视图；Stage II 从 3DGS 渲染中间视图，并用 appearance-guided inpainting 修复遮挡；Stage III 用 global consistency alignment 消除曝光和风格不一致。生成过程中，高斯球同步初始化在 mesh 表面并构建 3DGS 场景。相对 ScenDi 的体素到高斯，MeSS 使用已有城市 mesh 约束跨视图几何。
\textbf{实现细节：} MeSS（2025.08） 利用城市网格模型作为几何先验，是城市场景生成路线的代表性方法。与全景图路线不同，MeSS 的输入是具有语义和实例标签但无纹理的 3D 城市网格模型。方法采用稀疏到密的视角生成策略：第一阶段通过级联 ControlNet（ControlNet-S 用于几何条件，ControlNet-n 用于外推生成）沿预定义相机路径自回归生成关键视角图像，深度以色图编码保留远近场深度差异；第二阶段从 3DGS 中间渲染缺失视角，通过外观引导修复（Appearance Guided Inpainting, AGInpaint）填补伪影；第三阶段通过全局一致性对齐消除不同视角学习的高斯泼溅面（Gaussian surfels）之间的视觉不一致。
\textbf{方法信息：} Huawei 参与了这篇论文。
\textbf{设计动机与边界：} 相对原始 3DGS 以独立高斯和图像重建误差为主的设定，MeSS 把问题落在以下可实现的改动上：方法分为三个阶段：Stage I使用级联外扩ControlNet生成几何一致的稀疏关键视图，ControlNet-S接收深度、语义和法线作为控制信号生成初始关键帧，后续关键帧通过前帧扭曲作为额外条件进行自回归生成；Stage II通过AGInpaint（外观引导修复）模块从3DGS高斯场渲染中间视图并修复遮挡伪影；Stage III使用GCAlign（全局一致性对齐）模块消除视觉不一致（如曝光差异）。在生成过程中同步将高斯球初始化在网格表面上构建3DGS场景，支持重光照和风格迁移。其设计目的不是单独增加一个模块，而是用上述表示、监督或渲染约束消除原流程在目标任务中的具体失配。提出MeSS框架，利用城市网格模型作为几何先验，通过三阶段管线（级联外扩ControlNet生成稀疏关键视图、AGInpaint传播密集中间视图、GCAlign全局一致性对齐）生成跨视角一致的户外场景，并同步重建3DGS场景。解决了视频扩散模型难以遵循预设相机路径、图像扩散模型缺乏跨视角一致性的问题，在几何对齐和生成质量上超越现有方法。适用时还应保留论文给出的边界：依赖城市网格模型的可用性和质量；生成质量受限于图像扩散模型的能力；对于极度复杂的城市拓扑结构可能仍有几何对齐问题。
% METHOD: 2601.15221
\paragraph{ScenDi（3D 到 2D 级联城市场景扩散）（2026.01）}
ScenDi~\cite{2601_15221} 面向可控城市场景生成。第一阶段训练 3D latent diffusion model：Voxel-to-3DGS VQ-VAE 把体素网格编码为 latent，再解码为 3D Gaussians，并可由 3D boxes、道路图或文本控制。第二阶段用 2D video diffusion 对低分辨率 3DGS 渲染进行外观细化，补高频和背景内容。相对纯 3D 方法，ScenDi 用 2D 扩散补外观；相对纯 2D 视频生成，它保留 3DGS 带来的相机可控性。
3D 阶段的关键是把城市先压成体素 latent，而不是直接让扩散模型处理海量高斯。Voxel-to-3DGS VQ-VAE 从 2D 监督和单目深度估计构造伪 3D 监督，学习“体素网格到高斯原语”的解码器；扩散模型在这个 latent 空间里接受 3D 边界框、道路图或文本控制，先决定道路、建筑和车辆等三维布局。2D 阶段再把低分辨率 3DGS 渲染交给视频扩散模型，填补高频纹理、远景背景和体素外内容。这个级联设计把几何可控性留给 3DGS，把视觉细节交给 2D 视频扩散，适合 Waymo、KITTI-360 这类城市场景。
\textbf{方法信息：} Alibaba 参与了这篇论文。
\textbf{设计动机与边界：} 相对原始 3DGS 以独立高斯和图像重建误差为主的设定，ScenDi 把问题落在以下可实现的改动上：方法分为两个阶段：(1)3D生成阶段：设计3D隐扩散模型(LDM)，包含Voxel-to-3DGS VQ-VAE和3D扩散模型。VQ-VAE从2D监督和单目深度估计获得的伪3D GT中学习，将体素网格编码为隐变量再解码为3D高斯原语。可通过3D边界框、道路图和文本指令可选控制生成。(2)2D细化阶段：训练2D视频扩散模型，以3DGS渲染的低分辨率图像为条件，填充高频细节并合成体素外背景内容。在Waymo和KITTI-360数据集上验证。其设计目的不是单独增加一个模块，而是用上述表示、监督或渲染约束消除原流程在目标任务中的具体失配。提出ScenDi级联式3D-to-2D场景扩散框架，结合3D和2D扩散模型生成城市场景。首先训练3D隐扩散模型生成粗3DGS场景（可通过边界框、道路图或文本控制），然后训练2D视频扩散模型增强外观细节。解决了纯3D方法外观细节退化和纯2D方法相机可控性不足的问题。适用时还应保留论文给出的边界：3D生成分辨率有限导致远处区域细节不足；依赖深度估计的准确性；2D细化可能引入与3D不一致的细节。
% METHOD: 2604.14268
\paragraph{HY-World 2.0（多模态世界生成与重建）（2026.04）}
HY-World 2.0~\cite{2604_14268} 不是单个 text-to-3D 模型，而是把“生成世界”和“重建世界”放到同一套 3DGS 系统里。LucidDreamer 从一个初始视角沿轨迹扩展场景，重点是用 inpainting 和深度对齐逐步补空间；HY-World 2.0 的尺度更大，输入可以是文本、单图、多视图图像或视频，输出是可导航、可交互的 3DGS 世界。它更像一条完整生产线：稀疏输入走生成分支，密集输入走重建分支，二者最后都落到 WorldMirror 和 3DGS 表示。

当输入是文本或单图时，系统先用 HY-Pano 2.0 生成或扩展全景图，得到一个覆盖 360 度视野的世界初始画布。随后 WorldNav 做轨迹规划：它不只是给相机随机绕一圈，而是结合场景语义、可见区域、遮挡和碰撞约束，选择能补全空间的路径。这个轨迹决定后续哪些位置需要生成关键帧，哪些区域需要从新视角看见，避免只在初始全景附近形成薄壳。

WorldStereo 2.0 负责世界扩展。它在关键帧 latent 空间运行相机引导的视频/多视角生成，并加入一致性记忆，让不同轨迹段生成的图像共享几何和外观线索。生成得到的关键帧、多视图图像或视频再交给 WorldMirror 2.0，后者是前馈式 3D 预测/重建模型，用这些观测恢复深度、法线、点云、相机和最终高斯表示。若输入本身就是多视图或视频，系统可以跳过前面的想象扩展，直接进入 WorldMirror 做几何一致重建。

最后的 WorldLens 是运行和交互层。它提供 3DGS 渲染、自动 IBL 光照、碰撞检测和角色探索接口，使生成结果不只是一段视频，而是能被浏览和交互的世界资产。相对 DreamGaussian 这种逐对象优化锚点，HY-World 2.0 的参照意义在于：3DGS 不再只是生成结果的表示，而是连接全景生成、轨迹规划、视频扩展、前馈重建和实时交互的统一世界格式。
\textbf{实现细节：} HY-World 2.0（2026.04） 是首个统一"生成"与"重建"的开源多模态世界模型框架，支持文本、单图、多视角图像和视频四种输入模态。文本/单图输入触发世界生成模式，通过四阶段管线生成可导航的 3DGS 场景：（1）HY-Pano 2.0 在大规模真实+合成混合全景数据上训练，通过无需显式几何变换的隐式学习实现从透视到全景的空间映射；（2）WorldNav 进行语义感知的轨迹规划，结合 42 视角密集深度采样和 NavMesh 碰撞检测；（3）WorldStereo 2.0 在关键帧潜在空间中运行相机引导视频扩散，通过全局几何记忆（GGM）和空间立体记忆（SSM）确保多轨迹一致性；（4）WorldMirror 2.0 通过前馈 3D 重建将关键帧融合为最终 3DGS。多视角/视频输入触发世界重建模式，WorldMirror 2.0 直接从密集观测恢复精确 3D 结构。
\textbf{锚点细化：} HY-World 2.0：是首个统一"生成"与"重建"的开源多模态世界模型框架，支持文本、单图、多视角图像和视频四种输入模态。文本/单图输入触发世界生成模式，通过四阶段管线生成可导航的 3DGS 场景：（1）HY-Pano 2.0 在大规模真实+合成混合全景数据上训练，通过无需显式几何变换的隐式学习实现从透视到全景的空间映射；（2）WorldNav 进行语义感知的轨迹规划，结合 42 视角密集深度采样和 NavMesh 碰撞检测；（3）WorldStereo 2.0 在关键帧潜在空间中运行相机引导视频扩散，通过全局几何记忆（GGM）和空间立体记忆（SSM）确保多轨迹一致性；（4）WorldMirror 2.0 通过前馈 3D 重建将关键帧融合为最终 3DGS。多视角/视频输入触发世界重建模式，WorldMirror 2.0 直接从密集观测恢复精确 3D 结构。
\textbf{设计动机与边界：} 相对原始 3DGS 以独立高斯和图像重建误差为主的设定，HY-World 2.0 把问题落在以下可实现的改动上：HY-World 2.0的四大模块：（1）HY-Pano 2.0：从文本或单视角图像生成高质量全景图；（2）WorldNav：3D场景理解和轨迹规划，实现全景到3D导航的桥接；（3）WorldStereo 2.0：基于关键帧的视图生成模型，引入一致性记忆机制，从多视角立体匹配扩展场景；（4）WorldMirror 2.0：前馈式通用3D预测模型，从多视角图像或视频重建3D世界，改进模型架构和学习策略。WorldLens渲染平台：引擎无关架构、自动IBL光照、高效碰撞检测、训练-渲染协同设计，支持角色交互探索。其设计目的不是单独增加一个模块，而是用上述表示、监督或渲染约束消除原流程在目标任务中的具体失配。提出HY-World 2.0，首个开源多模态世界模型框架，统一3D世界的生成和重建。支持文本、单视角图像、多视角图像和视频输入，通过四阶段方法（HY-Pano 2.0全景生成、WorldNav轨迹规划、WorldStereo 2.0视图生成、WorldMirror 2.0世界合成）生成3DGS场景。引入WorldLens高性能3DGS渲染平台，支持IBL光照和碰撞检测。适用时还应保留论文给出的边界：与闭源模型仍有性能差距；四阶段管道复杂度高；全景生成质量限制最终3D世界质量。
\textbf{方法信息：} Tencent 参与了这篇论文；代码：\url{https://github.com/filliptm/HY-World-2.0-Playground}；引用很多，Semantic Scholar 约 12 次。该 Playground 是围绕 Tencent HY-World 2.0/WorldMirror 的工程化运行界面，并非论文中全部上游组件的完整官方源码。



\section{单图到 3D 重建}

单图到 3D 介于重建和生成之间：可见面来自输入图，不可见面来自扩散或视频先验。高斯方法的具体差异在于，是否逐像素生成一个或多个高斯、是否用视频扩散探索新视角、是否再转 mesh 或 4D 动画。


% METHOD: 2405.20310
\paragraph{Hierarchical Splatter Image（一个像素多个高斯）（2024.05）}
A Pixel Is Worth More Than One 3D Gaussians~\cite{2405_20310} 扩展 Splatter Image：单图中一个像素不应只对应一个高斯，因为不可见面和遮挡区域需要额外表达。方法为每个像素预测父高斯和多个子高斯。父高斯由 U-Net 直接预测，保证可见面重建；子高斯由轻量 MLP 根据父高斯投影特征和目标相机嵌入生成，学习分配到不可见区域。训练先父后子。相对 GPS-Gaussian 或 pixel-aligned 3DGS，该方法的核心是层次化像素表达，专门补遮挡内容。
\textbf{实现细节：} 与 LeanGaussian 的无对应关系设计形成鲜明对比，针对 Splatter Image 像素→单一高斯一对一映射无法重建被遮挡背面区域的局限，Hierarchical Splatter Image（2024.05） 保留了像素对应关系，但将其从一对一推广为一对多的层级结构：每个像素对应一个父高斯和若干子高斯（child Gaussian）。子高斯通过轻量级条件 MLP 生成，其输入为父高斯的图像特征 $G^{\mathrm{Parent}}(I, P)$ 和目标相机视角的嵌入。层级结构的设计动机在于：原始的 Splatter Image 虽然能精确预测可见区域的高斯参数，但对输入视角中被遮挡的背面细节重建效果较差。子高斯专门负责补充这些不可见区域，通过显式编码父高斯与目标相机位姿之间的空间关系来推断遮挡部分。该模型采用两阶段训练：前 10K 次迭代仅训练 U-Net 以确保父高斯的合理预测，之后加入条件 MLP 进行端到端联合训练。
\textbf{方法信息：} Alibaba 参与了这篇论文。
\textbf{设计动机与边界：} 相对原始 3DGS 以独立高斯和图像重建误差为主的设定，A Pixel Is Worth More Than One 3D Gaussians in Single-View 3D Reconstruction 把问题落在以下可实现的改动上：方法基于Splatter Image框架扩展，采用类别特定的像素到3D-Gaussian映射管线。核心创新是每个像素不再映射到单个3D Gaussian，而是映射到一个父Gaussian和多个子Gaussian。父Gaussian通过U-Net图像编码器直接预测，与原始Splatter Image相同。子Gaussian通过轻量MLP生成，MLP以父Gaussian的投影图像特征和目标相机视角的嵌入作为输入条件。训练采用分阶段策略：先训练父Gaussian确保基本重建，再联合优化父子Gaussian。输入图像特征和目标相机位置的联合条件使子Gaussian学会分配到不可见区域，恢复父Gaussian遗漏的被遮挡细节。推理时只需单张输入图像和目标位姿即可生成新视角。其设计目的不是单独增加一个模块，而是用上述表示、监督或渲染约束消除原流程在目标任务中的具体失配。提出层次化Splatter Image方法，每个像素用多个3D Gaussian表示（父Gaussian和子Gaussian），父Gaussian按原始方式学习，子Gaussian通过轻量MLP基于父Gaussian的投影特征和目标相机嵌入生成。这种层次结构使模型能恢复被遮挡区域的细节，在ShapeNet-SRN和CO3D数据集上取得SOTA性能，尤其在被遮挡内容重建方面表现突出。适用时还应保留论文给出的边界：依赖类别级训练数据，泛化到未见类别能力有限；子Gaussian数量需手动设置；对极端遮挡（如物体背面）的恢复仍有限；需要已知相机内参。
% METHOD: 2407.04237
\paragraph{GSD（视角引导高斯扩散重建）（2024.07）}
GSD~\cite{2407_04237} 先将密集视角数据转换成固定数量 GS 数据集，在 16 维高斯特征上训练无条件 DDPM。推理时虽然模型无条件，但通过 splatting function 把输入图像细粒度 2D 特征传播到 3D Gaussian 空间，引导去噪采样重建目标对象；最后用辅助 2D diffusion 对渲染图做纹理增强。相对 DiffGS 的函数扩散，GSD 更依赖 2D 输入特征投影来引导高斯 denoising。
训练阶段先把密集视角对象重建成固定数量的高斯集合，并把每个高斯写成 16 维向量：位置、尺度、6D 旋转、不透明度和颜色。DDPM 在这些高斯向量上学习无条件先验，骨干可以用普通 Transformer。推理阶段的条件信息不是简单拼接到 token 上，而是把输入图像的细粒度 2D 特征通过 splatting function 投到候选高斯上，让去噪过程知道哪些三维位置和输入视图一致。最后的 2D 扩散增强只负责修纹理，不改变前面的高斯几何先验。
\textbf{方法信息：} Huawei 参与了这篇论文。
\textbf{设计动机与边界：} 相对原始 3DGS 以独立高斯和图像重建误差为主的设定，GSD 把问题落在以下可实现的改动上：GSD包含三个核心组件：(1) GS生成先验建模：将密集视角图像数据集转换为GS数据集（通过约束密集化保持每场景固定数量的GS椭球体），在此基础上训练无条件DDPM扩散模型。每个GS椭球体用16维特征表示（位置3维、缩放3维、旋转6维、不透明度1维、颜色3维），使用vanilla Transformer作为骨干网络；(2) 视角引导去噪采样：虽然模型是无条件训练的，但通过splatting函数将2D输入视图的细粒度特征传播到3D GS空间，引导去噪过程进行视角条件重建。具体通过将输入图像特征投射到每个GS椭球体实现特征聚合；(3) 2D扩散增强：使用辅助2D扩散模型对渲染图像进行精炼和复用，提升纹理和细节。其设计目的不是单独增加一个模块，而是用上述表示、监督或渲染约束消除原流程在目标任务中的具体失配。提出GSD框架，将Gaussian Splatting作为3D表示与扩散模型结合，通过splatting函数传播细粒度2D特征引导去噪采样，无需微调即可实现视角引导的3D重建，再用2D扩散模型提升渲染质量。在CO3D上达到SOTA。适用时还应保留论文给出的边界：每场景固定GS数量可能限制复杂物体表示；扩散模型推理较慢；对遮挡区域细节有限
% METHOD: 2409.12193
\paragraph{Vista3D（可见与不可见面双隐函数）（2024.09）}
Vista3D~\cite{2409_12193} 从单图生成物体，并把可见面与不可见面的建模解耦。粗阶段用 3DGS 和 Zero-1-to-3 XL 快速生成初始几何，并用 Top-K densification、尺度和透射率正则加速收敛。精阶段从高斯提取 SDF，用 FlexiCubes 优化几何和纹理，同时用两个隐函数分别处理可见面和不可见面。角度扩散先验组合协调 2D 与 3D-aware diffusion 梯度，平衡一致性和多样性。相对 ContrastiveGaussian，它更强调后处理网格和双隐函数几何。
粗阶段的重点是尽快得到可用几何：Zero-1-to-3 XL 提供新视角先验，3DGS 负责快速可微渲染；Top-K densification 只扩展梯度最大的高斯，尺度和透射率正则抑制过大的半透明团块。精阶段不继续在高斯上硬拉，而是从高斯提取 SDF，并用 FlexiCubes 做可微网格优化。两个隐函数分别拟合可见面和不可见面，使输入图能确定的区域与扩散补全的背面区域不互相污染。角度扩散先验组合则在 2D 扩散和 3D-aware 扩散之间调节梯度，避免只追求单视角纹理或只追求过度保守的几何。
\textbf{实现细节：} Vista3D（2024.09） 采用了粗-精两阶段策略，并提出了"3D暗面探索"（3D darkside exploration）的视角：将不可见区域的生成视为从单张图像发掘隐藏维度的问题，而非简单的重建任务。粗阶段利用3DGS快速（约30秒）生成初始几何，采用 Top-K 致密化和透射率/尺度正则化策略避免过度膨胀；精阶段直接从3DGS中提取 SDF，使用 FlexiCubes 可微分等值面表示进一步优化几何和纹理。Vista3D 的另一个核心贡献是扩散先验组合（prior composition）：通过梯度幅度约束同时使用两个独立的扩散先验——Zero-1-to-3 负责可见面的保真度，Stable Diffusion 负责不可见面的多样性——实现了生成一致性与多样性的平衡。
\textbf{方法信息：} Huawei 参与了这篇论文；仓库：\url{https://github.com/florinshen/vista3d}，当前约 57 stars，但核验时只有 README，完整代码暂未确认。
\textbf{设计动机与边界：} 相对原始 3DGS 以独立高斯和图像重建误差为主的设定，Vista3D 把问题落在以下可实现的改动上：Vista3D采用两阶段策略：(1)粗阶段：利用3DGS从单图像快速生成初始几何（约1分钟），使用Zero-1-to-3 XL作为2D扩散先验，采用基于梯度的Top-K密集化策略（聚焦最高梯度的高斯点），引入高斯尺度和透射率正则化加速收敛；(2)精阶段：从高斯泼溅直接提取SDF，使用FlexiCubes可微等面表示进一步优化几何和纹理，设计两个独立隐函数分别捕捉可见面和不可见面；(3)角度扩散先验组合：协调2D扩散先验和3D感知扩散先验的梯度，通过角度约束实现生成的一致性与多样性平衡。其设计目的不是单独增加一个模块，而是用上述表示、监督或渲染约束消除原流程在目标任务中的具体失配。Vista3D在5分钟内从单图像生成高质量3D物体，创新性地使用双隐函数解耦可见面和不可见面，并通过角度扩散先验组合协调两种扩散先验。适用时还应保留论文给出的边界：5分钟生成时间仍有一定开销；对极度复杂拓扑结构的恢复有限。
% METHOD: 2410.06756
\paragraph{DreamMesh4D（mesh-Gaussian 混合 4D）（2024.10）}
DreamMesh4D~\cite{2410_06756} 从单目视频生成 4D 内容。静态阶段先用 image-to-3D 获得粗 mesh，再用 SuGaR 把高斯绑定到 mesh 三角面上，通过光度和 SDS 优化静态表面。动态阶段在 mesh 表面采样稀疏控制点构建 deformation graph，变形网络预测控制点变换，再用 LBS 和 dual quaternion skinning 混合驱动 mesh 顶点和表面高斯。相对纯 4DGS，它的增量是 mesh 兼容图形管线，且控制点让动态优化更稳定。
静态阶段先把输入视频或参考图提升成粗 mesh，再用 SuGaR 式表面高斯把渲染表示贴到三角面附近；光度损失锁住参考视图，SDS 损失补未见区域。动态阶段不直接给每个高斯预测独立运动，而是在 mesh 表面采样稀疏控制点形成 deformation graph。网络预测控制点的刚性或近刚性变换，随后用 LBS 的简单插值和 dual quaternion skinning 的低伪影旋转混合共同驱动 mesh 顶点和表面高斯。这样，运动场被控制点约束在可解释的表面结构上，减少纯 4DGS 中每个高斯各自动导致的抖动。
\textbf{方法信息：} 引用很多，Semantic Scholar 约 66 次。
\textbf{设计动机与边界：} 相对原始 3DGS 以独立高斯和图像重建误差为主的设定，DreamMesh4D 把问题落在以下可实现的改动上：DreamMesh4D包含两个阶段：(1)静态阶段：通过image-to-3D生成获取粗mesh，使用SuGaR将高斯溅射绑定到mesh三角面上，通过参考视图光度损失和SDS损失优化静态表面高斯和mesh顶点；(2)动态阶段：在mesh表面均匀采样稀疏控制点构建变形图，使用变形网络预测控制点变换，通过混合几何蒙皮算法（结合LBS的简单性和DQS的无伪影特性）驱动mesh顶点和表面高斯的变形。稀疏控制点机制大幅提升计算效率，变形图提供额外约束。两阶段训练分别优化静态和动态参数。其设计目的不是单独增加一个模块，而是用上述表示、监督或渲染约束消除原流程在目标任务中的具体失配。提出DreamMesh4D，首个结合mesh表示和高斯溅射的4D生成框架，通过稀疏控制点和混合几何蒙皮（LBS+DQS）算法实现高效时空一致的4D内容生成，与现代图形管线兼容适用时还应保留论文给出的边界：依赖初始mesh的质量，对大幅非刚性变形的处理可能受限；稀疏控制点数量需要手动设置
% METHOD: 2410.09467
\paragraph{Hybrid Diffusion Priors（频率分解的单图优化）（2024.10）}
Enhancing Single Image to 3D Generation~\cite{2410_09467} 用 hybrid frequency score distillation 结合 2D 和 3D diffusion priors。流程先用 Zero123++ 生成多视图，再由 LGM 初始化粗 3DGS；优化时把图像转到频率域，低频使用 3D diffusion prior 保持几何一致，高频使用 2D diffusion prior 精炼纹理。相对 GaussianDreamer 同时依赖 2D 和 3D 扩散，但没有明确频段分工；该方法把“几何”和“纹理”用频率掩码拆开。
\textbf{实现细节：} 方法采用两阶段流程：（1）使用Zero123++从单图生成多视角图像，通过LGM大重建模型初始化粗3DGS表示；（2）提出混合频率分数蒸馏损失(hf-SDL)优化粗3DGS。hf-SDL的核心是将图像转换到频率域，使用离散傅里叶变换(DFT)。在低频部分使用3D扩散模型（Zero123）的几何先验保持跨视角一致性，在高频部分使用2D扩散模型（Stable Diffusion）精炼纹理保真度。通过频率掩码实现两个模型在不同频段的互补。 提出混合频率分数蒸馏损失(hf-SDL)，桥接2D和3D扩散模型的输出差距，在低频使用3D先验保持几何一致性，在高频使用2D先验精炼纹理，达到SOTA的单图3D生成效果。 仍需数分钟优化时间；依赖预训练扩散模型质量；对极端形状物体可能受限
\textbf{方法信息：} Amazon 参与了这篇论文。
\textbf{设计动机与边界：} 相对原始 3DGS 以独立高斯和图像重建误差为主的设定，Enhancing Single Image to 3D Generation using Gaussian Splatting and Hybrid Diffusion Prio 把问题落在以下可实现的改动上：方法采用两阶段流程：（1）使用Zero123++从单图生成多视角图像，通过LGM大重建模型初始化粗3DGS表示；（2）提出混合频率分数蒸馏损失(hf-SDL)优化粗3DGS。hf-SDL的核心是将图像转换到频率域，使用离散傅里叶变换(DFT)。在低频部分使用3D扩散模型（Zero123）的几何先验保持跨视角一致性，在高频部分使用2D扩散模型（Stable Diffusion）精炼纹理保真度。通过频率掩码实现两个模型在不同频段的互补。其设计目的不是单独增加一个模块，而是用上述表示、监督或渲染约束消除原流程在目标任务中的具体失配。提出混合频率分数蒸馏损失(hf-SDL)，桥接2D和3D扩散模型的输出差距，在低频使用3D先验保持几何一致性，在高频使用2D先验精炼纹理，达到SOTA的单图3D生成效果。适用时还应保留论文给出的边界：仍需数分钟优化时间；依赖预训练扩散模型质量；对极端形状物体可能受限
% METHOD: 2412.09597
\paragraph{LiftImage3D（视频生成先验释放到 3DGS）（2024.12）}
LiftImage3D~\cite{2412_09597} 从单图生成大视角场景。视频扩散模型在大相机运动下容易退化，方法用 joint trajectory strategy 把大运动拆成小步，并用 MotionCtrl 控制相机轨迹。生成帧再由 MASt3R 估计位姿和粗点云。为处理每帧生成的独立畸变，3DGS 表示分成规范高斯和帧特定 distortion 参数，优化时把畸变影响分离出去。相对 NavCrafter 的交互探索，LiftImage3D 更强调从单图释放视频扩散先验并恢复规范 3DGS。
流程上，它先规划一条由小位移拼接成的大视角轨迹，避免视频扩散一次处理过大相机运动。MotionCtrl 负责把轨迹条件注入 SVD，生成覆盖更多视角的图像序列；MASt3R 再从这些帧估计相机和粗点云，为 3DGS 初始化提供几何支撑。优化高斯时，方法显式区分规范高斯和每帧畸变参数：规范高斯承载最终三维资产，畸变参数吸收视频扩散在不同帧中引入的局部形变。这样，生成视频的不稳定性不会直接固化成错误几何。
\textbf{实现细节：} 前文介绍的前馈重建路线（如 LeanGaussian、DiffusionGS）通过单次网络推理直接输出3DGS参数，速度快但受限于训练数据的几何先验——当输入视角稀缺时，不可见区域的几何只能靠网络"猜"，质量天花板明显。优化式路线走了另一条路：不直接预测3D，而是先用扩散模型生成多视角图像作为"伪监督"，再通过迭代优化将这些2D观测提升为3D。LiftImage3D（2024.12） 正是这条路线的代表性工作，它将视频扩散先验与3DGS重建统一为端到端管线，系统性地解决了视频扩散先验用于单图3D重建时的三大挑战：（1）大相机运动下视频生成质量退化、（2）生成帧缺乏精确的相机位姿、（3）扩散去噪过程引入的几何扭曲破坏3D一致性。核心改动：传统的"视频到3D"管线直接用运动可控的视频扩散模型（LVDM，Latent Video Diffusion Model，潜在视频扩散模型）生成完整视频序列，再用 COLMAP 估计位姿。但大相机运动会导致 LVDM 生成模糊帧，而 COLMAP 对扩散模型产生的几何不一致极度脆弱，往往直接失败。
\textbf{方法信息：} Huawei 参与了这篇论文。
\textbf{设计动机与边界：} 相对原始 3DGS 以独立高斯和图像重建误差为主的设定，LiftImage3D 把问题落在以下可实现的改动上：LiftImage3D包含三个核心组件：(1) 关节轨迹策略：将大相机运动分解为可控的小步骤，利用MotionCtrl的相机运动控制模块引导SVD生成，保证生成质量的同时扩展视角覆盖范围；(2) 基于学习的视觉里程计：使用MASt3R从生成帧直接估计相机位姿和粗几何（点云），通过Procrustes对齐获得相对位姿；(3) 程变感知3DGS表示：为每个帧学习独立的畸变参数，同时优化输出无畸变的规范高斯。具体地，每个高斯的属性分解为规范部分和帧特定畸变部分，通过优化分离畸变影响。最终从规范高斯实现高质量新视角合成。其设计目的不是单独增加一个模块，而是用上述表示、监督或渲染约束消除原流程在目标任务中的具体失配。提出LiftImage3D，通过有效释放LVDM生成先验同时确保3D一致性来解决单图像3D重建挑战。核心创新包括：关节轨迹策略分解大相机运动、使用MASt3R直接估计相机轨迹、畸变感知3DGS表示处理帧间独立畸变。在LLFF、DL3DV和Tanks and Temples上达到SOTA。适用时还应保留论文给出的边界：依赖视频扩散模型的生成质量；对极度复杂场景可能引入生成伪影；多步骤生成流水线增加总处理时间。
% METHOD: 2504.08100
\paragraph{ContrastiveGaussian（对比学习精炼 SDS 单图生成）（2025.04）}
ContrastiveGaussian~\cite{2504_08100} 从单张图像生成 3DGS，解决的是单图扩散先验带来的视角不一致。普通流程通常先用 Zero-1-to-3 XL 之类的新视角扩散生成若干监督图，再用 SDS 或重建损失优化高斯；问题是扩散结果本身质量参差不齐，错误视角如果被同等当作监督，会把 3DGS 拉向伪影。

ContrastiveGaussian 的第一阶段仍保留这条基线：先用超分模型增强输入图，再用 Zero-1-to-3 XL 生成新视角，优化高斯时用 reference loss 锁住输入视角。第二阶段的改变是把扩散输出按质量变成正负样本。高质量扩散图或超分图作为正样本，不一致、模糊或有伪影的新视角作为负样本；LPIPS 用来度量渲染图与这些样本的感知距离，Quantity-Aware Triplet Loss 再根据正负样本数量和质量调权。这样 3DGS 不只是被扩散图牵引，还会主动远离扩散模型产生的坏视角。相对普通 SDS，它把“扩散先验不稳定”从噪声来源改造成训练信号。
\textbf{实现细节：} ContrastiveGaussian（2025.04） 将对比学习引入3DGS生成过程以解决 SDS 优化导致的纹理模糊和几何不一致问题。其核心思路是：利用2D扩散模型为新视角生成参考图像，通过超分辨率增强输入质量后，将渲染图像与参考图像构成正样本对、与无关图像构成负样本对，用感知损失（perceptual loss）引导高斯参数优化。该方法还提出了 Quantity-Aware Triplet Loss（QA-Triplet Loss），根据正负样本在特征空间中的距离动态调整 triplet 损失的边际，解决了训练中正负样本数量不均衡的问题。最后通过扩散模型的多步 MSE 损失细化粗纹理。
\textbf{方法信息：} 代码：\url{https://github.com/yanllan-ljb/contrastivegaussian}。



% METHOD: 2512.06158
\paragraph{Track4DGen（基础跟踪器运动先验）（2025.12）}
Tracking-Guided 4D Generation~\cite{2512_06158} 从静态 3D 模型生成时间一致动画。DreamMesh4D 用 mesh 控制点约束运动，Track4DGen 则把基础跟踪器学到的密集点对应当作运动先验。第一阶段基于 Animate3D 的多视角视频扩散骨干，并结合 MVDream 的多视角注意力和 AnimateDiff 的时间注意力；训练时使用 CoTracker 这类跟踪模型提取的点轨迹，在特征层强制生成视频保持跨帧对应关系。这样扩散特征本身携带“同一个点在时间中怎么移动”的信息，而不只是生成看起来连贯的 RGB 帧。

第二阶段做混合 4D-GS 重建。系统以静态 3D 高斯作为规范资产，从多视角视频估计时间相关运动场；运动编码时把第一阶段的跟踪感知扩散特征与 Hex-plane 特征拼接，再预测高斯位移和增强的 4D 球谐外观。相对只用 Hex-plane 学运动场的 4DGS，它的增量是把跟踪器先验送进扩散和重建两端，减少外观漂移并增强跨视角一致性。
\textbf{实现细节：} Animate3D 等方法用多视角视频扩散模型（MV-VDM）生成动画视频再重建 4D 资产，但仅依赖像素/潜空间中的视频扩散损失，缺少特征级的时间感知监督，导致长时间生成中出现外观漂移（物体逐帧退化或改变）。Track4DGen（2025.12） 的核心思想是利用基础点跟踪器 CoTracker3 的密集帧间对应关系作为"运动专家"，将其知识注入扩散特征空间，从而在生成阶段就建立强时间一致性。第一阶段升级 MV-VDM。Track4DGen 在标准 U-Net 去噪架构上，通过系统性的逐块探测发现：U-Net 解码器中第二个时空上采样块的时间运动模块的特征具有最强的长时间对应能力。因此在每个视角的第一帧图像上叠加 $15 \times 15$ 网格、排除物体外区域后，随机选取 8 个查询点/视角，用 CoTracker3 跟踪其跨帧轨迹，然后将跟踪点映射到该特征空间的对应位置并提取描述子。对应跟踪损失强制特征级点对应一致：$\mathcal{L}_{corr} = \frac{1}{|Q|}\sum_{q \in Q} \sum_{t} [1 - \cos(f_t(p_q^t), f_1(p_q^{}1))]$，其中 $f_t$ 为第 $t$ 帧的扩散特征、$p_q^t$ 为查询点 $q$ 在帧 $t$ 的 2D 位置。多视角扩散训练用 AdamW（学习率 $10^{-6}$）在 16 块 H100 上训练 20 epoch，批大小 1024（$16 \times 4$ 视角 $\times 16$ 帧），推理时 25 步采样。第二阶段混合 4D-GS 重建。经典 4D-GS 仅用 Hex-plane 特征表示运动场，Track4DGen 将第一阶段训练的扩散特征（已编码跟踪先验）与 Hex-plane 特征在共置点上拼接，形成混合运动编码。Hex-plane 分辨率 $[100, 100, 8]$、特征维度 16，先用光度/重建损失进行 750 次运动重建（批大小 64），再启用 4D-SDS 和 ARAP 正则额外 250 次迭代。
\textbf{方法信息：} Microsoft 参与了这篇论文。
\textbf{设计动机与边界：} 相对原始 3DGS 以独立高斯和图像重建误差为主的设定，Tracking-Guided 4D Generation 把问题落在以下可实现的改动上：Track4DGen采用两阶段管道：（1）多视角视频扩散与密集点跟踪：基于Animate3D的MV-VDM骨干，结合MVDream的多视角3D注意力和AnimateDiff的时间注意力。关键创新是在扩散生成器内强制执行密集的特征级点对应关系监督，利用基础跟踪模型（如CoTracker）提取的运动先验，产生时间一致的特征抑制外观漂移并增强跨视角一致性。扩散特征编码了跟踪感知的对应关系信息。（2）混合4D-GS重建：以静态3D高斯为规范资产，从多视角视频估计时间依赖的运动场。创新点在于混合运动编码——将同位置的扩散特征（携带Stage-One跟踪先验）与Hex-plane特征拼接，并增强4D球谐函数实现更高保真度的动态建模。相比纯Hex-plane方法，混合编码利用了丰富的语义运动信息。同时发布Sketchfab28高质量数据集用于4D生成基准测试。其设计目的不是单独增加一个模块，而是用上述表示、监督或渲染约束消除原流程在目标任务中的具体失配。提出将基础点跟踪器的运动先验显式注入多视角视频扩散和4D-GS重建两个阶段，在多视角视频生成和4D生成基准上均超越基线，并发布Sketchfab28高质量4D生成基准数据集。适用时还应保留论文给出的边界：两阶段管道计算开销较大；依赖多视角视频扩散模型质量；对快速复杂运动可能仍有时序漂移。
% METHOD: 2603.17131
\paragraph{SMAL-pets（单图可动画宠物高斯）（2026.03）}
SMAL-pets~\cite{2603_17131} 把单图到 3D avatar 从人扩展到狗等宠物。人体头像和人体 avatar 常借助 SMPL/SMPL-X 提供骨架、形状和蒙皮先验；该方法把对应先验换成 SMAL，让四足动物的体型、关节和姿态变化有可控参数。输入是一张宠物图像，系统先生成合成多视角监督，再进入两阶段重建：第一阶段把高斯绑定到 SMAL 网格上，利用形状系数、姿态系数、全局平移和可学习顶点偏移确定粗几何；第二阶段解除部分束缚，用 3DGS 补体积细节和毛发外观。

动画阶段，高斯通过 SMAL 蒙皮跟随底层网格变形，因此新姿态不是重新优化一套 3DGS，而是由参数化动物模型驱动已有高斯。文本提示可用于外观精化或动作编辑。相对人体 SMPL 路线，它的核心增量是把“参数化结构先验 + 高斯体积细节 + 蒙皮动画”迁移到宠物对象上。
\textbf{实现细节：} SMAL-pets（2026.03） 将3DGS与 SMAL（Skinned Multi-Animal Linear）参数化模型结合，从单张图像生成高质量、可动画化的动物（狗）化身。该方法采用混合表示：SMAL 网格提供结构先验（形状、姿态、可变形顶点偏移），3DGS 提供高保真渲染。重建分为两个阶段——网格约束的初始优化阶段用 SMAL 顶点位置初始化高斯中心，无约束的体积细化阶段释放高斯以捕捉超出网格分辨率的细节。蒙皮过程将高斯基元绑定到 SMAL 骨骼上，通过操纵 SMAL 参数即可驱动化身执行新姿态。
\textbf{方法信息：} Huawei 参与了这篇论文；项目链接：\url{https://github.com/piotr310100/smal-pets}，当前未确认完整代码。
\textbf{设计动机与边界：} 相对原始 3DGS 以独立高斯和图像重建误差为主的设定，SMAL-pets 把问题落在以下可实现的改动上：SMAL-pets采用混合架构：(1)从单张输入图像生成合成多视角数据集；(2)两阶段重建：第一阶段为网格绑定的初始优化（SMAL参数化模型提供几何和结构先验），第二阶段为无束缚的体素精化（3DGS提供高保真体积渲染）；(3)蒙皮过程将高斯泼溅绑定到底层网格实现新姿态动画。SMAL模型参数包括形状系数、姿态系数、全局平移和可学习顶点偏移。支持通过文本提示进行多模态编辑（外观精化和复杂动画执行）。其设计目的不是单独增加一个模块，而是用上述表示、监督或渲染约束消除原流程在目标任务中的具体失配。提出SMAL-pets框架，将3D高斯泼溅与SMAL参数化模型结合，实现从单张图像生成高质量可编辑动物虚拟形象，支持通过自然语言提示控制外观和动画，提供多模态编辑套件。适用时还应保留论文给出的边界：对非标准姿态的输入图像质量下降（建议先规范化）；SMAL低面数限制细节捕获；依赖合成多视角数据生成质量；仅验证狗类动物。
% METHOD: 2604.02828
\paragraph{NavCrafter（单图可导航场景探索）（2026.04）}
NavCrafter~\cite{2604_02828} 从单图生成可探索 3D 场景。它用 Pluecker ray embeddings 表示相机轨迹，三阶段控制架构把轨迹注入视频 diffusion：3D 卷积适配器编码轨迹，LoRA attention modulation 调整 QKV 投影，随机参考帧增强跨视角一致。碰撞感知规划器选择安全轨迹，迭代生成新视角扩展场景覆盖。最后用增强 3DGS 管线重建，加入深度对齐、结构正则和扩散细化。相对 LiftImage3D，NavCrafter 更强调可控导航和逐步探索。
\textbf{实现细节：} NavCrafter（2026.04） 从视频扩散模型中学习3D先验，旨在解决单图场景探索中的两大挑战：可控性和一致性。其方法包含三个模块：可控新视角合成模块采用三阶段架构（双分支相机注入、LoRA 调制、注意力引导），通过 Pl\"{u}cker 射线嵌入将相机轨迹注入 DiT 去噪器，实现精确的相机控制；碰撞感知轨迹规划器迭代地探索新视角，通过可微变换将物体级坐标系与场景级坐标系对齐；增强版3DGS管线引入深度对齐监督（确保几何精度）和结构正则化（防止过拟合），并结合图像扩散细化模块提升视觉保真度。
\textbf{实现补充：} 实现上还要核对输入、输出、参数、特征、初始化、训练损失、正则、渲染或推理接口，以及优化中的剪枝、分裂、采样、编码或解码步骤；这些抓手决定该方法到底改了 3DGS 管线的哪一环。
\textbf{技术实现补充：} NavCrafter（2026.04） 从视频扩散模型中学习3D先验，旨在解决单图场景探索中的两大挑战：可控性和一致性。其方法包含三个模块：可控新视角合成模块采用三阶段架构（双分支相机注入、LoRA 调制、注意力引导），通过 Pl\"{u}cker 射线嵌入将相机轨迹注入 DiT 去噪器，实现精确的相机控制；碰撞感知轨迹规划器迭代地探索新视角，通过可微变换将物体级坐标系与场景级坐标系对齐；增强版3DGS管线引入深度对齐监督（确保几何精度）和结构正则化（防止过拟合），并结合图像扩散细化模块提升视觉保真度。
\textbf{方法信息：} Huawei 参与了这篇论文。
\textbf{设计动机与边界：} 相对原始 3DGS 以独立高斯和图像重建误差为主的设定，NavCrafter 把问题落在以下可实现的改动上：三模块管线：(1)可控新视角合成：Plucker射线嵌入表示相机轨迹，三阶段控制架构——双分支相机注入（3D卷积适配器编码轨迹嵌入注入视频token）+LoRA注意力调制（轻量编码器生成LoRA控制token调制QKV投影）+随机参考帧增强跨视角一致性；(2)迭代视角合成与碰撞感知规划：碰撞感知轨迹规划器避免场景碰撞并修正轨迹，迭代探索新视角逐步扩展场景覆盖；(3)几何感知3D重建：增强3DGS管线包含深度对齐监督确保几何精度、结构正则化防止过拟合、图像扩散细化提升视觉保真度。其设计目的不是单独增加一个模块，而是用上述表示、监督或渲染约束消除原流程在目标任务中的具体失配。NavCrafter通过多阶段相机控制实现视频扩散模型的精确轨迹控制，碰撞感知规划器和增强3DGS管线实现从单图到高保真3D场景的完整流程，在大视角变化下达到SOTA新视角合成质量。适用时还应保留论文给出的边界：迭代合成和3DGS训练的计算开销较大；对极端遮挡的3D推理有限；扩散模型可能引入不一致的细节。

\section{4D 动态生成}

4D 生成的难点是时间一致和多视角一致同时成立。高斯方法通常先生成环绕视频或多视角视频，再重建 4DGS；或者直接学习高斯 variation field，让静态资产动起来。
% METHOD: 2312.17142
\paragraph{DreamGaussian4D（静态 GS 加 HexPlane 变形）（2023.12）}
DreamGaussian4D~\cite{2312_17142} 将 DreamGaussian 扩展到 4D。它先用改进的 DreamGaussianHD 从单图生成静态 3DGS，再优化 HexPlane deformation field 预测每个时间戳的高斯位移；运动来自生成视频，而不是直接用 video SDS。最后用 image-to-video diffusion 在 UV 空间进行时间一致纹理精炼。相对 Diffusion4D 的训练式视频扩散，DreamGaussian4D 更像优化式 pipeline：先静态，再变形，再纹理。
第一阶段的 DreamGaussianHD 仍是单图到静态 3DGS，但优化策略比 DreamGaussian 更重视未见区域：每次采样更多视角参与 SDS/重建约束，并用固定白背景减少扩散噪声对透明区域的干扰。第二阶段把静态高斯放进 HexPlane 变形场，网络输入时间戳后预测每个高斯的位移或变形参数，让同一套规范高斯在不同时间形成动画。运动监督来自 image-to-video 模型生成的视频，而不是每一步直接对 4D 表示施加 video SDS；随后在 UV 空间用视频扩散做纹理精炼，保证同一表面在时间上不闪烁。
\textbf{锚点细化：} DreamGaussian4D：针对 MAV3D、4D-fy 等依赖隐式 NeRF 和视频 SDS 导致生成 4D 场景需数小时且运动难控的局限，用\textbf{显式 4D 表示（3DGS+HexPlane 形变场）+ 驱动视频显式监督}替代隐式 NeRF+视频 SDS，核心思想是\textbf{用 HexPlane 将 4D 时空场分解为六个特征平面，把复杂度从 $O(N^4)$ 降至 $O(N^2)$}——参考视角用驱动视频 MSE 损失提供运动监督、其他视角用 Zero-1-to-3 的 SDS 损失传播运动到未观测区域，零初始化变形网络最后几层使初始动态与静态一致，10 分钟内生成 14 帧 4D 动画。
\textbf{方法信息：} 代码：\url{https://github.com/jiawei-ren/dreamgaussian4d}；GitHub 关注度高，约 617 stars；引用很多，Semantic Scholar 约 220 次。



% METHOD: 2405.16645
\paragraph{Diffusion4D（视频扩散到显式 4D 高斯）（2024.05）}
Diffusion4D~\cite{2405_16645} 是 4D 内容生成的重要锚点。DreamGaussian4D 这类方法通常先做静态 3D，再用视频或变形场把它动起来；Diffusion4D 的判断是，4D 生成最缺的是“同一对象在时间和视角上同时一致”的监督，因此不应只依赖逐样本 SDS，而应训练一个能一次输出环绕动态视频的 4D-aware video diffusion model。

具体流程先策展数据：从 Objaverse-1.0 和 Objaverse-XL 中筛出动态 3D 资产，用固定前视时间帧的 SSIM 判断运动是否太弱或太剧烈，再渲染 24 帧环绕视频作为训练样本。模型输入可以是文本、单图或静态 3D 内容，输出是相机绕动态对象一圈运动的 orbital video；相机方位角随时间均匀变化，所以视频的每一帧同时对应一个时间点和一个新视角。为了让用户能控制“动得多大”，方法定义 3D-to-4D motion magnitude，并用运动幅度重建损失和 3D-aware classifier-free guidance 强化动态强度和几何一致性。得到多视角一致的视频后，系统再用 coarse-to-fine 的 Gaussian splatting 优化显式 4D 表示。相对 DreamGaussian4D 的“单图静态资产 + 变形优化”，Diffusion4D 把时空一致性提前放进视频扩散模型，再让 3DGS 承接重建。
\textbf{实现细节：} 与 DreamGaussian4D 的 SDS 优化策略不同，Diffusion4D（2024.05） 和 EG4D（2024.05） 提出无需分数蒸馏（score distillation）的4D生成方法。EG4D 同样不依赖 SDS，而是通过显式生成多视角视频，利用注意力注入策略保证时序一致，再基于3DGS的鲁棒动态重建方法生成高质量4D资产。
\textbf{实现补充：} 实现上还要核对输入、输出、参数、特征、初始化、训练损失、正则、渲染或推理接口，以及优化中的剪枝、分裂、采样、编码或解码步骤；这些抓手决定该方法到底改了 3DGS 管线的哪一环。
\textbf{设计动机与边界：} 相对原始 3DGS 以独立高斯和图像重建误差为主的设定，Diffusion4D 把问题落在以下可实现的改动上：Diffusion4D的核心思路是将视频扩散模型的时间一致性迁移为4D生成所需的时空一致性。方法分为三个阶段：1）数据策展：从Objaverse-1.0和Objaverse-XL中筛选高质量动态3D资产。使用SSIM评估时序动态性，移除运动过于微弱或过于剧烈的资产，保留运动幅度适中的资产。2）4D感知视频扩散模型：在策展数据集上训练，输入为文本/单图/静态3D内容，输出为围绕动态3D资产的环绕视频（T个多视角图像，相机沿均匀递增的方位角轨迹运动）。引入3D-to-4D运动幅度度量来控制动态强度，设计运动幅度重建损失确保生成内容与输入运动提示一致。使用3D感知分类器自由引导（classifier-free guidance）增强几何一致性。3）显式4D构建：利用生成的多视角一致环绕视频，通过高斯溅射从粗到精地构建4D表示。粗阶段快速建立初始4D结构，精阶段添加细节和优化动态。整个过程仅需几分钟。其设计目的不是单独增加一个模块，而是用上述表示、监督或渲染约束消除原流程在目标任务中的具体失配。提出Diffusion4D框架，通过精心策展的动态3D数据集训练4D感知视频扩散模型，生成动态3D资产的环绕视角视频。引入3D到4D运动幅度度量作为引导，提出运动幅度重建损失和3D感知分类器自由引导优化动态学习。最终以从粗到精的方式用高斯溅射进行显式4D构建，几分钟内生成高质量多样化4D资产。适用时还应保留论文给出的边界：4D感知视频扩散模型的训练需要大规模策展数据集；对复杂运动的建模仍有提升空间；生成的4D资产分辨率有限。
\textbf{方法信息：} Microsoft 参与了这篇论文；引用很多，Semantic Scholar 约 110 次。



% METHOD: 2405.18132
\paragraph{EG4D（不用 SDS 的显式 4D 生成）（2024.05）}
EG4D~\cite{2405_18132} 也称 DG4D，它和 Diffusion4D 一样反对把 4D 生成完全交给 SDS，但切入点更偏单图到 4D。DreamGaussian4D 用单图先生成静态高斯，再依赖视频先验和变形场；EG4D 的改变是先显式生成多视角动态视频，把看不见的视角和动作都变成可监督图像，再用 4DGS 重建。

实现上，第一阶段用 SVD-XL 从参考图生成视频，再用 SV3D 产生多视角结果。为了避免不同时间帧各自扩散导致身份漂移，EG4D 在空间注意力里注入跨帧 latent 的指数移动平均，让当前多视角推理能感知历史时间信息。第二阶段用 4DGS 做 coarse reconstruction，代码中对应 \texttt{train.py} 和 \texttt{sv4d.py} 配置；这里不仅优化几何变形，还加入颜色变换来吸收视频扩散带来的帧间纹理差异，避免把颜色闪烁误当成几何运动。第三阶段用扩散先验做 semantic refinement，修复生成视频中的语义缺陷并继续细化 4D 高斯。相对 Diffusion4D 的训练式 4D 视频模型，EG4D 更像一条工程化显式管线：复用现有视频/多视角扩散模型，用注意力注入和 4DGS 优化把它们粘成稳定资产。

从优化流程看，EG4D 的输入是一张参考图，输出是可以从新视角和新时间渲染的动态高斯资产。它先把 SVD-XL/SV3D 生成的视频整理成多时间、多视角监督，再初始化规范空间中的高斯集合；渲染某个时间戳时，变形模块预测每个高斯的位移、旋转和尺度变化，颜色变换模块则修正同一高斯在不同帧中的外观偏移。训练损失主要约束渲染图像与生成多视角帧的一致性，并通过时间邻近约束抑制运动抖动。这样，几何运动由 4DGS 变形场承担，视频扩散只负责提供可见监督和后期语义 refinement，而不是像 SDS 那样在每一步用噪声梯度直接牵引整个 4D 表示。
\textbf{实现细节：} 与 DreamGaussian4D 的 SDS 优化策略不同，Diffusion4D（2024.05） 和 EG4D（2024.05） 提出无需分数蒸馏（score distillation）的4D生成方法。EG4D 同样不依赖 SDS，而是通过显式生成多视角视频，利用注意力注入策略保证时序一致，再基于3DGS的鲁棒动态重建方法生成高质量4D资产。
\textbf{方法信息：} 代码：\url{https://github.com/jasongzy/eg4d}。
\textbf{设计动机与边界：} 相对原始 3DGS 以独立高斯和图像重建误差为主的设定，EG4D 把问题落在以下可实现的改动上：EG4D（也称DG4D）是一个多阶段4D生成框架。（1）时间一致多视角视频合成：从单张输入图像，利用视频扩散模型生成多视角视频。为解决帧间不一致问题，设计注意力注入机制，让每次多视角扩散推理通过跨帧潜在指数移动平均（EMA）感知时间信息。这种无需训练的策略有效缓解视频级别的不一致性。（2）4DGS动态重建：选择4D高斯泼溅作为4D表示，利用其高效训练和渲染能力。针对现有4DGS方法假设外观变化仅由几何变形导致的局限性，提出更鲁棒的重建方法处理帧间纹理差异。（3）语义修复精炼：使用扩散先验修复生成结果中的语义缺陷，提升视觉质量。整个框架完全使用开源工具实现。其设计目的不是单独增加一个模块，而是用上述表示、监督或渲染约束消除原流程在目标任务中的具体失配。提出首个无需分数蒸馏（SDS）的高质量4D生成框架。通过显式生成多视角视频替代隐式SDS蒸馏，避免过饱和和Janus问题。注意力注入策略实现时间一致性多视角视频合成。基于4DGS的动态重建高效且高质量。用户偏好研究显示生成质量大幅超越基线。适用时还应保留论文给出的边界：依赖视频扩散模型的生成质量；多阶段流程增加处理时间；对复杂运动模式效果有限。
% METHOD: 2504.02045
\paragraph{WorldPrompter（360 度视频作为 3D 场景中间表示）（2025.04）}
Generating 360 Video is What You Need For a 3D Scene~\cite{2504_02045} 提出先生成完整 360 度全景视频，再用前馈重建器转成 Gaussian scene。FastScene 先生成全景图，再沿相机轨迹逐步修补新视角；WorldPrompter 的判断是，逐步 outpainting 会不断积累全局失真，因此应该一次性生成覆盖完整环境的全景视频。视频生成器根据文本产生一段模拟人在虚拟环境中行走拍摄的全景序列，训练时混合图像和视频数据来同时保持空间内容和时间连贯。

有了全景视频后，系统再用前馈 3D 重建器快速恢复 Gaussian scene。这里 360 度视频不是最终输出，而是全局一致的中间表示：它把室内或街景的环绕内容、相机移动和遮挡变化先组织成一段连续观测，再交给 3DGS 承载可导航资产。相对逐视角扩展方法，WorldPrompter 的增量是把“场景全局布局”提前交给视频扩散一次性决定。
\textbf{实现细节：} WorldPrompter（2025.04） 验证了360度全景视频作为3D场景生成中间表示的有效性。该方法收集了约1700个全景视频和20万+全景图像数据集，训练基于 DiT 的全景视频生成器，可生成长达10.6秒（128帧）的360度全景视频，模拟人在虚拟环境中行走的效果。关键设计包括混合训练策略（混合室内全景视频和户外全景图像增强泛化能力）和掩码扩散损失（确保生成无伪影）。生成的全景视频经过 COLMAP 估计相机姿态，再通过 Long-LRM 前馈3D重建器转化为高斯泼溅。
\textbf{方法信息：} Adobe 参与了这篇论文。
\textbf{设计动机与边界：} 相对原始 3DGS 以独立高斯和图像重建误差为主的设定，Generating 360° Video is What You Need For a 3D Scene 把问题落在以下可实现的改动上：WorldPrompter管线包含两个核心组件：(1)条件化360度全景视频生成器：使用混合图像和视频数据训练的扩散模型，根据文本提示生成128帧全景视频，模拟人在虚拟环境中行走并拍摄。采用图像和视频数据混合训练策略以实现空间和时间一致性。(2)前馈3D重建器：将生成的全景视频快速重建为高斯泼溅表示。关键创新在于放弃传统的逐步扩展（outpainting）方法，转而一次性生成完整场景，确保全局视觉内容的一致性和无失真。该方法避免了逐视角生成导致的全局失真和不一致问题。其设计目的不是单独增加一个模块，而是用上述表示、监督或渲染约束消除原流程在目标任务中的具体失配。提出以360度全景视频作为中间场景表示的生成管线WorldPrompter，通过条件化360度全景视频生成器合成128帧视频模拟人在虚拟环境中的行走，再通过前馈3D重建器重建为高斯泼溅，实现可自由导航的完整3D场景生成。COLMAP匹配率达94.6%，超越现有360度视频生成器和3D场景生成模型。适用时还应保留论文给出的边界：主要面向静态场景；全景视频生成可能受限于训练数据的多样性；场景空间范围仍受全景视频覆盖限制。
% METHOD: 2507.23785
\paragraph{Gaussian Variation Field Diffusion（变化场 VAE 到 4D）（2025.07）}
Gaussian Variation Field Diffusion~\cite{2507_23785} 不逐实例拟合 4DGS，而是从 3D 动画数据直接编码规范高斯和属性变化场。Direct 4DMesh-to-GS Variation Field VAE 使用 mesh-to-GS autoencoder 得到规范 GS，并用 Perceiver-style transformer 把动画序列运动聚合成固定 latent；随后时间感知 DiT 以输入视频和规范 GS 为条件生成 variation field。相对 Diffusion4D 的“视频生成后重建”，该方法把 4DGS 的变化场本身作为扩散对象。
它的 VAE 分两部分编码：一部分是规范 3DGS，提供静态形状和外观；另一部分是 Gaussian Variation Fields，记录位置、旋转、尺度或外观随时间的变化。规范高斯由预训练 mesh-to-GS autoencoder 从 3D 动画资产中获得，运动序列则由 Perceiver-style transformer 聚合成固定长度 latent，并用 mesh-guided loss 让高斯运动贴合 mesh 真值。生成时，时间感知 DiT 不再生成视频帧，而是在输入视频和规范 GS 条件下直接生成 variation field。这样输出天然是 4DGS 的时变参数，而不是先得到视频再反解为动态高斯。
\textbf{方法信息：} Microsoft 参与了这篇论文。
\textbf{设计动机与边界：} 相对原始 3DGS 以独立高斯和图像重建误差为主的设定，Gaussian Variation Field Diffusion for High-fidelity Video-to-4D Synthesis 把问题落在以下可实现的改动上：方法包含两个核心组件：(1) Direct 4DMesh-to-GS Variation Field VAE：将规范3DGS表示与高斯属性时序变化场（Gaussian Variation Fields）分别编码。利用预训练mesh-to-GS自编码器获取规范GS，通过Perceiver-style transformer聚合3D动画序列的运动信息为固定长度潜表示，引入mesh-guided loss对齐GS运动与mesh GT运动。(2) Gaussian Variation Field扩散模型：基于时间感知DiT，以输入视频和规范GS为条件，在紧凑潜空间中学习变化场分布。仅使用Objaverse数据集中精选的可动画3D物体进行训练，但对野外视频输入展现出强泛化能力。核心创新在于绕过昂贵的逐实例4DGS拟合过程，直接从mesh动画数据高效编码高斯变化场。其设计目的不是单独增加一个模块，而是用上述表示、监督或渲染约束消除原流程在目标任务中的具体失配。提出Direct 4DMesh-to-GS Variation Field VAE，直接从3D动画数据编码规范高斯场及其时序变化，无需逐实例拟合；训练基于时间感知Diffusion Transformer的高斯变化场扩散模型，仅用合成数据训练即可泛化到野外视频输入，实现高保真Video-to-4D生成。适用时还应保留论文给出的边界：依赖Objaverse数据集的合成3D动画数据训练，生成质量受限于训练数据的多样性和质量；仅处理单体物体的4D生成，未涉及复杂场景。
% METHOD: 2508.07409
\paragraph{CharacterShot（姿态控制 4D 角色动画）（2025.08）}
CharacterShot~\cite{2508_07409} 从单张角色图和 2D pose sequence 生成 4D 角色动画。Track4DGen 一类方法更依赖跟踪器或视频中的隐式运动线索，CharacterShot 的改变是把姿态序列作为显式控制条件，让用户指定角色怎样动。第一阶段基于 CogVideoX 训练姿态控制 2D 动画：输入角色图和姿态序列，姿态通过 3D VAE 压到视频 latent 中，使视频生成器在保持角色身份的同时跟随骨架动作。

第二阶段把单视角动画扩展成多视角视频。方法加入双注意力模块，同时在时间帧和相机视角之间交换信息，并结合相机先验约束不同视角看到的是同一个角色，而不是彼此独立的视频。第三阶段才进入 4DGS：多视角、多时间帧作为监督，优化随时间变化的高斯表示；邻居约束要求相邻时间和相邻视角下的高斯变化连续，减少动画闪烁和视角跳变。相对先生成视频再粗略重建的流程，CharacterShot 的重点是把“姿态可控性”一路传到多视角生成和 4DGS 优化中。
\textbf{实现细节：} CharacterShot（2025.08） 的核心贡献在于解决了4D角色动画生成中一个长期被忽视的问题——可控性：从用户提供的单张角色图像和2D姿态序列出发，生成连贯、可控的动态3D角色动画。前述4D生成方法（如 Stag4D、SV4D）的典型管线是：先用扩散模型生成单视角角色视频，再基于该视频优化4D高斯表示。然而，这条管线存在两个根本缺陷。第一，从自定义角色图像生成特定运动的视频本身就极其复杂且代价高昂，普通用户难以获得满意的驱动视频。第二，即使获得了多视角视频，这些视频之间往往存在细微的跨视角不对齐（slight misalignment），导致4D高斯优化过程中出现离群点和噪声——具体表现为角色手部突然消失或可见的几何伪影，这正是逐帧独立优化缺乏几何约束的后果。CharacterShot 的三阶段管线。为解决上述问题，CharacterShot 设计了从姿态控制到4D输出的端到端框架，分为三个阶段： 第一阶段是姿态驱动的 2D 角色动画。
\textbf{技术实现补充：} CharacterShot（2025.08） 的核心贡献在于解决了4D角色动画生成中一个长期被忽视的问题——可控性：从用户提供的单张角色图像和2D姿态序列出发，生成连贯、可控的动态3D角色动画。Baseline 做法及其局限。前述4D生成方法（如 Stag4D、SV4D）的典型管线是：先用扩散模型生成单视角角色视频，再基于该视频优化4D高斯表示。然而，这条管线存在两个根本缺陷。第一，从自定义角色图像生成特定运动的视频本身就极其复杂且代价高昂，普通用户难以获得满意的驱动视频。第二，即使获得了多视角视频，这些视频之间往往存在细微的跨视角不对齐（slight misalignment），导致4D高斯优化过程中出现离群点和噪声——具体表现为角色手部突然消失或可见的几何伪影，这正是逐帧独立优化缺乏几何约束的后果。CharacterShot 的三阶段管线。为解决上述问题，CharacterShot 设计了从姿态控制到4D输出的端到端框架，分为三个阶段： \textit{阶段一：姿态驱动的2D角色动画。}以 DiT（Diffusion Transformer）架构的图生视频模型 CogVideoX 为基座，将姿态信息编码为与视频隐变量同维度的条件隐变量 $z_p \in \mathbb{R}^{f \times h \times w \times 16}$，并与视频隐变量 $z_i$ 和参考图像隐变量 $z_r$ 在通道维度拼接： 其中 $z_p'$ 是参考图像对应的姿态隐变量，$f$ 为帧数，$h \times w$ 为空间分辨率。训练时随机替换参考帧及其姿态，增强模型对姿态输入不对齐的鲁棒性。\textit{阶段二：多视角视频生成。}将阶段一的2D输出扩展到 $V$ 个视角，引入相机编码器将 Pl\"ucker 嵌入（一种将相机内外参编码为 $6 \times 8h \times 8w$ 特征图的方法）编码为相机 token，注入每个视角的输入序列。双注意力模块（dual-attention module）是此阶段的关键创新：先前方法（如 SV4D）分别计算空间注意力、时间注意力和视角注意力，这种分离式设计不利于跨维度信息的隐式传递。
\textbf{方法信息：} 代码：\url{https://github.com/jeoyal/charactershot}。
\textbf{设计动机与边界：} 相对原始 3DGS 以独立高斯和图像重建误差为主的设定，CharacterShot 把问题落在以下可实现的改动上：CharacterShot框架包含三个阶段。(1)姿态控制2D角色动画：基于CogVideoX(DiT)预训练，将姿态序列通过3D VAE压缩为潜在表示并与视频潜在连接作为条件，使模型支持任意2D姿态序列作为控制信号。(2)多视角视频生成：引入双注意力模块和相机先验，将2D动画模型提升为3D，生成具有时空和空间视角一致性的多视角视频。双注意力模块同时处理时间帧和空间视角的注意力计算。(3)邻居约束4DGS优化：在多视角视频上进行新颖的邻居约束4D高斯溅射优化，利用时空和视角邻域约束确保4D表示的连续性和稳定性。同时构建Character4D大规模数据集，包含13115个具有多样外观和动作的独特角色，从多视角渲染。其设计目的不是单独增加一个模块，而是用上述表示、监督或渲染约束消除原流程在目标任务中的具体失配。提出CharacterShot框架，从单张参考角色图像和2D姿态序列生成可控一致的4D角色动画。预训练DiT基础2D角色动画模型，引入双注意力模块和多视角视频生成，设计邻居约束4D高斯溅射优化，构建Character4D数据集（13115个角色）。适用时还应保留论文给出的边界：多阶段流水线复杂度高；对极端姿态的泛化性有限；4DGS优化的计算成本较大。

\section{其他 3D 生成方法}

这里收纳城市、点云、地球尺度和多模态预训练等方向。它们说明 3DGS 生成并不限于单物体，关键是找到适合大规模结构的中间表示：城市更适合 BEV 或 patch 序列，地球尺度需要结构和纹理分开压缩，点云上色需要先有几何锚定，多模态预训练则要让高斯属性进入语言-图像空间。读这些方法时，应先问它们把高斯放在生成链路的哪个位置。
% METHOD: 2406.06526
\paragraph{GaussianCity（无界城市的 BEV-Point 高斯生成）（2024.06）}
Generative Gaussian Splatting for Unbounded 3D City Generation~\cite{2406_06526} 也称 GaussianCity。普通 3DGS 面向有限场景时可以把所有高斯都放进显存，但无界城市生成会让高斯数量随地图面积增长。GaussianCity 的核心改变是引入 BEV-Point 表示：位置相关属性压缩到 BEV 地图，风格相关属性放到 lookup table；渲染时只激活当前相机可见的 BEV points，因此显存主要取决于视野，而不是整座城市的面积。

生成流程从城市条件图开始，例如语义图、高度场和密度图。系统先生成 BEV points，并用射线相交过滤不可见点；随后 Point Serializer 同时读取局部 3D patch 和全局上下文，把道路、建筑高度、街区密度等结构信息编码成序列，再解码为每个可见点的 Gaussian 属性。这样，BEV 决定城市平面和高度结构，lookup table 控制外观风格，3DGS rasterization 负责实时可视化。相对 ScenDi 的体素 latent 城市扩散，GaussianCity 更像把城市生成压到 BEV 空间中，再按可见性动态展开成高斯。
\textbf{实现细节：} 当生成目标从单个物体扩展到城市级别时，3D高斯的存储和计算面临量级挑战——无界场景要求将高斯数量扩展到十亿级，VRAM需求可达数百GB。现有的NeRF方法（如InfiniCity和CityDreamer）沿射线以相同密度采样并聚合所有点的颜色值，导致推理效率低下且细节丢失；而基于3DGS的方法（如DreamGaussian、GaussianDreamer、GaussianEditor）仅生成有限尺度的物体或场景。GaussianCity（2024.06） 是首个面向无界3D城市生成的3DGS生成模型，其核心洞察是引入BEV-Point（鸟瞰点）作为高度紧凑的中间表示，将高斯属性解耦为位置相关属性和风格相关属性两部分。BEV-Point表示。在BEV-Point表示中，只有当前帧可见的鸟瞰点被保留，因此VRAM使用量不随场景尺度增长。具体实现上，鸟瞰语义图 $S$ 和高度场 $H$ 定义BEV点坐标： 其中二值密度图 $D$ 根据语义类别调整采样密度（道路/水面等简单纹理区域降低密度，建筑立面等复杂纹理区域增加密度）。
\textbf{设计动机与边界：} 相对原始 3DGS 以独立高斯和图像重建误差为主的设定，Generative Gaussian Splatting for Unbounded 3D City Generation 把问题落在以下可实现的改动上：GaussianCity解决将3DGS从有限尺度扩展到无限尺度城市生成的挑战。核心创新：(1) BEV-Point表示：将高斯属性解耦为位置相关属性（可压缩为BEV地图）和风格相关属性（压缩为风格查找表）。仅可见BEV点参与渲染和优化，VRAM使用保持恒定。通过语义图、高度场和密度图生成BEV点，再用射线交叉过滤不可见点。(2) 空间感知BEV-Point解码器：利用Point Serializer捕获BEV点的结构特征（局部3D patch采样）和上下文特征（全局特征聚合）。(3) 生成管线：从噪声或语义图生成BEV-Point特征，再解码为3D高斯属性。其设计目的不是单独增加一个模块，而是用上述表示、监督或渲染约束消除原流程在目标任务中的具体失配。提出首个面向无界3D城市生成的高斯泼溅框架GaussianCity。核心创新：(1) BEV-Point作为紧凑中间表示，确保VRAM使用不随场景规模增长而增加；(2) 空间感知高斯属性解码器，利用Point Serializer整合BEV点的结构和上下文特征。相比CityDreamer速度提升60倍。适用时还应保留论文给出的边界：生成质量受BEV-Point分辨率限制；对复杂建筑结构的细节建模仍有提升空间。
\textbf{方法信息：} 代码：\url{https://github.com/hzxie/gaussiancity}；GitHub 关注度较高，约 332 stars。



% METHOD: 2502.17860
\paragraph{UniGS（语言-图像-3D 高斯预训练）（2025.02）}
UniGS~\cite{2502_17860} 不属于常规“从文本生成一个 3D 物体”的路线，而是把 3DGS 引入语言-图像-3D 预训练。它针对的是 Uni3D、ULIP 等点云预训练的表示瓶颈：点云只有离散几何点，缺少颜色、透明度、尺度、旋转和可渲染外观；而 2D 图像是密集像素，两者之间存在明显模态差距。UniGS 的改变是用优化后的 3DGS 替代点云作为 3D 输入，使三维表示既有几何，又能通过 splatting 和图像空间建立对应。

训练流程先固定一个预训练视觉语言模型，例如 CLIP 的图像编码器和文本编码器，作为共享的语言-图像目标空间。每个训练样本包含文本、图像和对应 3DGS。3D 分支需要把高斯集合编码成一个 embedding，并通过 Language-3DGS Alignment 和 Image-3DGS Alignment 两个对比损失，分别拉近文本-高斯、图像-高斯的表示距离。这样，训练完成后，文本可以检索 3DGS，图像也可以对齐 3DGS，3D 表示能够参与开放词表分类和检索。

核心模块是 Gaussian-Aware Guidance。直接把 3DGS 的中心位置当点云送进 Uni3D 并不充分，因为高斯不一定落在真实物体表面，而且每个高斯还有颜色、不透明度、尺度、旋转等属性。UniGS 使用并行双分支 3D encoder：基础分支利用冻结的点云编码器提供空间先验，高级分支处理高斯的完整显式属性；高级分支通过 cross-attention 从基础分支聚合先验，以避免新属性学习和已有点云先验互相冲突。这个设计让模型既利用点云编码器的空间归纳偏置，又能读到 3DGS 特有的外观和形状信息。

数据上，UniGS 可以从多视图图像或 COLMAP 数据优化 3DGS，也可以用已有点云数据集初始化高斯位置。论文在 Objaverse、ABO、MVImgNet 和 SUN RGBD 上验证，报告相对 Uni3D 在零样本分类、文本检索和开放世界理解任务上有明显提升。相对本章其他生成方法，UniGS 的作用不是直接产出资产，而是把 3DGS 变成可被语言和图像共同索引的基础表征。
\textbf{实现细节：} UniGS（2025.02） 提出统一的语言-图像-3D 预训练框架，利用 3DGS 作为 3D 表示桥接多模态学习。核心思想是将每个高斯基元的全部属性（位置、颜色、缩放、旋转、不透明度）拼接为高维特征 $g = (\mu, c, \alpha, s, R)$，通过跨模态对比学习在语言、图像和3D三个模态间建立共享嵌入空间。高斯感知引导（Gaussian-Aware Guidance）。UniGS的关键创新是利用3DGS渲染过程中的固有信息指导表征学习：在渲染管线中，高斯基元经投影和$\alpha$-混合产生像素级颜色，这一可微过程天然将3D几何与2D图像关联。UniGS通过高斯感知引导模块将渲染梯度反向传播到3D表征，使模型不仅学习到高斯属性的统计分布，还能理解属性之间的几何语义关系（如缩放和旋转如何共同影响渲染结果）。数据规模与泛化性。UniGS通过集成现有3D数据集（ShapeNet、Objaverse等）构建了大规模3DGS训练语料，支持从10000个3D点的密集采样。
\textbf{设计动机与边界：} 相对原始 3DGS 以独立高斯和图像重建误差为主的设定，UniGS 把问题落在以下可实现的改动上：UniGS用3DGS替代传统点云作为3D表示。3DGS将3D世界建模为带颜色和不透明度的3D高斯集合，建立了与2D图像的自然连接（通过泼溅渲染）。方法分两阶段：(1)从预训练的视觉语言模型建立共享视觉-文本空间；(2)使用并行结构双分支3D编码器（冻结的点云基础编码器提供先验+高级编码器建模高层信息）将3DGS与语言-图像表示对齐。核心创新是Gaussian-Aware Guidance模块：利用高斯的显式空间属性（位置、颜色、不透明度、尺度、旋转）指导细粒度3D特征学习，解决3DGS不一定分布在物体表面导致的特征捕捉不足问题。支持从多视图图像或COLMAP数据高效优化3DGS，现有点云数据集可直接作为高斯位置初始化。其设计目的不是单独增加一个模块，而是用上述表示、监督或渲染约束消除原流程在目标任务中的具体失配。首次将3DGS作为3D表示引入语言-图像-3D预训练，提出Gaussian-Aware Guidance模块增强跨模态对齐，在零样本分类(+9.36%)、文本检索(+4.3%)和开放世界理解(+7.92%)上大幅超越Uni3D。适用时还应保留论文给出的边界：3DGS优化增加预处理开销；对大规模3D数据集的3DGS优化可能耗时；基础编码器的容量限制上限性能。
\textbf{方法信息：} Huawei 参与了这篇论文；代码：\url{https://github.com/li-hao-yuan/unigs}；仓库目前约 38 stars。
% METHOD: 2507.16535
\paragraph{EarthCrafter（双稀疏地球尺度生成）（2025.07）}
EarthCrafter~\cite{2507_16535} 面向大规模地理场景。它构建 Aerial-Earth3D 航空数据集，并用双稀疏架构分离结构和纹理。StructVAE 压缩几何体素，TexVAE 压缩体素化 2DGS；条件感知 flow matching 模型分别生成结构 latent 和纹理 latent，支持语义、图像或无条件地形生成。相对 GaussianCity 的 BEV-Point 城市生成，EarthCrafter 更偏地球表面尺度，并显式分开结构体素和 2DGS 纹理。

数据层面，Aerial-Earth3D 把大范围航空区域切成固定尺度场景，并为多视图 Google Earth 帧补充姿态、深度、法线、语义分割和相机位姿；结构真值来自重建网格的体素化。表示层面，StructVAE 将全分辨率结构体素压缩到稀疏 latent，利用伪稀疏到稀疏模块做精确上采样；TexVAE 则处理体素化 2DGS，专门承载外观纹理。生成时，结构分支决定地形、建筑和道路的空间骨架，纹理分支补充可渲染外观，最后解码为可浏览的大范围 3D 场景。这个拆分让模型不必在同一个 latent 里同时记住地貌拓扑和细节纹理，也使语义条件、图像条件和无条件采样可以接到同一套渲染输出上。
\textbf{方法信息：} Alibaba 参与了这篇论文。
\textbf{设计动机与边界：} 相对原始 3DGS 以独立高斯和图像重建误差为主的设定，EarthCrafter 把问题落在以下可实现的改动上：EarthCrafter的双创新：(1)Aerial-Earth3D数据集：50,028个600mx600m场景，覆盖美国大陆，45M多视图Google Earth帧。包含姿态标注、深度图、法线、语义分割和相机位姿。通过启发式相机姿态设计和InstantNGP重建3D网格，体素化作为结构真值；(2)StructVAE：利用空间压缩结构潜空间，将全分辨率体素坐标V\_c压缩到1/256（c\_s=32）。创新性地提出伪稀疏到稀疏(PSS)块，通过稀疏像素重排和体素分类实现精确上采样，达到97.1%结构重建精度；(3)TexVAE：将体素化2DGS压缩到紧凑潜空间；(4)条件感知流匹配模型：StructFM生成结构潜变量，TexFM生成纹理潜变量，支持混合输入条件（语义、图像或无条件）的灵活生成。其设计目的不是单独增加一个模块，而是用上述表示、监督或渲染约束消除原流程在目标任务中的具体失配。提出了Aerial-Earth3D（迄今最大3D航空数据集，50k场景、45M多视图帧）和EarthCrafter框架。核心创新是双稀疏架构分离结构和纹理生成：StructVAE压缩几何体素（1/256压缩），TexVAE压缩体素化2DGS，条件感知流匹配模型支持灵活生成。支持语义引导城市布局生成和无条件地形合成。适用时还应保留论文给出的边界：依赖Google Earth数据可能存在时效性限制；600m场景粒度可能对某些应用不够精细；双VAE训练需要大量计算资源。
% METHOD: 2508.05631
\paragraph{GAP（无颜色点云到高斯）（2025.08）}
GAP~\cite{2508_05631} 从无颜色原始点云生成高保真 3D Gaussians。它先用 CAP-UDF 学习 unsigned distance field 和点法线，用 2DGS 初始化表面对齐的高斯盘；再通过多视角优化和深度感知文本到图像扩散生成一致外观；难观测区域由扩散修复补全。表面锚定机制约束高斯位于 3D 形状表面。相对文本到 3D，它的输入是已有几何点云，目标是“给点云上色并高斯化”。
第一步先把点云变成可放置高斯的表面。CAP-UDF 从无颜色点云学习 unsigned distance field，并推导表面法线；GAP 使用 2DGS 而不是普通 3DGS 初始化高斯盘，因为 2D 高斯盘可以直接对齐局部切平面，旋转矩阵由法线确定。第二步做渐进式外观生成：从多个视角调用深度感知文本到图像扩散模型，合成与已有几何一致的外观图，再通过高斯优化把颜色和不透明度写回表面对齐高斯。第三步针对遮挡和难观测区域，用扩散修复补全缺失外观。表面锚定约束贯穿优化，防止扩散生成的纹理把高斯拉离输入点云形状。
\textbf{实现细节：} 原始3D点云是最常见、最易获取的3D数据形式——LiDAR、深度相机等设备输出的就是点云。然而，这些点云既没有颜色属性，也缺乏拓扑连接信息，无法直接用于高质量渲染。传统做法是先将点云重建为网格、再进行UV展开和纹理贴图，但自动UV展开常导致纹理重叠、碎片化和扭曲，严重制约最终渲染质量。3DGS天然免去了UV参数化的需要，因此"将原始无颜色点云直接转换为高保真3D高斯"成为一个极具吸引力的目标。GAP（2025.08） 正是为此设计的方法：给定文本提示词$c$和点云$P = \{p_i\}_{i=1}^N$，目标是生成高保真3D高斯$G = \{g_i\}_{i=1}^M$。高斯初始化。GAP的初始化策略充分利用了点云固有的几何信息。首先，高斯中心直接设为输入点云坐标，提供了对3D形状表面的粗略几何描述。在此基础上，方法利用CAP-UDF从点云学习一个神经无符号距离场（UDF）$f_u$，并通过梯度推断法线：$n_i = \nabla f_u(p_i) / \|\nabla f_u(p_i)\|$。值得注意的是，GAP选择2DGS（oriented disk）而非标准3DGS作为表示——2DGS将3D高斯椭球体替换为2D定向盘片，天然将法线编码为盘片方向，在表达精细局部几何时性能更优。每个盘片的旋转矩阵根据推断的法线初始化，确保盘片方向与局部表面对齐。渐进式多视角生成。核心挑战在于：点云没有颜色，高斯的外观从何而来？GAP的答案是利用文本引导的深度感知扩散模型逐视角生成外观图像，作为高斯优化的监督信号。具体而言，对于预设的视角序列$\{v_j\}_{j=1}^K$，方法渐进地为每个视角生成外观：利用学习到的UDF场通过射线行进计算深度图$D_j$，将当前渲染图像$I_j$、深度图$D_j$、掩码$M_j$和文本提示$c$一起送入基于ControlNet的深度感知修复扩散模型，生成几何一致的外观图像。这里的关键设计是三区域掩码策略——掩码被划分为生成区$M_{\mathrm{generate}}$（从未生成过的空白区域）、保留区$M_{\mathrm{keep}}$（已处理且当前视角无更优观察角度的区域）和更新区$M_{\mathrm{update}}$（当前视角提供更优观察角度的区域）。
\textbf{技术实现补充：} 原始3D点云是最常见、最易获取的3D数据形式——LiDAR、深度相机等设备输出的就是点云。然而，这些点云既没有颜色属性，也缺乏拓扑连接信息，无法直接用于高质量渲染。传统做法是先将点云重建为网格、再进行UV展开和纹理贴图，但自动UV展开常导致纹理重叠、碎片化和扭曲，严重制约最终渲染质量。3DGS天然免去了UV参数化的需要，因此"将原始无颜色点云直接转换为高保真3D高斯"成为一个极具吸引力的目标。GAP（2025.08） 正是为此设计的方法：给定文本提示词$c$和点云$P = \{p_i\}_{i=1}^N$，目标是生成高保真3D高斯$G = \{g_i\}_{i=1}^M$。高斯初始化。GAP的初始化策略充分利用了点云固有的几何信息。首先，高斯中心直接设为输入点云坐标，提供了对3D形状表面的粗略几何描述。在此基础上，方法利用CAP-UDF从点云学习一个神经无符号距离场（UDF）$f_u$，并通过梯度推断法线：$n_i = \nabla f_u(p_i) / \|\nabla f_u(p_i)\|$。值得注意的是，GAP选择2DGS（oriented disk）而非标准3DGS作为表示——2DGS将3D高斯椭球体替换为2D定向盘片，天然将法线编码为盘片方向，在表达精细局部几何时性能更优。每个盘片的旋转矩阵根据推断的法线初始化，确保盘片方向与局部表面对齐。渐进式多视角生成。核心挑战在于：点云没有颜色，高斯的外观从何而来？GAP的答案是利用文本引导的深度感知扩散模型逐视角生成外观图像，作为高斯优化的监督信号。具体而言，对于预设的视角序列$\{v_j\}_{j=1}^K$，方法渐进地为每个视角生成外观：利用学习到的UDF场通过射线行进计算深度图$D_j$，将当前渲染图像$I_j$、深度图$D_j$、掩码$M_j$和文本提示$c$一起送入基于ControlNet的深度感知修复扩散模型，生成几何一致的外观图像。这里的关键设计是三区域掩码策略——掩码被划分为生成区$M_{\mathrm{generate}}$（从未生成过的空白区域）、保留区$M_{\mathrm{keep}}$（已处理且当前视角无更优观察角度的区域）和更新区$M_{\mathrm{update}}$（当前视角提供更优观察角度的区域）。
\textbf{方法信息：} Baidu 参与了这篇论文。
\textbf{设计动机与边界：} 相对原始 3DGS 以独立高斯和图像重建误差为主的设定，GAP 把问题落在以下可实现的改动上：GAP方法核心流程包括三个阶段。首先，采用2DGS替代3DGS作为表示基础，通过CAP-UDF学习神经无符号距离场并推导点法线，用于初始化每个2D高斯盘的旋转矩阵，使其正确对齐表面方向。其次，设计渐进式高斯生成方案，从多个视角利用深度感知的文本到图像扩散模型合成一致外观图像，通过高斯优化策略从生成图像中学习高斯属性。最后，针对难以观测区域，提出基于扩散的高斯修复策略，利用扩散模型的生成能力补全缺失部分。表面锚定机制在优化过程中有效约束高斯位于3D形状表面，确保几何精度。该方法支持从合成点云、真实扫描到大尺度场景等不同复杂度级别的点云转高斯任务。其设计目的不是单独增加一个模块，而是用上述表示、监督或渲染约束消除原流程在目标任务中的具体失配。提出GAP方法，首次实现从无颜色属性的原始点云直接生成高保真3D高斯。通过多视角优化框架结合深度感知图像扩散模型生成一致外观，引入表面锚定机制约束高斯位于物体表面，并采用基于扩散的修复策略完成难以观测区域，覆盖从合成到真实扫描和大尺度场景多种复杂度。适用时还应保留论文给出的边界：依赖文本提示引导外观生成，可能存在文本与预期外观不匹配的情况；大尺度场景的优化时间较长；生成质量受扩散模型能力上限约束。

\section{SDS 与 Score 蒸馏方法}

SDS 方法虽然不是最终趋势，但它们是 3DGS 生成的入口锚点。它们证明显式高斯比 NeRF 更适合优化式 text-to-3D，因为 densification 和 rasterization 让几何和纹理收敛更快。
% METHOD: 2309.16653
\paragraph{DreamGaussian（生成式高斯喷洒锚点）（2023.09）}
DreamGaussian~\cite{2309_16653} 是把 3DGS 带入生成任务的基础锚点。它要解决的是 DreamFusion 类方法的效率瓶颈：SDS 可以把 2D 扩散模型的梯度传给 3D 表示，但若 3D 表示是 NeRF，每一步都要做昂贵体渲染，单个资产通常需要较长优化时间。DreamGaussian 的核心判断是，生成任务里不一定要用隐式场承载 SDS；原始 3DGS 的 tile-based 光栅化和显式高斯更新更适合快速试错。

第一阶段是生成式 Gaussian Splatting。方法在球体内随机采样少量 3D Gaussians，初始尺度较简单，颜色只保留 RGB，禁用球谐以避免在生成早期引入不必要的视角依赖。图像到 3D 时使用 Zero-1-to-3 XL 作为视角条件 2D 先验，并在参考视图加入 RGB 和 alpha 重建损失；文本到 3D 时使用 Stable Diffusion 的 SDS。每一步随机采样环绕相机，渲染当前高斯图像和透明度图，再把扩散模型的噪声预测误差反传到高斯位置、尺度、旋转、不透明度和颜色。它还采用时间步退火，让早期 SDS 提供粗形状梯度，后期逐渐转向细节。

DreamGaussian 真正重要的工程改动是重新设计 densification 节奏。重建任务里的 3DGS 会根据真实多视角误差逐渐克隆、分裂和剪枝；生成任务只有含糊的 SDS 梯度，NeRF 式 occupancy pruning 不容易判断哪里是空、哪里该保留。DreamGaussian 选择从更少基元开始，但更频繁地致密化，使新高斯随着生成过程逐步补到高梯度区域。论文报告图像到 3D 的高斯阶段约 500 步即可得到粗资产，代码 README 也把主流程拆成 \texttt{main.py} 的高斯阶段和 \texttt{main2.py} 的 mesh 细化阶段。

第二阶段是把生成好的高斯转成可用资产。由于 SDS 直接优化出的高斯往往纹理发糊，方法先通过局部密度查询和阈值化确定表面区域，再用 Marching Cubes 提取 mesh，把高斯颜色烘焙到 UV 纹理图上。随后在 UV 空间做可微渲染细化，用图像空间监督改善纹理细节。这个设计让输出不只是一个可实时渲染的高斯场，也能导出 textured mesh，方便进入 Blender 等下游流程。后续 GaussianDreamer 改它的初始化，MicroDreamer 改它的优化采样，GaussianDreamerPro 等方法继续约束高斯增长；这些方法都以 DreamGaussian 的“3DGS 承载 SDS”作为参照。
\textbf{锚点细化：} DreamGaussian：针对 DreamFusion 的 NeRF 射线采样导致 SDS 优化需数分钟的局限，用\textbf{3DGS 光栅化渲染}替代 NeRF 体渲染作为 SDS 优化的载体，核心思想是\textbf{利用 3DGS 基于 tile 的光栅化推理速度比 NeRF 快两个数量级的优势}——配合``少基元、高频致密化''策略（每 50--100 步致密化一次，梯度阈值降至 0.01）、禁用球谐系数仅保留 RGB、时间步线性退火三项适配，将单图到带纹理网格的生成时间从约 10 分钟压缩到约 2 分钟。
\textbf{方法信息：} Baidu 参与了这篇论文；代码：\url{https://github.com/dreamgaussian/dreamgaussian}；GitHub 关注度高，约 4342 stars；引用很多，Semantic Scholar 约 1016 次。



% METHOD: 2404.19525
\paragraph{MicroDreamer（分数式迭代重建加速）（2024.04）}
MicroDreamer~\cite{2404_19525} 针对的是 DreamGaussian/DreamFusion 式 SDS 的效率瓶颈。SDS 每一步都要调用扩散模型给当前 3D 渲染提供梯度，NFE 很高；如果在 latent diffusion 空间做，还要经过 VAE 编码器反传。MicroDreamer 的改变是提出 Score-based Iterative Reconstruction（SIR）：不是每一步只吃一次扩散梯度，而是先从多视角 score diffusion model 采样一组较一致的目标图像，再像普通多视角重建那样多次优化 3D 表示去拟合这组图像。

具体到 3DGS 路线，MicroDreamer 继承 DreamGaussian 的工程结构：\texttt{main.py} 负责高斯阶段，\texttt{main2.py} 负责 mesh refinement。但它把高斯阶段从约 500 次 SDS 优化改成少量 SIR 迭代，官方代码示例中高斯阶段约 20 次迭代、mesh 阶段约 3 次细化即可导出结果。训练策略还可以交替使用 SIR 与 SDS，并逐步增加采样视角数，让早期快速确定粗形状，后期补视角一致性。相对 DreamGaussian，它不是改 densification 或 mesh 提取，而是改“扩散监督如何喂给 3DGS”：把昂贵的逐步蒸馏改为少量多视角样本驱动的像素空间重建。
\textbf{实现细节：} 如果说 DreamGaussian 证明了 3DGS 与 SDS 的可行性，那么 MicroDreamer（2024.04） 则追问了一个更本质的问题：SDS 本身能否被重写为更高效的范式？动机。SDS 每次参数更新都需要一次扩散模型前向传播（NFE 过高），且损失在潜在空间计算需要反向传播通过 VAE 编码器，这两者共同导致优化速度缓慢；MicroDreamer 将"蒸馏"改写为"重建"范式来根本性地降低 NFE。MicroDreamer 的核心改动是将"蒸馏"范式改写为"重建"范式，以大幅降低 NFE（Number of Function Evaluations，函数评估次数）。该方法与 DreamGaussian 共享 3DGS 表示，但优化策略完全不同。方法指出 SDS 慢的两个根因：（1）每次3D参数更新都需要一次扩散模型前向传播（即NFE过高），整个优化需数万步扩散推理；（2）SDS损失在潜在空间计算，需要反向传播通过VAE编码器，进一步拖慢速度。直觉上，SDS像是"每走一步都要问一次老师"，而MicroDreamer改为"问一次老师后自己练习很多步"。为此，MicroDreamer提出Score-based Iterative Reconstruction（SIR，基于分数的迭代重建）。动机。SDS 的根本瓶颈在于每次 3D 参数更新都需调用扩散模型（NFE 过高）；SIR 将扩散推理与 3D 参数更新解耦，每次扩散推理服务多次 3D 参数更新，将 NFE 从数万降至约 160。SIR模仿可微分3D重建过程：在每次外循环迭代中，先用扩散模型精炼出一组多视角图像作为"伪真值"，然后在保持相机位姿不变的情况下对3D参数做多步内循环优化： 其中$c^{(k)}$是第$k$次迭代的相机位姿集合，$\hat{x}^{(k)}$是经扩散模型精炼的多视角图像。白话注解：SIR将"蒸馏"改写为"重建"。先让扩散模型生成一组一致的多视角图像，再用标准的$L_1$重建损失去拟合它们。关键优势在于：重建损失在像素空间直接计算，无需反向传播通过VAE；且每次扩散推理可服务多次3D参数更新，NFE从数万降至约160。
\textbf{设计动机与边界：} 相对原始 3DGS 以独立高斯和图像重建误差为主的设定，MicroDreamer 把问题落在以下可实现的改动上：MicroDreamer的核心是Score-based Iterative Reconstruction（SIR）算法。与SDS的单步优化不同，SIR给定一组从多视角分数扩散模型采样的图像，反复优化3D参数而非单步优化。具体流程：(1)从多视角扩散模型采样一组图像；(2)用这些图像反复优化3D表示（NeRF或3DGS）直至收敛；(3)训练改进包括交替使用SIR和SDS损失、逐步增加采样视角数等。SIR的关键优势在于可以在像素空间进行优化（而非LDM的潜在空间），避免了通过VAE编码器的反向传播，同时保持与各种3D表示的通用兼容性。其设计目的不是单独增加一个模块，而是用上述表示、监督或渲染约束消除原流程在目标任务中的具体失配。提出分数式迭代重建（SIR）算法，通过重复优化3D参数模拟可微3D重建过程；支持在像素空间而非潜在空间进行优化；相比SDS减少5-20倍NFE；在单A100 GPU上约20秒从3DGS生成mesh，比DreamGaussian快一倍且性能更优。适用时还应保留论文给出的边界：生成质量受多视角扩散模型限制；对复杂场景的生成仍有改进空间；mesh提取质量依赖GS质量。
\textbf{方法信息：} 代码：\url{https://github.com/ml-gsai/microdreamer}；GitHub 关注度较高，约 123 stars。
\section{最新直接高斯扩散}

最新一批方法开始减少多级转换，直接让扩散模型面对高斯资产或稀疏 3D latent。相对 DreamGaussian 的优化式 SDS 和 DiffSplat 的 latent 适配，这条路线更像“把 3DGS 变成扩散模型的原生输出格式”：训练时学习高斯属性的分布，采样时直接得到可渲染资产。它们仍很新，外部关注度还不高，但代表了生成式 3DGS 从后处理重建走向端到端资产生成的趋势。



% METHOD: 2309.16585
\paragraph{Text-to-3D using Gaussian Splatting（生成式 3DGS 与扩散先验）（2023.09）}
Text-to-3D using Gaussian Splatting~\cite{2309_16585} 关注的问题是：结合扩散或生成模型进行文本/图像到 3D 生成 Text-to-3D using Gaussian Splatting 是宽召回补全中发现的论文，主要根据标题、摘要和 arXiv 类别判定其与 3DGS 的关系。
相对原始 3DGS，Text-to-3D using Gaussian Splatting 的变化不是只替换损失函数，而是把表示、初始化、约束或网络预测路径重新组织到任务流程中：该工作被召回是因为论文文本包含 Gaussian/3D Gaussian 相关表述；当前分类先确定其任务域和与 3DGS 的相关性，后续证据流水线再补代码、企业和引用信息。
\textbf{实现细节：} 具体实现可以按输入、核心模块和渲染/优化路径来读。输入端首先取得任务所需的图像、相机、文本、时间序列、语义或几何线索，并把这些线索转换成可以约束高斯中心、尺度、旋转、不透明度、颜色或附加特征的条件。核心结构中，方法使用的主要机制是：从摘要层面看，Text-to-3D using Gaussian Splatting 围绕“结合扩散或生成模型进行文本/图像到 3D 生成”展开。若其正文或代码后续可读，应进一步核对它具体改动的是表示、初始化、渲染、优化、损失、动态建模、压缩或应用接口中的哪一环。 训练或推理时，这些模块通常先给出初始化或中间特征，再通过可微 Gaussian rasterization、重建损失、感知/语义约束、几何正则或任务损失回传到高斯参数；如果方法包含前馈网络，则网络输出的是高斯属性、变形场、压缩码、语义特征或控制变量，而不是只输出最终 2D 图像。渲染路径仍以显式高斯投影、深度排序和 alpha 合成为基础，只是在采样、属性预测、局部编辑、动态变形、压缩编码或下游任务接口处加入额外约束。这样写入本章时，它与前后方法的关系是：相对原始 3DGS 和本章已有锚点方法，它可以理解为 3DGS extends\_or\_uses 宽召回分类阶段的关系判断，后续需要按论文方法部分复核。
\textbf{证据与边界：} 本轮召回分类尚未逐项核对实验表格；后续证据补全阶段需要读取论文源码/PDF或项目页确认数据集、指标和对比方法。 该条来自宽召回批次；源码下载受 arXiv e-print 网络限制，部分论文暂时只能依据标题、摘要和类别分类，需要在入可信书前做人工/LLM精读复核。 本书只采用本地 paper pack、数据库分类、源码/引用登记和论文摘要中可核验的信息；没有本地证据支撑的速度、排行榜或开源完整性不在正文中强化。\textbf{方法信息：} 代码：\url{https://github.com/gsgen3d/gsgen}；引用信号：381 citations (11.4/month, domain-relative cutoff 3.8/month, semantic\_scholar)；代码关注：gsgen3d/gsgen has 845 stars (25.3/month, domain-relative cutoff 45.9/month)。
\textbf{设计动机与边界：} 相对原始 3DGS 以独立高斯和图像重建误差为主的设定，Text-to-3D using Gaussian Splatting 把问题落在以下可实现的改动上：从摘要层面看，Text-to-3D using Gaussian Splatting 围绕“结合扩散或生成模型进行文本/图像到 3D 生成”展开。若其正文或代码后续可读，应进一步核对它具体改动的是表示、初始化、渲染、优化、损失、动态建模、压缩或应用接口中的哪一环。其设计目的不是单独增加一个模块，而是用上述表示、监督或渲染约束消除原流程在目标任务中的具体失配。该工作被召回是因为论文文本包含 Gaussian/3D Gaussian 相关表述；当前分类先确定其任务域和与 3DGS 的相关性，后续证据流水线再补代码、企业和引用信息。适用时还应保留论文给出的边界：该条来自宽召回批次；源码下载受 arXiv e-print 网络限制，部分论文暂时只能依据标题、摘要和类别分类，需要在入可信书前做人工/LLM精读复核。
% METHOD: 2311.11284
\paragraph{LucidDreamer（生成式 3DGS 与扩散先验）（2023.11）}
LucidDreamer~\cite{2311_11284} 关注的问题是：结合扩散或生成模型进行文本/图像到 3D 生成 LucidDreamer 是宽召回补全中发现的论文，主要根据标题、摘要和 arXiv 类别判定其与 3DGS 的关系。
相对原始 3DGS，LucidDreamer 的变化不是只替换损失函数，而是把表示、初始化、约束或网络预测路径重新组织到任务流程中：该工作被召回是因为论文文本包含 Gaussian/3D Gaussian 相关表述；当前分类先确定其任务域和与 3DGS 的相关性，后续证据流水线再补代码、企业和引用信息。
\textbf{实现细节：} 具体实现可以按输入、核心模块和渲染/优化路径来读。输入端首先取得任务所需的图像、相机、文本、时间序列、语义或几何线索，并把这些线索转换成可以约束高斯中心、尺度、旋转、不透明度、颜色或附加特征的条件。核心结构中，方法使用的主要机制是：从摘要层面看，LucidDreamer 围绕“结合扩散或生成模型进行文本/图像到 3D 生成”展开。若其正文或代码后续可读，应进一步核对它具体改动的是表示、初始化、渲染、优化、损失、动态建模、压缩或应用接口中的哪一环。 训练或推理时，这些模块通常先给出初始化或中间特征，再通过可微 Gaussian rasterization、重建损失、感知/语义约束、几何正则或任务损失回传到高斯参数；如果方法包含前馈网络，则网络输出的是高斯属性、变形场、压缩码、语义特征或控制变量，而不是只输出最终 2D 图像。渲染路径仍以显式高斯投影、深度排序和 alpha 合成为基础，只是在采样、属性预测、局部编辑、动态变形、压缩编码或下游任务接口处加入额外约束。这样写入本章时，它与前后方法的关系是：相对原始 3DGS 和本章已有锚点方法，它可以理解为 3DGS extends\_or\_uses 宽召回分类阶段的关系判断，后续需要按论文方法部分复核。
\textbf{证据与边界：} 本轮召回分类尚未逐项核对实验表格；后续证据补全阶段需要读取论文源码/PDF或项目页确认数据集、指标和对比方法。 该条来自宽召回批次；源码下载受 arXiv e-print 网络限制，部分论文暂时只能依据标题、摘要和类别分类，需要在入可信书前做人工/LLM精读复核。 本书只采用本地 paper pack、数据库分类、源码/引用登记和论文摘要中可核验的信息；没有本地证据支撑的速度、排行榜或开源完整性不在正文中强化。\textbf{方法信息：} 代码：\url{https://github.com/luciddreamer-cvlab/LucidDreamer}；引用信号：332 citations (10.5/month, domain-relative cutoff 3.8/month, semantic\_scholar)；代码关注：luciddreamer-cvlab/luciddreamer has 1523 stars (48.0/month, domain-relative cutoff 45.9/month)。
\textbf{设计动机与边界：} 相对原始 3DGS 以独立高斯和图像重建误差为主的设定，LucidDreamer 把问题落在以下可实现的改动上：从摘要层面看，LucidDreamer 围绕“结合扩散或生成模型进行文本/图像到 3D 生成”展开。若其正文或代码后续可读，应进一步核对它具体改动的是表示、初始化、渲染、优化、损失、动态建模、压缩或应用接口中的哪一环。其设计目的不是单独增加一个模块，而是用上述表示、监督或渲染约束消除原流程在目标任务中的具体失配。该工作被召回是因为论文文本包含 Gaussian/3D Gaussian 相关表述；当前分类先确定其任务域和与 3DGS 的相关性，后续证据流水线再补代码、企业和引用信息。适用时还应保留论文给出的边界：该条来自宽召回批次；源码下载受 arXiv e-print 网络限制，部分论文暂时只能依据标题、摘要和类别分类，需要在入可信书前做人工/LLM精读复核。
% METHOD: 2311.16037
\paragraph{GaussianEditor（生成式 3DGS 与扩散先验）（2023.11）}
GaussianEditor~\cite{2311_16037} 关注的问题是：结合扩散或生成模型进行文本/图像到 3D 生成 GaussianEditor 是宽召回补全中发现的论文，主要根据标题、摘要和 arXiv 类别判定其与 3DGS 的关系。
相对原始 3DGS，GaussianEditor 的变化不是只替换损失函数，而是把表示、初始化、约束或网络预测路径重新组织到任务流程中：该工作被召回是因为论文文本包含 Gaussian/3D Gaussian 相关表述；当前分类先确定其任务域和与 3DGS 的相关性，后续证据流水线再补代码、企业和引用信息。
\textbf{实现细节：} 具体实现可以按输入、核心模块和渲染/优化路径来读。输入端首先取得任务所需的图像、相机、文本、时间序列、语义或几何线索，并把这些线索转换成可以约束高斯中心、尺度、旋转、不透明度、颜色或附加特征的条件。核心结构中，方法使用的主要机制是：从摘要层面看，GaussianEditor 围绕“结合扩散或生成模型进行文本/图像到 3D 生成”展开。若其正文或代码后续可读，应进一步核对它具体改动的是表示、初始化、渲染、优化、损失、动态建模、压缩或应用接口中的哪一环。 训练或推理时，这些模块通常先给出初始化或中间特征，再通过可微 Gaussian rasterization、重建损失、感知/语义约束、几何正则或任务损失回传到高斯参数；如果方法包含前馈网络，则网络输出的是高斯属性、变形场、压缩码、语义特征或控制变量，而不是只输出最终 2D 图像。渲染路径仍以显式高斯投影、深度排序和 alpha 合成为基础，只是在采样、属性预测、局部编辑、动态变形、压缩编码或下游任务接口处加入额外约束。这样写入本章时，它与前后方法的关系是：相对原始 3DGS 和本章已有锚点方法，它可以理解为 3DGS extends\_or\_uses 宽召回分类阶段的关系判断，后续需要按论文方法部分复核。
\textbf{证据与边界：} 本轮召回分类尚未逐项核对实验表格；后续证据补全阶段需要读取论文源码/PDF或项目页确认数据集、指标和对比方法。 该条来自宽召回批次；源码下载受 arXiv e-print 网络限制，部分论文暂时只能依据标题、摘要和类别分类，需要在入可信书前做人工/LLM精读复核。 本书只采用本地 paper pack、数据库分类、源码/引用登记和论文摘要中可核验的信息；没有本地证据支撑的速度、排行榜或开源完整性不在正文中强化。\textbf{方法信息：} 代码：\url{https://github.com/buaacyw/gaussianeditor}；引用信号：204 citations (6.5/month, domain-relative cutoff 3.8/month, semantic\_scholar)；代码关注：buaacyw/gaussianeditor has 1444 stars (45.9/month, domain-relative cutoff 45.9/month)。
\textbf{设计动机与边界：} 相对原始 3DGS 以独立高斯和图像重建误差为主的设定，GaussianEditor 把问题落在以下可实现的改动上：从摘要层面看，GaussianEditor 围绕“结合扩散或生成模型进行文本/图像到 3D 生成”展开。若其正文或代码后续可读，应进一步核对它具体改动的是表示、初始化、渲染、优化、损失、动态建模、压缩或应用接口中的哪一环。其设计目的不是单独增加一个模块，而是用上述表示、监督或渲染约束消除原流程在目标任务中的具体失配。该工作被召回是因为论文文本包含 Gaussian/3D Gaussian 相关表述；当前分类先确定其任务域和与 3DGS 的相关性，后续证据流水线再补代码、企业和引用信息。适用时还应保留论文给出的边界：该条来自宽召回批次；源码下载受 arXiv e-print 网络限制，部分论文暂时只能依据标题、摘要和类别分类，需要在入可信书前做人工/LLM精读复核。
% METHOD: 2311.17061
\paragraph{HumanGaussian（生成式 3DGS 与扩散先验）（2023.11）}
HumanGaussian~\cite{2311_17061} 关注的问题是：结合扩散或生成模型进行文本/图像到 3D 生成 HumanGaussian 是宽召回补全中发现的论文，主要根据标题、摘要和 arXiv 类别判定其与 3DGS 的关系。
相对原始 3DGS，HumanGaussian 的变化不是只替换损失函数，而是把表示、初始化、约束或网络预测路径重新组织到任务流程中：该工作被召回是因为论文文本包含 Gaussian/3D Gaussian 相关表述；当前分类先确定其任务域和与 3DGS 的相关性，后续证据流水线再补代码、企业和引用信息。
\textbf{实现细节：} 具体实现可以按输入、核心模块和渲染/优化路径来读。输入端首先取得任务所需的图像、相机、文本、时间序列、语义或几何线索，并把这些线索转换成可以约束高斯中心、尺度、旋转、不透明度、颜色或附加特征的条件。核心结构中，方法使用的主要机制是：从摘要层面看，HumanGaussian 围绕“结合扩散或生成模型进行文本/图像到 3D 生成”展开。若其正文或代码后续可读，应进一步核对它具体改动的是表示、初始化、渲染、优化、损失、动态建模、压缩或应用接口中的哪一环。 训练或推理时，这些模块通常先给出初始化或中间特征，再通过可微 Gaussian rasterization、重建损失、感知/语义约束、几何正则或任务损失回传到高斯参数；如果方法包含前馈网络，则网络输出的是高斯属性、变形场、压缩码、语义特征或控制变量，而不是只输出最终 2D 图像。渲染路径仍以显式高斯投影、深度排序和 alpha 合成为基础，只是在采样、属性预测、局部编辑、动态变形、压缩编码或下游任务接口处加入额外约束。这样写入本章时，它与前后方法的关系是：相对原始 3DGS 和本章已有锚点方法，它可以理解为 3DGS extends\_or\_uses 宽召回分类阶段的关系判断，后续需要按论文方法部分复核。
\textbf{证据与边界：} 本轮召回分类尚未逐项核对实验表格；后续证据补全阶段需要读取论文源码/PDF或项目页确认数据集、指标和对比方法。 该条来自宽召回批次；源码下载受 arXiv e-print 网络限制，部分论文暂时只能依据标题、摘要和类别分类，需要在入可信书前做人工/LLM精读复核。 本书只采用本地 paper pack、数据库分类、源码/引用登记和论文摘要中可核验的信息；没有本地证据支撑的速度、排行榜或开源完整性不在正文中强化。\textbf{方法信息：} 代码：\url{https://github.com/alvinliu0/humangaussian}；引用信号：119 citations (3.8/month, domain-relative cutoff 3.8/month, semantic\_scholar)；代码关注：alvinliu0/humangaussian has 492 stars (15.7/month, domain-relative cutoff 45.9/month)。
\textbf{设计动机与边界：} 相对原始 3DGS 以独立高斯和图像重建误差为主的设定，HumanGaussian 把问题落在以下可实现的改动上：从摘要层面看，HumanGaussian 围绕“结合扩散或生成模型进行文本/图像到 3D 生成”展开。若其正文或代码后续可读，应进一步核对它具体改动的是表示、初始化、渲染、优化、损失、动态建模、压缩或应用接口中的哪一环。其设计目的不是单独增加一个模块，而是用上述表示、监督或渲染约束消除原流程在目标任务中的具体失配。该工作被召回是因为论文文本包含 Gaussian/3D Gaussian 相关表述；当前分类先确定其任务域和与 3DGS 的相关性，后续证据流水线再补代码、企业和引用信息。适用时还应保留论文给出的边界：该条来自宽召回批次；源码下载受 arXiv e-print 网络限制，部分论文暂时只能依据标题、摘要和类别分类，需要在入可信书前做人工/LLM精读复核。
% METHOD: 2311.17907
\paragraph{CG3D（生成式 3DGS 与扩散先验）（2023.11）}
CG3D~\cite{2311_17907} 关注的问题是：结合扩散或生成模型进行文本/图像到 3D 生成 CG3D 是宽召回补全中发现的论文，主要根据标题、摘要和 arXiv 类别判定其与 3DGS 的关系。
相对原始 3DGS，CG3D 的变化不是只替换损失函数，而是把表示、初始化、约束或网络预测路径重新组织到任务流程中：该工作被召回是因为论文文本包含 Gaussian/3D Gaussian 相关表述；当前分类先确定其任务域和与 3DGS 的相关性，后续证据流水线再补代码、企业和引用信息。
\textbf{实现细节：} 具体实现可以按输入、核心模块和渲染/优化路径来读。输入端首先取得任务所需的图像、相机、文本、时间序列、语义或几何线索，并把这些线索转换成可以约束高斯中心、尺度、旋转、不透明度、颜色或附加特征的条件。核心结构中，方法使用的主要机制是：从摘要层面看，CG3D 围绕“结合扩散或生成模型进行文本/图像到 3D 生成”展开。若其正文或代码后续可读，应进一步核对它具体改动的是表示、初始化、渲染、优化、损失、动态建模、压缩或应用接口中的哪一环。 训练或推理时，这些模块通常先给出初始化或中间特征，再通过可微 Gaussian rasterization、重建损失、感知/语义约束、几何正则或任务损失回传到高斯参数；如果方法包含前馈网络，则网络输出的是高斯属性、变形场、压缩码、语义特征或控制变量，而不是只输出最终 2D 图像。渲染路径仍以显式高斯投影、深度排序和 alpha 合成为基础，只是在采样、属性预测、局部编辑、动态变形、压缩编码或下游任务接口处加入额外约束。这样写入本章时，它与前后方法的关系是：相对原始 3DGS 和本章已有锚点方法，它可以理解为 3DGS extends\_or\_uses 宽召回分类阶段的关系判断，后续需要按论文方法部分复核。
\textbf{证据与边界：} 本轮召回分类尚未逐项核对实验表格；后续证据补全阶段需要读取论文源码/PDF或项目页确认数据集、指标和对比方法。 该条来自宽召回批次；源码下载受 arXiv e-print 网络限制，部分论文暂时只能依据标题、摘要和类别分类，需要在入可信书前做人工/LLM精读复核。 本书只采用本地 paper pack、数据库分类、源码/引用登记和论文摘要中可核验的信息；没有本地证据支撑的速度、排行榜或开源完整性不在正文中强化。\textbf{方法信息：} 公开链接：\url{https://github.com/asvilesov/cg3d}，完整性按本地证据保守处理；引用信号：56 citations (1.8/month, domain-relative cutoff 3.8/month, semantic\_scholar)。
\textbf{设计动机与边界：} 相对原始 3DGS 以独立高斯和图像重建误差为主的设定，CG3D 把问题落在以下可实现的改动上：从摘要层面看，CG3D 围绕“结合扩散或生成模型进行文本/图像到 3D 生成”展开。若其正文或代码后续可读，应进一步核对它具体改动的是表示、初始化、渲染、优化、损失、动态建模、压缩或应用接口中的哪一环。其设计目的不是单独增加一个模块，而是用上述表示、监督或渲染约束消除原流程在目标任务中的具体失配。该工作被召回是因为论文文本包含 Gaussian/3D Gaussian 相关表述；当前分类先确定其任务域和与 3DGS 的相关性，后续证据流水线再补代码、企业和引用信息。适用时还应保留论文给出的边界：该条来自宽召回批次；源码下载受 arXiv e-print 网络限制，部分论文暂时只能依据标题、摘要和类别分类，需要在入可信书前做人工/LLM精读复核。
% METHOD: 2312.09147
\paragraph{Triplane Meets Gaussian Splatting（生成式 3DGS 与扩散先验）（2023.12）}
Triplane Meets Gaussian Splatting~\cite{2312_09147} 关注的问题是：结合扩散或生成模型进行文本/图像到 3D 生成 Triplane Meets Gaussian Splatting 是宽召回补全中发现的论文，主要根据标题、摘要和 arXiv 类别判定其与 3DGS 的关系。
相对原始 3DGS，Triplane Meets Gaussian Splatting 的变化不是只替换损失函数，而是把表示、初始化、约束或网络预测路径重新组织到任务流程中：该工作被召回是因为论文文本包含 Gaussian/3D Gaussian 相关表述；当前分类先确定其任务域和与 3DGS 的相关性，后续证据流水线再补代码、企业和引用信息。
\textbf{实现细节：} 具体实现可以按输入、核心模块和渲染/优化路径来读。输入端首先取得任务所需的图像、相机、文本、时间序列、语义或几何线索，并把这些线索转换成可以约束高斯中心、尺度、旋转、不透明度、颜色或附加特征的条件。核心结构中，方法使用的主要机制是：从摘要层面看，Triplane Meets Gaussian Splatting 围绕“结合扩散或生成模型进行文本/图像到 3D 生成”展开。若其正文或代码后续可读，应进一步核对它具体改动的是表示、初始化、渲染、优化、损失、动态建模、压缩或应用接口中的哪一环。 训练或推理时，这些模块通常先给出初始化或中间特征，再通过可微 Gaussian rasterization、重建损失、感知/语义约束、几何正则或任务损失回传到高斯参数；如果方法包含前馈网络，则网络输出的是高斯属性、变形场、压缩码、语义特征或控制变量，而不是只输出最终 2D 图像。渲染路径仍以显式高斯投影、深度排序和 alpha 合成为基础，只是在采样、属性预测、局部编辑、动态变形、压缩编码或下游任务接口处加入额外约束。这样写入本章时，它与前后方法的关系是：相对原始 3DGS 和本章已有锚点方法，它可以理解为 3DGS extends\_or\_uses 宽召回分类阶段的关系判断，后续需要按论文方法部分复核。
\textbf{证据与边界：} 本轮召回分类尚未逐项核对实验表格；后续证据补全阶段需要读取论文源码/PDF或项目页确认数据集、指标和对比方法。 该条来自宽召回批次；源码下载受 arXiv e-print 网络限制，部分论文暂时只能依据标题、摘要和类别分类，需要在入可信书前做人工/LLM精读复核。 本书只采用本地 paper pack、数据库分类、源码/引用登记和论文摘要中可核验的信息；没有本地证据支撑的速度、排行榜或开源完整性不在正文中强化。\textbf{方法信息：} 代码：\url{https://github.com/one-2-3-45/one-2-3-45}；引用信号：322 citations (10.4/month, domain-relative cutoff 3.8/month, semantic\_scholar)；代码关注：vast-ai-research/triplanegaussian has 924 stars (29.9/month, domain-relative cutoff 45.9/month)。
\textbf{设计动机与边界：} 相对原始 3DGS 以独立高斯和图像重建误差为主的设定，Triplane Meets Gaussian Splatting 把问题落在以下可实现的改动上：从摘要层面看，Triplane Meets Gaussian Splatting 围绕“结合扩散或生成模型进行文本/图像到 3D 生成”展开。若其正文或代码后续可读，应进一步核对它具体改动的是表示、初始化、渲染、优化、损失、动态建模、压缩或应用接口中的哪一环。其设计目的不是单独增加一个模块，而是用上述表示、监督或渲染约束消除原流程在目标任务中的具体失配。该工作被召回是因为论文文本包含 Gaussian/3D Gaussian 相关表述；当前分类先确定其任务域和与 3DGS 的相关性，后续证据流水线再补代码、企业和引用信息。适用时还应保留论文给出的边界：该条来自宽召回批次；源码下载受 arXiv e-print 网络限制，部分论文暂时只能依据标题、摘要和类别分类，需要在入可信书前做人工/LLM精读复核。
% METHOD: 2401.08742
\paragraph{Efficient4D（生成式 3DGS 与扩散先验）（2024.01）}
Efficient4D~\cite{2401_08742} 关注的问题是：结合扩散或生成模型进行文本/图像到 3D 生成 Efficient4D 是宽召回补全中发现的论文，主要根据标题、摘要和 arXiv 类别判定其与 3DGS 的关系。
相对原始 3DGS，Efficient4D 的变化不是只替换损失函数，而是把表示、初始化、约束或网络预测路径重新组织到任务流程中：该工作被召回是因为论文文本包含 Gaussian/3D Gaussian 相关表述；当前分类先确定其任务域和与 3DGS 的相关性，后续证据流水线再补代码、企业和引用信息。
\textbf{实现细节：} 具体实现可以按输入、核心模块和渲染/优化路径来读。输入端首先取得任务所需的图像、相机、文本、时间序列、语义或几何线索，并把这些线索转换成可以约束高斯中心、尺度、旋转、不透明度、颜色或附加特征的条件。核心结构中，方法使用的主要机制是：从摘要层面看，Efficient4D 围绕“结合扩散或生成模型进行文本/图像到 3D 生成”展开。若其正文或代码后续可读，应进一步核对它具体改动的是表示、初始化、渲染、优化、损失、动态建模、压缩或应用接口中的哪一环。 训练或推理时，这些模块通常先给出初始化或中间特征，再通过可微 Gaussian rasterization、重建损失、感知/语义约束、几何正则或任务损失回传到高斯参数；如果方法包含前馈网络，则网络输出的是高斯属性、变形场、压缩码、语义特征或控制变量，而不是只输出最终 2D 图像。渲染路径仍以显式高斯投影、深度排序和 alpha 合成为基础，只是在采样、属性预测、局部编辑、动态变形、压缩编码或下游任务接口处加入额外约束。这样写入本章时，它与前后方法的关系是：相对原始 3DGS 和本章已有锚点方法，它可以理解为 3DGS extends\_or\_uses 宽召回分类阶段的关系判断，后续需要按论文方法部分复核。
\textbf{证据与边界：} 本轮召回分类尚未逐项核对实验表格；后续证据补全阶段需要读取论文源码/PDF或项目页确认数据集、指标和对比方法。 该条来自宽召回批次；源码下载受 arXiv e-print 网络限制，部分论文暂时只能依据标题、摘要和类别分类，需要在入可信书前做人工/LLM精读复核。 本书只采用本地 paper pack、数据库分类、源码/引用登记和论文摘要中可核验的信息；没有本地证据支撑的速度、排行榜或开源完整性不在正文中强化。\textbf{方法信息：} 代码：\url{https://github.com/ashawkey/stable-dreamfusion}；引用信号：55 citations (1.8/month, domain-relative cutoff 3.8/month, semantic\_scholar)；代码关注：fudan-zvg/efficient4d has 102 stars (3.4/month, domain-relative cutoff 45.9/month)。
\textbf{设计动机与边界：} 相对原始 3DGS 以独立高斯和图像重建误差为主的设定，Efficient4D 把问题落在以下可实现的改动上：从摘要层面看，Efficient4D 围绕“结合扩散或生成模型进行文本/图像到 3D 生成”展开。若其正文或代码后续可读，应进一步核对它具体改动的是表示、初始化、渲染、优化、损失、动态建模、压缩或应用接口中的哪一环。其设计目的不是单独增加一个模块，而是用上述表示、监督或渲染约束消除原流程在目标任务中的具体失配。该工作被召回是因为论文文本包含 Gaussian/3D Gaussian 相关表述；当前分类先确定其任务域和与 3DGS 的相关性，后续证据流水线再补代码、企业和引用信息。适用时还应保留论文给出的边界：该条来自宽召回批次；源码下载受 arXiv e-print 网络限制，部分论文暂时只能依据标题、摘要和类别分类，需要在入可信书前做人工/LLM精读复核。
% METHOD: 2402.07207
\paragraph{GALA3D（生成式 3DGS 与扩散先验）（2024.02）}
GALA3D~\cite{2402_07207} 关注的问题是：结合扩散或生成模型进行文本/图像到 3D 生成 GALA3D 是宽召回补全中发现的论文，主要根据标题、摘要和 arXiv 类别判定其与 3DGS 的关系。
相对原始 3DGS，GALA3D 的变化不是只替换损失函数，而是把表示、初始化、约束或网络预测路径重新组织到任务流程中：该工作被召回是因为论文文本包含 Gaussian/3D Gaussian 相关表述；当前分类先确定其任务域和与 3DGS 的相关性，后续证据流水线再补代码、企业和引用信息。
\textbf{实现细节：} 具体实现可以按输入、核心模块和渲染/优化路径来读。输入端首先取得任务所需的图像、相机、文本、时间序列、语义或几何线索，并把这些线索转换成可以约束高斯中心、尺度、旋转、不透明度、颜色或附加特征的条件。核心结构中，方法使用的主要机制是：从摘要层面看，GALA3D 围绕“结合扩散或生成模型进行文本/图像到 3D 生成”展开。若其正文或代码后续可读，应进一步核对它具体改动的是表示、初始化、渲染、优化、损失、动态建模、压缩或应用接口中的哪一环。 训练或推理时，这些模块通常先给出初始化或中间特征，再通过可微 Gaussian rasterization、重建损失、感知/语义约束、几何正则或任务损失回传到高斯参数；如果方法包含前馈网络，则网络输出的是高斯属性、变形场、压缩码、语义特征或控制变量，而不是只输出最终 2D 图像。渲染路径仍以显式高斯投影、深度排序和 alpha 合成为基础，只是在采样、属性预测、局部编辑、动态变形、压缩编码或下游任务接口处加入额外约束。这样写入本章时，它与前后方法的关系是：相对原始 3DGS 和本章已有锚点方法，它可以理解为 3DGS extends\_or\_uses 宽召回分类阶段的关系判断，后续需要按论文方法部分复核。
\textbf{证据与边界：} 本轮召回分类尚未逐项核对实验表格；后续证据补全阶段需要读取论文源码/PDF或项目页确认数据集、指标和对比方法。 该条来自宽召回批次；源码下载受 arXiv e-print 网络限制，部分论文暂时只能依据标题、摘要和类别分类，需要在入可信书前做人工/LLM精读复核。 本书只采用本地 paper pack、数据库分类、源码/引用登记和论文摘要中可核验的信息；没有本地证据支撑的速度、排行榜或开源完整性不在正文中强化。\textbf{方法信息：} 公开链接：\url{https://github.com/vdigpku/gala3d}，完整性按本地证据保守处理；引用信号：125 citations (4.3/month, domain-relative cutoff 3.8/month, semantic\_scholar)；代码关注：vdigpku/gala3d has 306 stars (10.6/month, domain-relative cutoff 45.9/month)。
\textbf{设计动机与边界：} 相对原始 3DGS 以独立高斯和图像重建误差为主的设定，GALA3D 把问题落在以下可实现的改动上：从摘要层面看，GALA3D 围绕“结合扩散或生成模型进行文本/图像到 3D 生成”展开。若其正文或代码后续可读，应进一步核对它具体改动的是表示、初始化、渲染、优化、损失、动态建模、压缩或应用接口中的哪一环。其设计目的不是单独增加一个模块，而是用上述表示、监督或渲染约束消除原流程在目标任务中的具体失配。该工作被召回是因为论文文本包含 Gaussian/3D Gaussian 相关表述；当前分类先确定其任务域和与 3DGS 的相关性，后续证据流水线再补代码、企业和引用信息。适用时还应保留论文给出的边界：该条来自宽召回批次；源码下载受 arXiv e-print 网络限制，部分论文暂时只能依据标题、摘要和类别分类，需要在入可信书前做人工/LLM精读复核。
% METHOD: 2402.08682
\paragraph{IM-3D（生成式 3DGS 与扩散先验）（2024.02）}
IM-3D~\cite{2402_08682} 关注的问题是：结合扩散或生成模型进行文本/图像到 3D 生成 IM-3D 是宽召回补全中发现的论文，主要根据标题、摘要和 arXiv 类别判定其与 3DGS 的关系。
相对原始 3DGS，IM-3D 的变化不是只替换损失函数，而是把表示、初始化、约束或网络预测路径重新组织到任务流程中：该工作被召回是因为论文文本包含 Gaussian/3D Gaussian 相关表述；当前分类先确定其任务域和与 3DGS 的相关性，后续证据流水线再补代码、企业和引用信息。
\textbf{实现细节：} 具体实现可以按输入、核心模块和渲染/优化路径来读。输入端首先取得任务所需的图像、相机、文本、时间序列、语义或几何线索，并把这些线索转换成可以约束高斯中心、尺度、旋转、不透明度、颜色或附加特征的条件。核心结构中，方法使用的主要机制是：从摘要层面看，IM-3D 围绕“结合扩散或生成模型进行文本/图像到 3D 生成”展开。若其正文或代码后续可读，应进一步核对它具体改动的是表示、初始化、渲染、优化、损失、动态建模、压缩或应用接口中的哪一环。 训练或推理时，这些模块通常先给出初始化或中间特征，再通过可微 Gaussian rasterization、重建损失、感知/语义约束、几何正则或任务损失回传到高斯参数；如果方法包含前馈网络，则网络输出的是高斯属性、变形场、压缩码、语义特征或控制变量，而不是只输出最终 2D 图像。渲染路径仍以显式高斯投影、深度排序和 alpha 合成为基础，只是在采样、属性预测、局部编辑、动态变形、压缩编码或下游任务接口处加入额外约束。这样写入本章时，它与前后方法的关系是：相对原始 3DGS 和本章已有锚点方法，它可以理解为 3DGS extends\_or\_uses 宽召回分类阶段的关系判断，后续需要按论文方法部分复核。
\textbf{证据与边界：} 本轮召回分类尚未逐项核对实验表格；后续证据补全阶段需要读取论文源码/PDF或项目页确认数据集、指标和对比方法。 该条来自宽召回批次；源码下载受 arXiv e-print 网络限制，部分论文暂时只能依据标题、摘要和类别分类，需要在入可信书前做人工/LLM精读复核。 本书只采用本地 paper pack、数据库分类、源码/引用登记和论文摘要中可核验的信息；没有本地证据支撑的速度、排行榜或开源完整性不在正文中强化。\textbf{方法信息：} 代码：\url{https://github.com/3dtopia/openlrm}；引用信号：97 citations (3.4/month, domain-relative cutoff 3.8/month, semantic\_scholar)。
\textbf{设计动机与边界：} 相对原始 3DGS 以独立高斯和图像重建误差为主的设定，IM-3D 把问题落在以下可实现的改动上：从摘要层面看，IM-3D 围绕“结合扩散或生成模型进行文本/图像到 3D 生成”展开。若其正文或代码后续可读，应进一步核对它具体改动的是表示、初始化、渲染、优化、损失、动态建模、压缩或应用接口中的哪一环。其设计目的不是单独增加一个模块，而是用上述表示、监督或渲染约束消除原流程在目标任务中的具体失配。该工作被召回是因为论文文本包含 Gaussian/3D Gaussian 相关表述；当前分类先确定其任务域和与 3DGS 的相关性，后续证据流水线再补代码、企业和引用信息。适用时还应保留论文给出的边界：该条来自宽召回批次；源码下载受 arXiv e-print 网络限制，部分论文暂时只能依据标题、摘要和类别分类，需要在入可信书前做人工/LLM精读复核。
% METHOD: 2403.12957
\paragraph{GVGEN（生成式 3DGS 与扩散先验）（2024.03）}
GVGEN~\cite{2403_12957} 关注的问题是：结合扩散或生成模型进行文本/图像到 3D 生成 GVGEN 是宽召回补全中发现的论文，主要根据标题、摘要和 arXiv 类别判定其与 3DGS 的关系。
相对原始 3DGS，GVGEN 的变化不是只替换损失函数，而是把表示、初始化、约束或网络预测路径重新组织到任务流程中：该工作被召回是因为论文文本包含 Gaussian/3D Gaussian 相关表述；当前分类先确定其任务域和与 3DGS 的相关性，后续证据流水线再补代码、企业和引用信息。
\textbf{实现细节：} 具体实现可以按输入、核心模块和渲染/优化路径来读。输入端首先取得任务所需的图像、相机、文本、时间序列、语义或几何线索，并把这些线索转换成可以约束高斯中心、尺度、旋转、不透明度、颜色或附加特征的条件。核心结构中，方法使用的主要机制是：从摘要层面看，GVGEN 围绕“结合扩散或生成模型进行文本/图像到 3D 生成”展开。若其正文或代码后续可读，应进一步核对它具体改动的是表示、初始化、渲染、优化、损失、动态建模、压缩或应用接口中的哪一环。 训练或推理时，这些模块通常先给出初始化或中间特征，再通过可微 Gaussian rasterization、重建损失、感知/语义约束、几何正则或任务损失回传到高斯参数；如果方法包含前馈网络，则网络输出的是高斯属性、变形场、压缩码、语义特征或控制变量，而不是只输出最终 2D 图像。渲染路径仍以显式高斯投影、深度排序和 alpha 合成为基础，只是在采样、属性预测、局部编辑、动态变形、压缩编码或下游任务接口处加入额外约束。这样写入本章时，它与前后方法的关系是：相对原始 3DGS 和本章已有锚点方法，它可以理解为 3DGS extends\_or\_uses 宽召回分类阶段的关系判断，后续需要按论文方法部分复核。
\textbf{证据与边界：} 本轮召回分类尚未逐项核对实验表格；后续证据补全阶段需要读取论文源码/PDF或项目页确认数据集、指标和对比方法。 该条来自宽召回批次；源码下载受 arXiv e-print 网络限制，部分论文暂时只能依据标题、摘要和类别分类，需要在入可信书前做人工/LLM精读复核。 本书只采用本地 paper pack、数据库分类、源码/引用登记和论文摘要中可核验的信息；没有本地证据支撑的速度、排行榜或开源完整性不在正文中强化。\textbf{方法信息：} 代码：\url{https://github.com/SOTAMak1r/GVGEN}；引用信号：63 citations (2.3/month, domain-relative cutoff 3.8/month, semantic\_scholar)；代码关注：sotamak1r/gvgen has 123 stars (4.4/month, domain-relative cutoff 45.9/month)。
\textbf{设计动机与边界：} 相对原始 3DGS 以独立高斯和图像重建误差为主的设定，GVGEN 把问题落在以下可实现的改动上：从摘要层面看，GVGEN 围绕“结合扩散或生成模型进行文本/图像到 3D 生成”展开。若其正文或代码后续可读，应进一步核对它具体改动的是表示、初始化、渲染、优化、损失、动态建模、压缩或应用接口中的哪一环。其设计目的不是单独增加一个模块，而是用上述表示、监督或渲染约束消除原流程在目标任务中的具体失配。该工作被召回是因为论文文本包含 Gaussian/3D Gaussian 相关表述；当前分类先确定其任务域和与 3DGS 的相关性，后续证据流水线再补代码、企业和引用信息。适用时还应保留论文给出的边界：该条来自宽召回批次；源码下载受 arXiv e-print 网络限制，部分论文暂时只能依据标题、摘要和类别分类，需要在入可信书前做人工/LLM精读复核。
% METHOD: 2403.14621
\paragraph{GRM（生成式 3DGS 与扩散先验）（2024.03）}
GRM~\cite{2403_14621} 关注的问题是：结合扩散或生成模型进行文本/图像到 3D 生成 GRM 是宽召回补全中发现的论文，主要根据标题、摘要和 arXiv 类别判定其与 3DGS 的关系。
相对原始 3DGS，GRM 的变化不是只替换损失函数，而是把表示、初始化、约束或网络预测路径重新组织到任务流程中：该工作被召回是因为论文文本包含 Gaussian/3D Gaussian 相关表述；当前分类先确定其任务域和与 3DGS 的相关性，后续证据流水线再补代码、企业和引用信息。
\textbf{实现细节：} 具体实现可以按输入、核心模块和渲染/优化路径来读。输入端首先取得任务所需的图像、相机、文本、时间序列、语义或几何线索，并把这些线索转换成可以约束高斯中心、尺度、旋转、不透明度、颜色或附加特征的条件。核心结构中，方法使用的主要机制是：从摘要层面看，GRM 围绕“结合扩散或生成模型进行文本/图像到 3D 生成”展开。若其正文或代码后续可读，应进一步核对它具体改动的是表示、初始化、渲染、优化、损失、动态建模、压缩或应用接口中的哪一环。 训练或推理时，这些模块通常先给出初始化或中间特征，再通过可微 Gaussian rasterization、重建损失、感知/语义约束、几何正则或任务损失回传到高斯参数；如果方法包含前馈网络，则网络输出的是高斯属性、变形场、压缩码、语义特征或控制变量，而不是只输出最终 2D 图像。渲染路径仍以显式高斯投影、深度排序和 alpha 合成为基础，只是在采样、属性预测、局部编辑、动态变形、压缩编码或下游任务接口处加入额外约束。这样写入本章时，它与前后方法的关系是：相对原始 3DGS 和本章已有锚点方法，它可以理解为 3DGS extends\_or\_uses 宽召回分类阶段的关系判断，后续需要按论文方法部分复核。
\textbf{证据与边界：} 本轮召回分类尚未逐项核对实验表格；后续证据补全阶段需要读取论文源码/PDF或项目页确认数据集、指标和对比方法。 该条来自宽召回批次；源码下载受 arXiv e-print 网络限制，部分论文暂时只能依据标题、摘要和类别分类，需要在入可信书前做人工/LLM精读复核。 本书只采用本地 paper pack、数据库分类、源码/引用登记和论文摘要中可核验的信息；没有本地证据支撑的速度、排行榜或开源完整性不在正文中强化。\textbf{方法信息：} 公开链接：\url{https://github.com/justimyhxu/GRM}，完整性按本地证据保守处理；引用信号：330 citations (12.0/month, domain-relative cutoff 3.8/month, semantic\_scholar)；代码关注：justimyhxu/grm has 641 stars (23.2/month, domain-relative cutoff 45.9/month)。
\textbf{设计动机与边界：} 相对原始 3DGS 以独立高斯和图像重建误差为主的设定，GRM 把问题落在以下可实现的改动上：从摘要层面看，GRM 围绕“结合扩散或生成模型进行文本/图像到 3D 生成”展开。若其正文或代码后续可读，应进一步核对它具体改动的是表示、初始化、渲染、优化、损失、动态建模、压缩或应用接口中的哪一环。其设计目的不是单独增加一个模块，而是用上述表示、监督或渲染约束消除原流程在目标任务中的具体失配。该工作被召回是因为论文文本包含 Gaussian/3D Gaussian 相关表述；当前分类先确定其任务域和与 3DGS 的相关性，后续证据流水线再补代码、企业和引用信息。适用时还应保留论文给出的边界：该条来自宽召回批次；源码下载受 arXiv e-print 网络限制，部分论文暂时只能依据标题、摘要和类别分类，需要在入可信书前做人工/LLM精读复核。
% METHOD: 2411.17423
\paragraph{DRiVE（生成式 3DGS 与扩散先验）（2024.11）}
DRiVE~\cite{2411_17423} 关注的问题是：结合扩散或生成模型进行文本/图像到 3D 生成 DRiVE 是宽召回补全中发现的论文，主要根据标题、摘要和 arXiv 类别判定其与 3DGS 的关系。
相对原始 3DGS，DRiVE 的变化不是只替换损失函数，而是把表示、初始化、约束或网络预测路径重新组织到任务流程中：该工作被召回是因为论文文本包含 Gaussian/3D Gaussian 相关表述；当前分类先确定其任务域和与 3DGS 的相关性，后续证据流水线再补代码、企业和引用信息。
\textbf{实现细节：} 具体实现可以按输入、核心模块和渲染/优化路径来读。输入端首先取得任务所需的图像、相机、文本、时间序列、语义或几何线索，并把这些线索转换成可以约束高斯中心、尺度、旋转、不透明度、颜色或附加特征的条件。核心结构中，方法使用的主要机制是：从摘要层面看，DRiVE 围绕“结合扩散或生成模型进行文本/图像到 3D 生成”展开。若其正文或代码后续可读，应进一步核对它具体改动的是表示、初始化、渲染、优化、损失、动态建模、压缩或应用接口中的哪一环。 训练或推理时，这些模块通常先给出初始化或中间特征，再通过可微 Gaussian rasterization、重建损失、感知/语义约束、几何正则或任务损失回传到高斯参数；如果方法包含前馈网络，则网络输出的是高斯属性、变形场、压缩码、语义特征或控制变量，而不是只输出最终 2D 图像。渲染路径仍以显式高斯投影、深度排序和 alpha 合成为基础，只是在采样、属性预测、局部编辑、动态变形、压缩编码或下游任务接口处加入额外约束。这样写入本章时，它与前后方法的关系是：相对原始 3DGS 和本章已有锚点方法，它可以理解为 3DGS extends\_or\_uses 宽召回分类阶段的关系判断，后续需要按论文方法部分复核。
\textbf{证据与边界：} 本轮召回分类尚未逐项核对实验表格；后续证据补全阶段需要读取论文源码/PDF或项目页确认数据集、指标和对比方法。 该条来自宽召回批次；源码下载受 arXiv e-print 网络限制，部分论文暂时只能依据标题、摘要和类别分类，需要在入可信书前做人工/LLM精读复核。 本书只采用本地 paper pack、数据库分类、源码/引用登记和论文摘要中可核验的信息；没有本地证据支撑的速度、排行榜或开源完整性不在正文中强化。\textbf{方法信息：} 代码：\url{https://github.com/mickeykang16/vr-drive}。
\textbf{设计动机与边界：} 相对原始 3DGS 以独立高斯和图像重建误差为主的设定，DRiVE 把问题落在以下可实现的改动上：从摘要层面看，DRiVE 围绕“结合扩散或生成模型进行文本/图像到 3D 生成”展开。若其正文或代码后续可读，应进一步核对它具体改动的是表示、初始化、渲染、优化、损失、动态建模、压缩或应用接口中的哪一环。其设计目的不是单独增加一个模块，而是用上述表示、监督或渲染约束消除原流程在目标任务中的具体失配。该工作被召回是因为论文文本包含 Gaussian/3D Gaussian 相关表述；当前分类先确定其任务域和与 3DGS 的相关性，后续证据流水线再补代码、企业和引用信息。适用时还应保留论文给出的边界：该条来自宽召回批次；源码下载受 arXiv e-print 网络限制，部分论文暂时只能依据标题、摘要和类别分类，需要在入可信书前做人工/LLM精读复核。
% METHOD: 2605.20752
\paragraph{GaussianDream（生成式 3DGS 与扩散先验）（2026.05）}
GaussianDream~\cite{2605_20752} 关注的问题是：结合扩散或生成模型进行文本/图像到 3D 生成 GaussianDream 是宽召回补全中发现的论文，主要根据标题、摘要和 arXiv 类别判定其与 3DGS 的关系。
相对原始 3DGS，GaussianDream 的变化不是只替换损失函数，而是把表示、初始化、约束或网络预测路径重新组织到任务流程中：该工作被召回是因为论文文本包含 Gaussian/3D Gaussian 相关表述；当前分类先确定其任务域和与 3DGS 的相关性，后续证据流水线再补代码、企业和引用信息。
\textbf{实现细节：} 具体实现可以按输入、核心模块和渲染/优化路径来读。输入端首先取得任务所需的图像、相机、文本、时间序列、语义或几何线索，并把这些线索转换成可以约束高斯中心、尺度、旋转、不透明度、颜色或附加特征的条件。核心结构中，方法使用的主要机制是：从摘要层面看，GaussianDream 围绕“结合扩散或生成模型进行文本/图像到 3D 生成”展开。若其正文或代码后续可读，应进一步核对它具体改动的是表示、初始化、渲染、优化、损失、动态建模、压缩或应用接口中的哪一环。 训练或推理时，这些模块通常先给出初始化或中间特征，再通过可微 Gaussian rasterization、重建损失、感知/语义约束、几何正则或任务损失回传到高斯参数；如果方法包含前馈网络，则网络输出的是高斯属性、变形场、压缩码、语义特征或控制变量，而不是只输出最终 2D 图像。渲染路径仍以显式高斯投影、深度排序和 alpha 合成为基础，只是在采样、属性预测、局部编辑、动态变形、压缩编码或下游任务接口处加入额外约束。这样写入本章时，它与前后方法的关系是：相对原始 3DGS 和本章已有锚点方法，它可以理解为 3DGS extends\_or\_uses 宽召回分类阶段的关系判断，后续需要按论文方法部分复核。
\textbf{证据与边界：} 本轮召回分类尚未逐项核对实验表格；后续证据补全阶段需要读取论文源码/PDF或项目页确认数据集、指标和对比方法。 该条来自宽召回批次；源码下载受 arXiv e-print 网络限制，部分论文暂时只能依据标题、摘要和类别分类，需要在入可信书前做人工/LLM精读复核。 本书只采用本地 paper pack、数据库分类、源码/引用登记和论文摘要中可核验的信息；没有本地证据支撑的速度、排行榜或开源完整性不在正文中强化。\textbf{方法信息：} 代码：\url{https://github.com/wuguanjun0627/gaussiandreamer}。
\textbf{设计动机与边界：} 相对原始 3DGS 以独立高斯和图像重建误差为主的设定，GaussianDream 把问题落在以下可实现的改动上：从摘要层面看，GaussianDream 围绕“结合扩散或生成模型进行文本/图像到 3D 生成”展开。若其正文或代码后续可读，应进一步核对它具体改动的是表示、初始化、渲染、优化、损失、动态建模、压缩或应用接口中的哪一环。其设计目的不是单独增加一个模块，而是用上述表示、监督或渲染约束消除原流程在目标任务中的具体失配。该工作被召回是因为论文文本包含 Gaussian/3D Gaussian 相关表述；当前分类先确定其任务域和与 3DGS 的相关性，后续证据流水线再补代码、企业和引用信息。适用时还应保留论文给出的边界：该条来自宽召回批次；源码下载受 arXiv e-print 网络限制，部分论文暂时只能依据标题、摘要和类别分类，需要在入可信书前做人工/LLM精读复核。
% METHOD: 2606.24874
\paragraph{FLUX3D（扩散对齐稀疏 3D latent）（2026.06）}
FLUX3D~\cite{2606_24874} 继承的是“稀疏 3D latent + 扩散生成”的路线，目标是从单张图像生成高保真 3DGS。它指出现有稀疏体素 latent 方法有两个瓶颈：第一，很多方法用 DINOv2 等判别式 2D 特征构建 3D latent，这类特征更适合语义抽象，不擅长保留纹理、文字、边缘等重建细节；第二，扩散阶段里，密集 2D image tokens 和稀疏 3D voxel latents 没有明确对齐，图像局部证据很难准确流入对应的 3D 单元。

为解决表示瓶颈，FLUX3D 提出 DA-SLAT（Diffusion-Aligned Structured Latents）。它不再用纯判别特征做稀疏体素 latent，而是从 FLUX 这类图像扩散模型中抽取更偏重合成和重建的特征，并把这些特征组织成稀疏 voxel latent。随后使用 decoder-only 架构直接把 structured latent 映射到 3DGS 输出，减少 encoder-decoder 压缩带来的外观损失。直观上，它让 3D latent 更像“能重建图像细节的 3D 容器”，而不只是语义标签的三维摘要。

为解决生成阶段的对齐瓶颈，FLUX3D 设计 SMDiT（Sparse-structure Multimodal Diffusion Transformer）和 MARoPE（Modal-Aware Rotary Positional Embedding）。SMDiT 在稀疏体素拓扑上处理 3D tokens，同时让 2D 图像 tokens 和 3D voxel tokens 做跨模态交互；MARoPE 把 2D tokens 放到虚拟 3D 平面中编码，同时保留体素自身坐标，使模型可以在不显式求几何匹配的情况下学习 2D-3D 对应。相对 DiffSplat 的 splat latent 方案，FLUX3D 更强调稀疏 3D 结构内部的局部对应和外观保真。
\textbf{方法信息：} ByteDance 参与了这篇论文；目前没有可靠完整公开代码记录。
\textbf{设计动机与边界：} 相对原始 3DGS 以独立高斯和图像重建误差为主的设定，FLUX3D 把问题落在以下可实现的改动上：使用更偏重重建线索的 2D 表示构建稀疏体素 latent，并在扩散 transformer 中加入密集 2D token 到稀疏 3D voxel latent 的对齐模块，最终生成 3D Gaussian 表示。其设计目的不是单独增加一个模块，而是用上述表示、监督或渲染约束消除原流程在目标任务中的具体失配。识别语义 2D 特征和稀疏 3D latent 之间的表示瓶颈，并设计可扩展机制对齐图像 token 与稀疏体素 latent。适用时还应保留论文给出的边界：依赖大规模扩散先验和稀疏体素设计，可能对训练成本、输入域泛化和复杂拓扑存在限制。
% METHOD: 2607.01803
\paragraph{PixGS（像素空间直接高斯扩散）（2026.07）}
PixGS~\cite{2607_01803} 是对 DiffSplat 类 latent 管线的进一步简化。DiffSplat 需要先训练高斯重建器，再训练能压缩高斯属性的 VAE，最后训练 latent diffusion；这个级联路线能复用图像扩散先验，但每一级都会引入容量瓶颈和误差传递。PixGS 的改变是直接使用 pixel-space diffusion，在像素空间的高斯属性张量上去噪，绕开 VAE latent 压缩。

具体表示上，它同样把 3D asset 写成 image-like Gaussian attribute tensors，每个像素或网格位置编码一个高斯的几何和外观字段。扩散模型每个时间步直接预测或修正这些高斯属性，而不是先生成多视图图像，也不是先生成压缩 latent 再解码。由于高斯属性本身可被光栅化渲染，训练损失可以直接落到最终 3DGS：渲染图像保证视图外观，深度和法线监督约束几何，LoG 等高频结构损失帮助保留文字、薄结构和局部纹理。

PixGS 还用三阶段训练策略稳定从重建到生成的过渡，使生成模型不完全受单个重建器收敛质量限制。论文报告其在文本和图像条件 3DGS 生成中能以约 1 秒速度产出资产。相对 FLUX3D 的“稀疏体素 latent + 2D-3D token 对齐”，PixGS 更强调端到端的直接 splat-level denoising：减少中间表征，直接在可渲染高斯字段上建模。
\textbf{实现细节：} PixGS。PixGS（2026.07） 可作为"Pixel-space diffusion for Gaussian splat generation"中的新锚点候选。相较于 3DGS generation，主要改变是：PixGS：避免复杂 cascade 和 latent-space 瓶颈的 pixel-space diffusion 3DGS 生成。3DGS 的显式高斯表示和可微渲染能力保持为基础。其设计动机是：针对高质量 3D 数据稀缺、2D 生成器视角不一致和复杂级联难扩展的问题，提出更直接的 3DGS 生成路径。具体而言，不同于把大规模 text-to-image latent diffusion 复杂适配到 3DGS 的方法，PixGS 通过 pixel-space diffusion 直接生成 Gaussian splat assets，以降低压缩 latent 和多组件串联造成的 artifacts。这一锚点判断来自当前分类中的机制关系：候选锚点：生成路线，且相关性 0.95；后续方法可围绕 直接生成 3D Gaussian Splats 而非通过复杂级联间接构造。继续比较和扩展。适用边界是：仍受训练数据规模、扩散采样成本和 3D 监督质量限制；复杂拓扑和高频细节需进一步验证。
\textbf{方法信息：} Qualcomm 参与了这篇论文；目前没有可靠完整公开代码记录。

% METHOD: 2607.03819
\paragraph{CGGS（生成式 3DGS 与扩散先验）（2026.07）}
CGGS~\cite{2607_03819} 关注的问题是：结合扩散或生成模型进行文本/图像到 3D 生成 CGGS 是宽召回补全中发现的论文，主要根据标题、摘要和 arXiv 类别判定其与 3DGS 的关系。
相对原始 3DGS，CGGS 的变化不是只替换损失函数，而是把表示、初始化、约束或网络预测路径重新组织到任务流程中：该工作被召回是因为论文文本包含 Gaussian/3D Gaussian 相关表述；当前分类先确定其任务域和与 3DGS 的相关性，后续证据流水线再补代码、企业和引用信息。
\textbf{实现细节：} 具体实现可以按输入、核心模块和渲染/优化路径来读。输入端首先取得任务所需的图像、相机、文本、时间序列、语义或几何线索，并把这些线索转换成可以约束高斯中心、尺度、旋转、不透明度、颜色或附加特征的条件。核心结构中，方法使用的主要机制是：从摘要层面看，CGGS 围绕“结合扩散或生成模型进行文本/图像到 3D 生成”展开。若其正文或代码后续可读，应进一步核对它具体改动的是表示、初始化、渲染、优化、损失、动态建模、压缩或应用接口中的哪一环。 训练或推理时，这些模块通常先给出初始化或中间特征，再通过可微 Gaussian rasterization、重建损失、感知/语义约束、几何正则或任务损失回传到高斯参数；如果方法包含前馈网络，则网络输出的是高斯属性、变形场、压缩码、语义特征或控制变量，而不是只输出最终 2D 图像。渲染路径仍以显式高斯投影、深度排序和 alpha 合成为基础，只是在采样、属性预测、局部编辑、动态变形、压缩编码或下游任务接口处加入额外约束。这样写入本章时，它与前后方法的关系是：相对原始 3DGS 和本章已有锚点方法，它可以理解为 3DGS extends\_or\_uses 宽召回分类阶段的关系判断，后续需要按论文方法部分复核。
\textbf{证据与边界：} 本轮召回分类尚未逐项核对实验表格；后续证据补全阶段需要读取论文源码/PDF或项目页确认数据集、指标和对比方法。 该条来自宽召回批次；源码下载受 arXiv e-print 网络限制，部分论文暂时只能依据标题、摘要和类别分类，需要在入可信书前做人工/LLM精读复核。 本书只采用本地 paper pack、数据库分类、源码/引用登记和论文摘要中可核验的信息；没有本地证据支撑的速度、排行榜或开源完整性不在正文中强化。\textbf{方法信息：} 代码：\url{https://github.com/cggs-26/cggs}。


本章方法可以按生成范式分成四代来理解。第一代是 DreamGaussian 和 GaussianDreamer：用 SDS 优化高斯，并用 3D 扩散或点云改善初始化。第二代是高斯重排和潜空间：DiffSplat、UVGS、Zero-1-to-G、DiffGS、Can3Tok、UniGS 把无结构高斯变成 latent、UV 图、函数或 token。第三代是把高斯放入扩散或视频模型内部：DiffusionGS、DSplats、SplatFlow、Lyra、Lyra 2.0 和 GGS 让 denoiser 或视频扩散直接预测高斯，并进一步处理长轨迹世界生成中的记忆和漂移问题。第四代扩展到场景和 4D：LucidDreamer、BloomScene、GaussianCity、ScenDi、Diffusion4D、DreamGaussian4D 和 CharacterShot 解决大场景、动态和可控动画。贯穿其中的共同变化是：3DGS 从一个可优化的渲染表示，逐渐变成生成模型可以直接建模、压缩、采样和编辑的三维资产格式。

\textbf{设计动机与边界：} 相对原始 3DGS 以独立高斯和图像重建误差为主的设定，CGGS 把问题落在以下可实现的改动上：从摘要层面看，CGGS 围绕“结合扩散或生成模型进行文本/图像到 3D 生成”展开。若其正文或代码后续可读，应进一步核对它具体改动的是表示、初始化、渲染、优化、损失、动态建模、压缩或应用接口中的哪一环。其设计目的不是单独增加一个模块，而是用上述表示、监督或渲染约束消除原流程在目标任务中的具体失配。该工作被召回是因为论文文本包含 Gaussian/3D Gaussian 相关表述；当前分类先确定其任务域和与 3DGS 的相关性，后续证据流水线再补代码、企业和引用信息。适用时还应保留论文给出的边界：该条来自宽召回批次；源码下载受 arXiv e-print 网络限制，部分论文暂时只能依据标题、摘要和类别分类，需要在入可信书前做人工/LLM精读复核。
\textbf{尾项实现澄清：} CGGS 的动机是第一人称视角重叠少、单人视角主导，导致文本生成的多视图内容与几何不一致。流程明确分三步：先以一致性增强损失微调多视图潜扩散模型，生成文本对齐的二维先验；再以光流和点轨迹对应估计深度并形成稠密点云布局；最后从该布局初始化高斯，由 MID 深度损失和分层优化精化几何与外观。因而它相对 SDS 类生成锚点的关键是把一致二维先验、布局和高斯精化串成可检查的训练流程。
